%% fileadv.tex
%%
%% Copyright (C) 2000-2018 Simone Piccardi.  Permission is granted to
%% copy, distribute and/or modify this document under the terms of the GNU Free
%% Documentation License, Version 1.1 or any later version published by the
%% Free Software Foundation; with the Invariant Sections being "Un preambolo",
%% with no Front-Cover Texts, and with no Back-Cover Texts.  A copy of the
%% license is included in the section entitled "GNU Free Documentation
%% License".
%%
\chapter{La gestione avanzata dei file}
\label{cha:file_advanced}

In questo capitolo affronteremo le tematiche relative alla gestione avanzata
dei file. Inizieremo con la trattazione delle problematiche del \textit{file
  locking} e poi prenderemo in esame le varie funzionalità avanzate che
permettono una gestione più sofisticata dell'I/O su file, a partire da quelle
che consentono di gestire l'accesso contemporaneo a più file esaminando le
varie modalità alternative di gestire l'I/O per concludere con la gestione dei
file mappati in memoria e le altre funzioni avanzate che consentono un
controllo più dettagliato delle modalità di I/O.


\section{Il \textit{file locking}}
\label{sec:file_locking}

\itindbeg{file~locking}

In sez.~\ref{sec:file_shared_access} abbiamo preso in esame le modalità in cui
un sistema unix-like gestisce l'accesso concorrente ai file da parte di
processi diversi. In quell'occasione si è visto come, con l'eccezione dei file
aperti in \textit{append mode}, quando più processi scrivono
contemporaneamente sullo stesso file non è possibile determinare la sequenza
in cui essi opereranno.

Questo causa la possibilità di una \textit{race condition}; in generale le
situazioni più comuni sono due: l'interazione fra un processo che scrive e
altri che leggono, in cui questi ultimi possono leggere informazioni scritte
solo in maniera parziale o incompleta; o quella in cui diversi processi
scrivono, mescolando in maniera imprevedibile il loro output sul file.

In tutti questi casi il \textit{file locking} è la tecnica che permette di
evitare le \textit{race condition}, attraverso una serie di funzioni che
permettono di bloccare l'accesso al file da parte di altri processi, così da
evitare le sovrapposizioni, e garantire la atomicità delle operazioni di
lettura o scrittura.


\subsection{L'\textit{advisory locking}}
\label{sec:file_record_locking}

La prima modalità di \textit{file locking} che è stata implementata nei
sistemi unix-like è quella che viene usualmente chiamata \textit{advisory
  locking},\footnote{Stevens in \cite{APUE} fa riferimento a questo argomento
  come al \textit{record locking}, dizione utilizzata anche dal manuale della
  \acr{glibc}; nelle pagine di manuale si parla di \textit{discrectionary file
    lock} per \func{fcntl} e di \textit{advisory locking} per \func{flock},
  mentre questo nome viene usato da Stevens per riferirsi al \textit{file
    locking} POSIX. Dato che la dizione \textit{record locking} è quantomeno
  ambigua, in quanto in un sistema Unix non esiste niente che possa fare
  riferimento al concetto di \textit{record}, alla fine si è scelto di
  mantenere il nome \textit{advisory locking}.} in quanto sono i singoli
processi, e non il sistema, che si incaricano di asserire e verificare se
esistono delle condizioni di blocco per l'accesso ai file. 

Questo significa che le funzioni \func{read} o \func{write} vengono eseguite
comunque e non risentono affatto della presenza di un eventuale \textit{lock};
pertanto è sempre compito dei vari processi che intendono usare il
\textit{file locking}, controllare esplicitamente lo stato dei file condivisi
prima di accedervi, utilizzando le relative funzioni.

In generale si distinguono due tipologie di \textit{file lock};\footnote{di
  seguito ci riferiremo sempre ai blocchi di accesso ai file con la
  nomenclatura inglese di \textit{file lock}, o più brevemente con
  \textit{lock}, per evitare confusioni linguistiche con il blocco di un
  processo (cioè la condizione in cui il processo viene posto in stato di
  \textit{sleep}).} la prima è il cosiddetto \textit{shared lock}, detto anche
\textit{read lock} in quanto serve a bloccare l'accesso in scrittura su un
file affinché il suo contenuto non venga modificato mentre lo si legge. Si
parla appunto di \textsl{blocco condiviso} in quanto più processi possono
richiedere contemporaneamente uno \textit{shared lock} su un file per
proteggere il loro accesso in lettura.

La seconda tipologia è il cosiddetto \textit{exclusive lock}, detto anche
\textit{write lock} in quanto serve a bloccare l'accesso su un file (sia in
lettura che in scrittura) da parte di altri processi mentre lo si sta
scrivendo. Si parla di \textsl{blocco esclusivo} appunto perché un solo
processo alla volta può richiedere un \textit{exclusive lock} su un file per
proteggere il suo accesso in scrittura.

In Linux sono disponibili due interfacce per utilizzare l'\textit{advisory
  locking}, la prima è quella derivata da BSD, che è basata sulla funzione
\func{flock}, la seconda è quella recepita dallo standard POSIX.1 (che è
derivata dall'interfaccia usata in System V), che è basata sulla funzione
\func{fcntl}.  I \textit{file lock} sono implementati in maniera completamente
indipendente nelle due interfacce (in realtà con Linux questo avviene solo
dalla serie 2.0 dei kernel) che pertanto possono coesistere senza
interferenze.

Entrambe le interfacce prevedono la stessa procedura di funzionamento: si
inizia sempre con il richiedere l'opportuno \textit{file lock} (un
\textit{exclusive lock} per una scrittura, uno \textit{shared lock} per una
lettura) prima di eseguire l'accesso ad un file.  Se il blocco viene acquisito
il processo prosegue l'esecuzione, altrimenti (a meno di non aver richiesto un
comportamento non bloccante) viene posto in stato di \textit{sleep}. Una volta
finite le operazioni sul file si deve provvedere a rimuovere il blocco.

La situazione delle varie possibilità che si possono verificare è riassunta in
tab.~\ref{tab:file_file_lock}, dove si sono riportati, a seconda delle varie
tipologie di blocco già presenti su un file, il risultato che si avrebbe in
corrispondenza di una ulteriore richiesta da parte di un processo di un blocco
nelle due tipologie di \textit{file lock} menzionate, con un successo o meno
della richiesta.

\begin{table}[htb]
  \centering
  \footnotesize
   \begin{tabular}[c]{|l|c|c|c|}
    \hline
    \textbf{Richiesta} & \multicolumn{3}{|c|}{\textbf{Stato del file}}\\
    \cline{2-4}
                &Nessun \textit{lock}&\textit{Read lock}&\textit{Write lock}\\
    \hline
    \hline
    \textit{Read lock} & esecuzione & esecuzione & blocco \\
    \textit{Write lock}& esecuzione & blocco & blocco \\
    \hline    
  \end{tabular}
  \caption{Tipologie di \textit{file locking}.}
  \label{tab:file_file_lock}
\end{table}

Si tenga presente infine che il controllo di accesso e la gestione dei
permessi viene effettuata quando si apre un file, l'unico controllo residuo
che si può avere riguardo il \textit{file locking} è che il tipo di blocco che
si vuole ottenere su un file deve essere compatibile con le modalità di
apertura dello stesso (in lettura per un \textit{read lock} e in scrittura per
un \textit{write lock}).

%%  Si ricordi che
%% la condizione per acquisire uno \textit{shared lock} è che il file non abbia
%% già un \textit{exclusive lock} attivo, mentre per acquisire un
%% \textit{exclusive lock} non deve essere presente nessun tipo di blocco.


\subsection{La funzione \func{flock}} 
\label{sec:file_flock}

La prima interfaccia per il \textit{file locking}, quella derivata da BSD,
permette di eseguire un blocco solo su un intero file; la funzione di sistema
usata per richiedere e rimuovere un \textit{file lock} è \funcd{flock}, ed il
suo prototipo è:

\begin{funcproto}{
\fhead{sys/file.h}
\fdecl{int flock(int fd, int operation)}
\fdesc{Applica o rimuove un \textit{file lock}.} 
}

{La funzione ritorna $0$ in caso di successo e $-1$ per un errore, nel qual
  caso \var{errno} assumerà uno dei valori: 
  \begin{errlist}
  \item[\errcode{EINTR}] la funzione è stata interrotta da un segnale
    nell'attesa dell'acquisizione di un \textit{file lock}.
  \item[\errcode{EINVAL}] si è specificato un valore non valido
    per \param{operation}.
  \item[\errcode{ENOLCK}] il kernel non ha memoria sufficiente per gestire il
    \textit{file lock}.
  \item[\errcode{EWOULDBLOCK}] il file ha già un blocco attivo, e si è
    specificato \const{LOCK\_NB}.
  \end{errlist}
  ed inoltre \errval{EBADF} nel suo significato generico.
}
\end{funcproto}

La funzione può essere usata per acquisire o rilasciare un \textit{file lock}
a seconda di quanto specificato tramite il valore dell'argomento
\param{operation}; questo viene interpretato come maschera binaria, e deve
essere passato costruendo il valore con un OR aritmetico delle costanti
riportate in tab.~\ref{tab:file_flock_operation}.

\begin{table}[htb]
  \centering
  \footnotesize
  \begin{tabular}[c]{|l|p{6cm}|}
    \hline
    \textbf{Valore} & \textbf{Significato} \\
    \hline
    \hline
    \constd{LOCK\_SH} & Richiede uno \textit{shared lock} sul file.\\ 
    \constd{LOCK\_EX} & Richiede un \textit{esclusive lock} sul file.\\
    \constd{LOCK\_UN} & Rilascia il \textit{file lock}.\\
    \constd{LOCK\_NB} & Impedisce che la funzione si blocchi nella
                        richiesta di un \textit{file lock}.\\
    \hline    
  \end{tabular}
  \caption{Valori dell'argomento \param{operation} di \func{flock}.}
  \label{tab:file_flock_operation}
\end{table}

I primi due valori, \const{LOCK\_SH} e \const{LOCK\_EX} permettono di
richiedere un \textit{file lock} rispettivamente condiviso o esclusivo, ed
ovviamente non possono essere usati insieme. Se con essi si specifica anche
\const{LOCK\_NB} la funzione non si bloccherà qualora il \textit{file lock}
non possa essere acquisito, ma ritornerà subito con un errore di
\errcode{EWOULDBLOCK}. Per rilasciare un \textit{file lock} si dovrà invece
usare direttamente \const{LOCK\_UN}.

Si tenga presente che non esiste una modalità per eseguire atomicamente un
cambiamento del tipo di blocco (da \textit{shared lock} a \textit{esclusive
  lock}), il blocco deve essere prima rilasciato e poi richiesto, ed è sempre
possibile che nel frattempo abbia successo un'altra richiesta pendente,
facendo fallire la riacquisizione.

Si tenga presente infine che \func{flock} non è supportata per i file
mantenuti su NFS, in questo caso, se si ha la necessità di utilizzare il
\textit{file locking}, occorre usare l'interfaccia del \textit{file locking}
POSIX basata su \func{fcntl} che è in grado di funzionare anche attraverso
NFS, a condizione ovviamente che sia il client che il server supportino questa
funzionalità.

La semantica del \textit{file locking} di BSD inoltre è diversa da quella del
\textit{file locking} POSIX, in particolare per quanto riguarda il
comportamento dei \textit{file lock} nei confronti delle due funzioni
\func{dup} e \func{fork}.  Per capire queste differenze occorre descrivere con
maggiore dettaglio come viene realizzato dal kernel il \textit{file locking}
per entrambe le interfacce.

In fig.~\ref{fig:file_flock_struct} si è riportato uno schema essenziale
dell'implementazione del \textit{file locking} in stile BSD su Linux. Il punto
fondamentale da capire è che un \textit{file lock}, qualunque sia
l'interfaccia che si usa, anche se richiesto attraverso un file descriptor,
agisce sempre su di un file; perciò le informazioni relative agli eventuali
\textit{file lock} sono mantenute dal kernel a livello di \textit{inode}, dato
che questo è l'unico riferimento in comune che possono avere due processi
diversi che aprono lo stesso file.

In particolare, come accennato in fig.~\ref{fig:file_flock_struct}, i
\textit{file lock} sono mantenuti in una \textit{linked list} di strutture
\kstructd{file\_lock}. La lista è referenziata dall'indirizzo di partenza
mantenuto dal campo \var{i\_flock} della struttura \kstruct{inode} (per le
definizioni esatte si faccia riferimento al file \file{include/linux/fs.h} nei
sorgenti del kernel).  Un bit del campo \var{fl\_flags} di specifica se si
tratta di un lock in semantica BSD (\constd{FL\_FLOCK}) o POSIX
(\constd{FL\_POSIX}) o un \textit{file lease} (\constd{FL\_LEASE}, vedi
sez.~\ref{sec:file_asyncronous_lease}).

\begin{figure}[!htb]
  \centering
  \includegraphics[width=12cm]{img/file_flock}
  \caption{Schema dell'architettura del \textit{file locking}, nel caso
    particolare del suo utilizzo da parte dalla funzione \func{flock}.}
  \label{fig:file_flock_struct}
\end{figure}

La richiesta di un \textit{file lock} prevede una scansione della lista per
determinare se l'acquisizione è possibile, ed in caso positivo l'aggiunta di
un nuovo elemento (cioè l'aggiunta di una nuova struttura
\kstruct{file\_lock}).  Nel caso dei blocchi creati con \func{flock} la
semantica della funzione prevede che sia \func{dup} che \func{fork} non creino
ulteriori istanze di un \textit{file lock} quanto piuttosto degli ulteriori
riferimenti allo stesso. Questo viene realizzato dal kernel secondo lo schema
di fig.~\ref{fig:file_flock_struct}, associando ad ogni nuovo \textit{file
  lock} un puntatore alla voce nella \textit{file table} da cui si è richiesto
il blocco, che così ne identifica il titolare. Il puntatore è mantenuto nel
campo \var{fl\_file} di \kstruct{file\_lock}, e viene utilizzato solo per i
\textit{file lock} creati con la semantica BSD.

Questa struttura prevede che, quando si richiede la rimozione di un
\textit{file lock}, il kernel acconsenta solo se la richiesta proviene da un
file descriptor che fa riferimento ad una voce nella \textit{file table}
corrispondente a quella registrata nel blocco.  Allora se ricordiamo quanto
visto in sez.~\ref{sec:file_dup} e sez.~\ref{sec:file_shared_access}, e cioè
che i file descriptor duplicati e quelli ereditati in un processo figlio
puntano sempre alla stessa voce nella \textit{file table}, si può capire
immediatamente quali sono le conseguenze nei confronti delle funzioni
\func{dup} e \func{fork}.

Sarà così possibile rimuovere un \textit{file lock} attraverso uno qualunque
dei file descriptor che fanno riferimento alla stessa voce nella \textit{file
  table}, anche se questo è diverso da quello con cui lo si è
creato,\footnote{attenzione, questo non vale se il file descriptor fa
  riferimento allo stesso file, ma attraverso una voce diversa della
  \textit{file table}, come accade tutte le volte che si apre più volte lo
  stesso file.} o se si esegue la rimozione in un processo figlio. Inoltre una
volta tolto un \textit{file lock} su un file, la rimozione avrà effetto su
tutti i file descriptor che condividono la stessa voce nella \textit{file
  table}, e quindi, nel caso di file descriptor ereditati attraverso una
\func{fork}, anche per processi diversi.

Infine, per evitare che la terminazione imprevista di un processo lasci attivi
dei \textit{file lock}, quando un file viene chiuso il kernel provvede anche a
rimuovere tutti i blocchi ad esso associati. Anche in questo caso occorre
tenere presente cosa succede quando si hanno file descriptor duplicati; in tal
caso infatti il file non verrà effettivamente chiuso (ed il blocco rimosso)
fintanto che non viene rilasciata la relativa voce nella \textit{file table};
e questo avverrà solo quando tutti i file descriptor che fanno riferimento
alla stessa voce sono stati chiusi.  Quindi, nel caso ci siano duplicati o
processi figli che mantengono ancora aperto un file descriptor, il
\textit{file lock} non viene rilasciato.
 

\subsection{Il \textit{file locking} POSIX}
\label{sec:file_posix_lock}

La seconda interfaccia per l'\textit{advisory locking} disponibile in Linux è
quella standardizzata da POSIX, basata sulla funzione di sistema
\func{fcntl}. Abbiamo già trattato questa funzione nelle sue molteplici
possibilità di utilizzo in sez.~\ref{sec:file_fcntl_ioctl}. Quando la si
impiega per il \textit{file locking} essa viene usata solo secondo il seguente
prototipo:

\begin{funcproto}{
\fhead{fcntl.h}
\fdecl{int fcntl(int fd, int cmd, struct flock *lock)}
\fdesc{Applica o rimuove un \textit{file lock}.} 
}

{La funzione ritorna $0$ in caso di successo e $-1$ per un errore, nel qual
  caso \var{errno} assumerà uno dei valori: 
  \begin{errlist}
    \item[\errcode{EACCES}] l'operazione è proibita per la presenza di
      \textit{file lock} da parte di altri processi.
    \item[\errcode{EDEADLK}] si è richiesto un \textit{lock} su una regione
      bloccata da un altro processo che è a sua volta in attesa dello sblocco
      di un \textit{lock} mantenuto dal processo corrente; si avrebbe pertanto
      un \textit{deadlock}. Non è garantito che il sistema riconosca sempre
      questa situazione.
    \item[\errcode{EINTR}] la funzione è stata interrotta da un segnale prima
      di poter acquisire un \textit{file lock}.
    \item[\errcode{ENOLCK}] il sistema non ha le risorse per il blocco: ci
      sono troppi segmenti di \textit{lock} aperti, si è esaurita la tabella
      dei \textit{file lock}, o il protocollo per il blocco remoto è fallito.
  \end{errlist}
  ed inoltre \errval{EBADF}, \errval{EFAULT} nel loro significato generico.}
\end{funcproto}

Al contrario di quanto avviene con l'interfaccia basata su \func{flock} con
\func{fcntl} è possibile bloccare anche delle singole sezioni di un file, fino
al singolo byte. Inoltre la funzione permette di ottenere alcune informazioni
relative agli eventuali blocchi preesistenti.  Per poter fare tutto questo la
funzione utilizza come terzo argomento una apposita struttura \struct{flock}
(la cui definizione è riportata in fig.~\ref{fig:struct_flock}) nella quale
inserire tutti i dati relativi ad un determinato blocco. Si tenga presente poi
che un \textit{file lock} fa sempre riferimento ad una regione, per cui si
potrà avere un conflitto anche se c'è soltanto una sovrapposizione parziale
con un'altra regione bloccata.

\begin{figure}[!htb]
  \footnotesize \centering
  \begin{minipage}[c]{0.90\textwidth}
    \includestruct{listati/flock.h}
  \end{minipage} 
  \normalsize 
  \caption{La struttura \structd{flock}, usata da \func{fcntl} per il
    \textit{file locking}.}
  \label{fig:struct_flock}
\end{figure}

I primi tre campi della struttura, \var{l\_whence}, \var{l\_start} e
\var{l\_len}, servono a specificare la sezione del file a cui fa riferimento
il blocco: \var{l\_start} specifica il byte di partenza, \var{l\_len} la
lunghezza della sezione e infine \var{l\_whence} imposta il riferimento da cui
contare \var{l\_start}. Il valore di \var{l\_whence} segue la stessa semantica
dell'omonimo argomento di \func{lseek}, coi tre possibili valori
\const{SEEK\_SET}, \const{SEEK\_CUR} e \const{SEEK\_END}, (si vedano le
relative descrizioni in tab.~\ref{tab:lseek_whence_values}).

Si tenga presente che un \textit{file lock} può essere richiesto anche per una
regione al di là della corrente fine del file, così che una eventuale
estensione dello stesso resti coperta dal blocco. Inoltre se si specifica un
valore nullo per \var{l\_len} il blocco si considera esteso fino alla
dimensione massima del file; in questo modo è possibile bloccare una qualunque
regione a partire da un certo punto fino alla fine del file, coprendo
automaticamente quanto eventualmente aggiunto in coda allo stesso.

Lo standard POSIX non richiede che \var{l\_len} sia positivo, ed a partire dal
kernel 2.4.21 è possibile anche indicare valori di \var{l\_len} negativi, in
tal caso l'intervallo coperto va da \var{l\_start}$+$\var{l\_len} a
\var{l\_start}$-1$, mentre per un valore positivo l'intervallo va da
\var{l\_start} a \var{l\_start}$+$\var{l\_len}$-1$. Si può però usare un
valore negativo soltanto se l'inizio della regione indicata non cade prima
dell'inizio del file, mentre come accennato con un valore positivo  si
può anche indicare una regione che eccede la dimensione corrente del file.

Il tipo di \textit{file lock} richiesto viene specificato dal campo
\var{l\_type}, esso può assumere i tre valori definiti dalle costanti
riportate in tab.~\ref{tab:file_flock_type}, che permettono di richiedere
rispettivamente uno \textit{shared lock}, un \textit{esclusive lock}, e la
rimozione di un blocco precedentemente acquisito. Infine il campo \var{l\_pid}
viene usato solo in caso di lettura, quando si chiama \func{fcntl} con
\const{F\_GETLK}, e riporta il \ids{PID} del processo che detiene il
\textit{file lock}.

\begin{table}[htb]
  \centering
  \footnotesize
  \begin{tabular}[c]{|l|l|}
    \hline
    \textbf{Valore} & \textbf{Significato} \\
    \hline
    \hline
    \constd{F\_RDLCK} & Richiede un blocco condiviso (\textit{read lock}).\\
    \constd{F\_WRLCK} & Richiede un blocco esclusivo (\textit{write lock}).\\
    \constd{F\_UNLCK} & Richiede l'eliminazione di un \textit{file lock}.\\
    \hline    
  \end{tabular}
  \caption{Valori possibili per il campo \var{l\_type} di \struct{flock}.}
  \label{tab:file_flock_type}
\end{table}

Oltre a quanto richiesto tramite i campi di \struct{flock}, l'operazione
effettivamente svolta dalla funzione è stabilita dal valore dall'argomento
\param{cmd} che, come già riportato in sez.~\ref{sec:file_fcntl_ioctl},
specifica l'azione da compiere; i valori utilizzabili relativi al \textit{file
  locking} sono tre:
\begin{basedescript}{\desclabelwidth{2.0cm}}
\item[\constd{F\_GETLK}] verifica se il \textit{file lock} specificato dalla
  struttura puntata da \param{lock} può essere acquisito: in caso negativo
  sovrascrive la struttura \param{flock} con i valori relativi al blocco già
  esistente che ne blocca l'acquisizione, altrimenti si limita a impostarne il
  campo \var{l\_type} con il valore \const{F\_UNLCK}.
\item[\constd{F\_SETLK}] se il campo \var{l\_type} della struttura puntata da
  \param{lock} è \const{F\_RDLCK} o \const{F\_WRLCK} richiede il
  corrispondente \textit{file lock}, se è \const{F\_UNLCK} lo rilascia; nel
  caso la richiesta non possa essere soddisfatta a causa di un blocco
  preesistente la funzione ritorna immediatamente con un errore di
  \errcode{EACCES} o di \errcode{EAGAIN}.
\item[\constd{F\_SETLKW}] è identica a \const{F\_SETLK}, ma se la richiesta di
  non può essere soddisfatta per la presenza di un altro blocco, mette il
  processo in stato di attesa fintanto che il blocco precedente non viene
  rilasciato; se l'attesa viene interrotta da un segnale la funzione ritorna
  con un errore di \errcode{EINTR}.
\end{basedescript}

Si noti che per quanto detto il comando \const{F\_GETLK} non serve a rilevare
una presenza generica di blocco su un file, perché se ne esistono altri
compatibili con quello richiesto, la funzione ritorna comunque impostando
\var{l\_type} a \const{F\_UNLCK}.  Inoltre a seconda del valore di
\var{l\_type} si potrà controllare o l'esistenza di un qualunque tipo di
blocco (se è \const{F\_WRLCK}) o di \textit{write lock} (se è
\const{F\_RDLCK}). Si consideri poi che può esserci più di un blocco che
impedisce l'acquisizione di quello richiesto (basta che le regioni si
sovrappongano), ma la funzione ne riporterà sempre soltanto uno, impostando
\var{l\_whence} a \const{SEEK\_SET} ed i valori \var{l\_start} e \var{l\_len}
per indicare quale è la regione bloccata.

Infine si tenga presente che effettuare un controllo con il comando
\const{F\_GETLK} e poi tentare l'acquisizione con \const{F\_SETLK} non è una
operazione atomica (un altro processo potrebbe acquisire un blocco fra le due
chiamate) per cui si deve sempre verificare il codice di ritorno di
\func{fcntl}\footnote{controllare il codice di ritorno delle funzioni invocate
  è comunque una buona norma di programmazione, che permette di evitare un
  sacco di errori difficili da tracciare proprio perché non vengono rilevati.}
quando la si invoca con \const{F\_SETLK}, per controllare che il blocco sia
stato effettivamente acquisito.

\begin{figure}[!htb]
  \centering \includegraphics[width=9cm]{img/file_lock_dead}
  \caption{Schema di una situazione di \textit{deadlock}.}
  \label{fig:file_flock_dead}
\end{figure}

Non operando a livello di interi file, il \textit{file locking} POSIX
introduce un'ulteriore complicazione; consideriamo la situazione illustrata in
fig.~\ref{fig:file_flock_dead}, in cui il processo A blocca la regione 1 e il
processo B la regione 2. Supponiamo che successivamente il processo A richieda
un lock sulla regione 2 che non può essere acquisito per il preesistente lock
del processo 2; il processo 1 si bloccherà fintanto che il processo 2 non
rilasci il blocco. Ma cosa accade se il processo 2 nel frattempo tenta a sua
volta di ottenere un lock sulla regione A? Questa è una tipica situazione che
porta ad un \textit{deadlock}, dato che a quel punto anche il processo 2 si
bloccherebbe, e niente potrebbe sbloccare l'altro processo.  Per questo motivo
il kernel si incarica di rilevare situazioni di questo tipo, ed impedirle
restituendo un errore di \errcode{EDEADLK} alla funzione che cerca di
acquisire un blocco che porterebbe ad un \textit{deadlock}.

Per capire meglio il funzionamento del \textit{file locking} in semantica
POSIX (che differisce alquanto rispetto da quello di BSD, visto
sez.~\ref{sec:file_flock}) esaminiamo più in dettaglio come viene gestito dal
kernel. Lo schema delle strutture utilizzate è riportato in
fig.~\ref{fig:file_posix_lock}; come si vede esso è molto simile all'analogo
di fig.~\ref{fig:file_flock_struct}. In questo caso nella figura si sono
evidenziati solo i campi di \kstructd{file\_lock} significativi per la
semantica POSIX, in particolare adesso ciascuna struttura contiene, oltre al
\ids{PID} del processo in \var{fl\_pid}, la sezione di file che viene bloccata
grazie ai campi \var{fl\_start} e \var{fl\_end}.  La struttura è comunque la
stessa, solo che in questo caso nel campo \var{fl\_flags} è impostato il bit
\const{FL\_POSIX} ed il campo \var{fl\_file} non viene usato. Il blocco è
sempre associato all'\textit{inode}, solo che in questo caso la titolarità non
viene identificata con il riferimento ad una voce nella \textit{file table},
ma con il valore del \ids{PID} del processo.

\begin{figure}[!htb]
  \centering \includegraphics[width=12cm]{img/file_posix_lock}
  \caption{Schema dell'architettura del \textit{file locking}, nel caso
    particolare del suo utilizzo secondo l'interfaccia standard POSIX.}
  \label{fig:file_posix_lock}
\end{figure}

Quando si richiede un \textit{file lock} il kernel effettua una scansione di
tutti i blocchi presenti sul file\footnote{scandisce cioè la \textit{linked
    list} delle strutture \kstruct{file\_lock}, scartando automaticamente
  quelle per cui \var{fl\_flags} non è \const{FL\_POSIX}, così che le due
  interfacce restano ben separate.}  per verificare se la regione richiesta
non si sovrappone ad una già bloccata, in caso affermativo decide in base al
tipo di blocco, in caso negativo il nuovo blocco viene comunque acquisito ed
aggiunto alla lista.

Nel caso di rimozione invece questa viene effettuata controllando che il
\ids{PID} del processo richiedente corrisponda a quello contenuto nel blocco.
Questa diversa modalità ha delle conseguenze precise riguardo il comportamento
dei \textit{file lock} POSIX. La prima conseguenza è che un \textit{file lock}
POSIX non viene mai ereditato attraverso una \func{fork}, dato che il processo
figlio avrà un \ids{PID} diverso, mentre passa indenne attraverso una
\func{exec} in quanto il \ids{PID} resta lo stesso.  Questo comporta che, al
contrario di quanto avveniva con la semantica BSD, quando un processo termina
tutti i \textit{file lock} da esso detenuti vengono immediatamente rilasciati.

La seconda conseguenza è che qualunque file descriptor che faccia riferimento
allo stesso file (che sia stato ottenuto con una \func{dup} o con una
\func{open} in questo caso non fa differenza) può essere usato per rimuovere
un blocco, dato che quello che conta è solo il \ids{PID} del processo. Da
questo deriva una ulteriore sottile differenza di comportamento: dato che alla
chiusura di un file i blocchi ad esso associati vengono rimossi, nella
semantica POSIX basterà chiudere un file descriptor qualunque per cancellare
tutti i blocchi relativi al file cui esso faceva riferimento, anche se questi
fossero stati creati usando altri file descriptor che restano aperti.

Dato che il controllo sull'accesso ai blocchi viene eseguito sulla base del
\ids{PID} del processo, possiamo anche prendere in considerazione un altro
degli aspetti meno chiari di questa interfaccia e cioè cosa succede quando si
richiedono dei blocchi su regioni che si sovrappongono fra loro all'interno
stesso processo. Siccome il controllo, come nel caso della rimozione, si basa
solo sul \ids{PID} del processo che chiama la funzione, queste richieste
avranno sempre successo.  Nel caso della semantica BSD, essendo i lock
relativi a tutto un file e non accumulandosi,\footnote{questa ultima
  caratteristica è vera in generale, se cioè si richiede più volte lo stesso
  \textit{file lock}, o più blocchi sulla stessa sezione di file, le richieste
  non si cumulano e basta una sola richiesta di rilascio per cancellare il
  blocco.}  la cosa non ha alcun effetto; la funzione ritorna con successo,
senza che il kernel debba modificare la lista dei \textit{file lock}.

Con i \textit{file lock} POSIX invece si possono avere una serie di situazioni
diverse: ad esempio è possibile rimuovere con una sola chiamata più
\textit{file lock} distinti (indicando in una regione che si sovrapponga
completamente a quelle di questi ultimi), o rimuovere solo una parte di un
blocco preesistente (indicando una regione contenuta in quella di un altro
blocco), creando un buco, o coprire con un nuovo blocco altri \textit{file
  lock} già ottenuti, e così via, a secondo di come si sovrappongono le
regioni richieste e del tipo di operazione richiesta.

Il comportamento seguito in questo caso è che la funzione ha successo ed
esegue l'operazione richiesta sulla regione indicata; è compito del kernel
preoccuparsi di accorpare o dividere le voci nella lista dei \textit{file
  lock} per far sì che le regioni bloccate da essa risultanti siano coerenti
con quanto necessario a soddisfare l'operazione richiesta.

\begin{figure}[!htbp]
  \footnotesize \centering
  \begin{minipage}[c]{\codesamplewidth}
    \includecodesample{listati/Flock.c}
  \end{minipage}
  \normalsize 
  \caption{Sezione principale del codice del programma \file{Flock.c}.}
  \label{fig:file_flock_code}
\end{figure}

Per fare qualche esempio sul \textit{file locking} si è scritto un programma che
permette di bloccare una sezione di un file usando la semantica POSIX, o un
intero file usando la semantica BSD; in fig.~\ref{fig:file_flock_code} è
riportata il corpo principale del codice del programma, (il testo completo è
allegato nella directory dei sorgenti, nel file \texttt{Flock.c}).

La sezione relativa alla gestione delle opzioni al solito si è omessa, come la
funzione che stampa le istruzioni per l'uso del programma, essa si cura di
impostare le variabili \var{type}, \var{start} e \var{len}; queste ultime due
vengono inizializzate al valore numerico fornito rispettivamente tramite gli
switch \code{-s} e \cmd{-l}, mentre il valore della prima viene impostato con
le opzioni \cmd{-w} e \cmd{-r} si richiede rispettivamente o un \textit{write
  lock} o \textit{read lock} (i due valori sono esclusivi, la variabile
assumerà quello che si è specificato per ultimo). Oltre a queste tre vengono
pure impostate la variabile \var{bsd}, che abilita la semantica omonima quando
si invoca l'opzione \cmd{-f} (il valore preimpostato è nullo, ad indicare la
semantica POSIX), e la variabile \var{cmd} che specifica la modalità di
richiesta del \textit{file lock} (bloccante o meno), a seconda dell'opzione
\cmd{-b}.

Il programma inizia col controllare (\texttt{\small 11-14}) che venga passato
un argomento (il file da bloccare), che sia stato scelto (\texttt{\small
  15-18}) il tipo di blocco, dopo di che apre (\texttt{\small 19}) il file,
uscendo (\texttt{\small 20-23}) in caso di errore. A questo punto il
comportamento dipende dalla semantica scelta; nel caso sia BSD occorre
reimpostare il valore di \var{cmd} per l'uso con \func{flock}; infatti il
valore preimpostato fa riferimento alla semantica POSIX e vale rispettivamente
\const{F\_SETLKW} o \const{F\_SETLK} a seconda che si sia impostato o meno la
modalità bloccante.

Nel caso si sia scelta la semantica BSD (\texttt{\small 25-34}) prima si
controlla (\texttt{\small 27-31}) il valore di \var{cmd} per determinare se
si vuole effettuare una chiamata bloccante o meno, reimpostandone il valore
opportunamente, dopo di che a seconda del tipo di blocco al valore viene
aggiunta la relativa opzione, con un OR aritmetico, dato che \func{flock}
vuole un argomento \param{operation} in forma di maschera binaria.  Nel caso
invece che si sia scelta la semantica POSIX le operazioni sono molto più
immediate si prepara (\texttt{\small 36-40}) la struttura per il lock, e lo
si esegue (\texttt{\small 41}).

In entrambi i casi dopo aver richiesto il blocco viene controllato il
risultato uscendo (\texttt{\small 44-46}) in caso di errore, o stampando un
messaggio (\texttt{\small 47-49}) in caso di successo. Infine il programma si
pone in attesa (\texttt{\small 50}) finché un segnale (ad esempio un \cmd{C-c}
dato da tastiera) non lo interrompa; in questo caso il programma termina, e
tutti i blocchi vengono rilasciati.

Con il programma possiamo fare varie verifiche sul funzionamento del
\textit{file locking}; cominciamo con l'eseguire un \textit{read lock} su un
file, ad esempio usando all'interno di un terminale il seguente comando:

\begin{Console}
[piccardi@gont sources]$ \textbf{./flock -r Flock.c}
Lock acquired
\end{Console}
%$
il programma segnalerà di aver acquisito un blocco e si bloccherà; in questo
caso si è usato il \textit{file locking} POSIX e non avendo specificato niente
riguardo alla sezione che si vuole bloccare sono stati usati i valori
preimpostati che bloccano tutto il file. A questo punto se proviamo ad
eseguire lo stesso comando in un altro terminale, e avremo lo stesso
risultato. Se invece proviamo ad eseguire un \textit{write lock} avremo:

\begin{Console}
[piccardi@gont sources]$ \textbf{./flock -w Flock.c}
Failed lock: Resource temporarily unavailable
\end{Console}
%$
come ci aspettiamo il programma terminerà segnalando l'indisponibilità del
blocco, dato che il file è bloccato dal precedente \textit{read lock}. Si noti
che il risultato è lo stesso anche se si richiede il blocco su una sola parte
del file con il comando:

\begin{Console}
[piccardi@gont sources]$ \textbf{./flock -w -s0 -l10 Flock.c}
Failed lock: Resource temporarily unavailable
\end{Console}
%$
se invece blocchiamo una regione con: 

\begin{Console}
[piccardi@gont sources]$ \textbf{./flock -r -s0 -l10 Flock.c}
Lock acquired
\end{Console}
%$
una volta che riproviamo ad acquisire il \textit{write lock} i risultati
dipenderanno dalla regione richiesta; ad esempio nel caso in cui le due
regioni si sovrappongono avremo che:

\begin{Console}
[piccardi@gont sources]$ \textbf{./flock -w -s5 -l15  Flock.c}
Failed lock: Resource temporarily unavailable
\end{Console}
%$
ed il blocco viene rifiutato, ma se invece si richiede una regione distinta
avremo che:

\begin{Console}
[piccardi@gont sources]$ \textbf{./flock -w -s11 -l15  Flock.c}
Lock acquired
\end{Console}
%$
ed il blocco viene acquisito. Se a questo punto si prova ad eseguire un
\textit{read lock} che comprende la nuova regione bloccata in scrittura:

\begin{Console}
[piccardi@gont sources]$ \textbf{./flock -r -s10 -l20 Flock.c}
Failed lock: Resource temporarily unavailable
\end{Console}
%$
come ci aspettiamo questo non sarà consentito.

Il programma di norma esegue il tentativo di acquisire il lock in modalità non
bloccante, se però usiamo l'opzione \cmd{-b} possiamo impostare la modalità
bloccante, riproviamo allora a ripetere le prove precedenti con questa
opzione:

\begin{Console}
[piccardi@gont sources]$ \textbf{./flock -r -b -s0 -l10 Flock.c} Lock acquired
\end{Console}
%$
il primo comando acquisisce subito un \textit{read lock}, e quindi non cambia
nulla, ma se proviamo adesso a richiedere un \textit{write lock} che non potrà
essere acquisito otterremo:

\begin{Console}
[piccardi@gont sources]$ \textbf{./flock -w -s0 -l10 Flock.c}
\end{Console}
%$
il programma cioè si bloccherà nella chiamata a \func{fcntl}; se a questo
punto rilasciamo il precedente blocco (terminando il primo comando un
\texttt{C-c} sul terminale) potremo verificare che sull'altro terminale il
blocco viene acquisito, con la comparsa di una nuova riga:

\begin{Console}
[piccardi@gont sources]$ \textbf{./flock -w -s0 -l10 Flock.c}
Lock acquired
\end{Console}
%$

Un'altra cosa che si può controllare con il nostro programma è l'interazione
fra i due tipi di blocco; se ripartiamo dal primo comando con cui si è
ottenuto un blocco in lettura sull'intero file, possiamo verificare cosa
succede quando si cerca di ottenere un blocco in scrittura con la semantica
BSD:

\begin{Console}
[root@gont sources]# \textbf{./flock -f -w Flock.c}
Lock acquired
\end{Console}
%$
che ci mostra come i due tipi di blocco siano assolutamente indipendenti; per
questo motivo occorre sempre tenere presente quale, fra le due semantiche
disponibili, stanno usando i programmi con cui si interagisce, dato che i
blocchi applicati con l'altra non avrebbero nessun effetto.

% \subsection{La funzione \func{lockf}}
% \label{sec:file_lockf}

Abbiamo visto come l'interfaccia POSIX per il \textit{file locking} sia molto
più potente e flessibile di quella di BSD, questo comporta anche una maggiore
complessità per via delle varie opzioni da passare a \func{fcntl}. Per questo
motivo è disponibile anche una interfaccia semplificata che utilizza la
funzione \funcd{lockf},\footnote{la funzione è ripresa da System V e per
  poterla utilizzare è richiesta che siano definite le opportune macro, una
  fra \macro{\_BSD\_SOURCE} o \macro{\_SVID\_SOURCE}, oppure
  \macro{\_XOPEN\_SOURCE} ad un valore di almeno 500, oppure
  \macro{\_XOPEN\_SOURCE} e \macro{\_XOPEN\_SOURCE\_EXTENDED}.} il cui
prototipo è:

\begin{funcproto}{
\fhead{unistd.h}
\fdecl{int lockf(int fd, int cmd, off\_t len)}
\fdesc{Applica, controlla o rimuove un \textit{file lock}.} 
}

{La funzione ritorna $0$ in caso di successo e $-1$ per un errore, nel qual
  caso \var{errno} assumerà uno dei valori: 
  \begin{errlist}
  \item[\errcode{EAGAIN}] il file è bloccato, e si sono richiesti
    \const{F\_TLOCK} o \const{F\_TEST} (in alcuni casi può dare anche
    \errcode{EACCESS}.
  \item[\errcode{EBADF}] \param{fd} non è un file descriptor aperto o si sono
    richiesti \const{F\_LOCK} o \const{F\_TLOCK} ma il file non è scrivibile.
  \item[\errcode{EINVAL}] si è usato un valore non valido per \param{cmd}.
  \end{errlist}
  ed inoltre \errcode{EDEADLK} e \errcode{ENOLCK} con lo stesso significato
  che hanno con \func{fcntl}.
}
\end{funcproto}
  
La funzione opera sul file indicato dal file descriptor \param{fd}, che deve
essere aperto in scrittura, perché utilizza soltanto \textit{lock}
esclusivi. La sezione di file bloccata viene controllata dal valore
di \param{len}, che indica la lunghezza della stessa, usando come riferimento
la posizione corrente sul file. La sezione effettiva varia a secondo del
segno, secondo lo schema illustrato in fig.~\ref{fig:file_lockf_boundary}, se
si specifica un valore nullo il file viene bloccato a partire dalla posizione
corrente fino alla sua fine presente o futura (nello schema corrisponderebbe
ad un valore infinito positivo).

\begin{figure}[!htb] 
  \centering
  \includegraphics[width=10cm]{img/lockf_boundary}
  \caption{Schema della sezione di file bloccata con \func{lockf}.}
  \label{fig:file_lockf_boundary}
\end{figure}

Il comportamento della funzione viene controllato dal valore
dell'argomento \param{cmd}, che specifica quale azione eseguire, i soli valori
consentiti sono i seguenti:

\begin{basedescript}{\desclabelwidth{2.0cm}}
\item[\constd{F\_LOCK}] Richiede un \textit{lock} esclusivo sul file, e blocca
  il processo chiamante se, anche parzialmente, la sezione indicata si
  sovrappone ad una che è già stata bloccata da un altro processo; in caso di
  sovrapposizione con un altro blocco già ottenuto le sezioni vengono unite.
\item[\constd{F\_TLOCK}] Richiede un \textit{exclusive lock}, in maniera
  identica a \const{F\_LOCK}, ma in caso di indisponibilità non blocca il
  processo restituendo un errore di \errval{EAGAIN}.
\item[\constd{F\_ULOCK}] Rilascia il blocco sulla sezione indicata, questo può
  anche causare la suddivisione di una sezione bloccata in precedenza nelle
  due parti eccedenti nel caso si sia indicato un intervallo più limitato.
\item[\constd{F\_TEST}] Controlla la presenza di un blocco sulla sezione di
  file indicata, \func{lockf} ritorna $0$ se la sezione è libera o bloccata
  dal processo stesso, o $-1$ se è bloccata da un altro processo, nel qual
  caso \var{errno} assume il valore \errval{EAGAIN} (ma su alcuni sistemi può
  essere restituito anche \errval{EACCESS}).
\end{basedescript}

La funzione è semplicemente una diversa interfaccia al \textit{file locking}
POSIX ed è realizzata utilizzando \func{fcntl}; pertanto la semantica delle
operazioni è la stessa di quest'ultima e quindi la funzione presenta lo stesso
comportamento riguardo gli effetti della chiusura dei file, ed il
comportamento sui file duplicati e nel passaggio attraverso \func{fork} ed
\func{exec}. Per questo stesso motivo la funzione non è equivalente a
\func{flock} e può essere usata senza interferenze insieme a quest'ultima.

% TODO trattare i POSIX file-private lock introdotti con il 3.15, 
% vedi http://lwn.net/Articles/586904/ correlato:
% http://www.samba.org/samba/news/articles/low_point/tale_two_stds_os2.html 

\subsection{Il \textit{mandatory locking}}
\label{sec:file_mand_locking}

\itindbeg{mandatory~locking}

Il \textit{mandatory locking} è una opzione introdotta inizialmente in SVr4,
per introdurre un \textit{file locking} che, come dice il nome, fosse
effettivo indipendentemente dai controlli eseguiti da un processo. Con il
\textit{mandatory locking} infatti è possibile far eseguire il blocco del file
direttamente al sistema, così che, anche qualora non si predisponessero le
opportune verifiche nei processi, questo verrebbe comunque rispettato.

Per poter utilizzare il \textit{mandatory locking} è stato introdotto un
utilizzo particolare del bit \acr{sgid} dei permessi dei file. Se si ricorda
quanto esposto in sez.~\ref{sec:file_special_perm}), esso viene di norma
utilizzato per cambiare il \ids{GID} effettivo con cui viene eseguito un
programma, ed è pertanto sempre associato alla presenza del permesso di
esecuzione per il gruppo. Impostando questo bit su un file senza permesso di
esecuzione in un sistema che supporta il \textit{mandatory locking}, fa sì che
quest'ultimo venga attivato per il file in questione. In questo modo una
combinazione dei permessi originariamente non contemplata, in quanto senza
significato, diventa l'indicazione della presenza o meno del \textit{mandatory
  locking}.\footnote{un lettore attento potrebbe ricordare quanto detto in
  sez.~\ref{sec:file_perm_management} e cioè che il bit \acr{sgid} viene
  cancellato (come misura di sicurezza) quando di scrive su un file, questo
  non vale quando esso viene utilizzato per attivare il \textit{mandatory
    locking}.}

L'uso del \textit{mandatory locking} presenta vari aspetti delicati, dato che
neanche l'amministratore può passare sopra ad un \textit{file lock}; pertanto
un processo che blocchi un file cruciale può renderlo completamente
inaccessibile, rendendo completamente inutilizzabile il sistema\footnote{il
  problema si potrebbe risolvere rimuovendo il bit \acr{sgid}, ma non è detto
  che sia così facile fare questa operazione con un sistema bloccato.}
inoltre con il \textit{mandatory locking} si può bloccare completamente un
server NFS richiedendo una lettura su un file su cui è attivo un blocco. Per
questo motivo l'abilitazione del \textit{mandatory locking} è di norma
disabilitata, e deve essere attivata filesystem per filesystem in fase di
montaggio, specificando l'apposita opzione di \func{mount} riportata in
sez.~\ref{sec:filesystem_mounting}, o con l'opzione \code{-o mand} per il
comando omonimo.

Si tenga presente inoltre che il \textit{mandatory locking} funziona solo
sull'interfaccia POSIX di \func{fcntl}. Questo ha due conseguenze: che non si
ha nessun effetto sui \textit{file lock} richiesti con l'interfaccia di
\func{flock}, e che la granularità del blocco è quella del singolo byte, come
per \func{fcntl}.

La sintassi di acquisizione dei blocchi è esattamente la stessa vista in
precedenza per \func{fcntl} e \func{lockf}, la differenza è che in caso di
\textit{mandatory lock} attivato non è più necessario controllare la
disponibilità di accesso al file, ma si potranno usare direttamente le
ordinarie funzioni di lettura e scrittura e sarà compito del kernel gestire
direttamente il \textit{file locking}.

Questo significa che in caso di \textit{read lock} la lettura dal file potrà
avvenire normalmente con \func{read}, mentre una \func{write} si bloccherà
fino al rilascio del blocco, a meno di non aver aperto il file con
\const{O\_NONBLOCK}, nel qual caso essa ritornerà immediatamente con un errore
di \errcode{EAGAIN}.

Se invece si è acquisito un \textit{write lock} tutti i tentativi di leggere o
scrivere sulla regione del file bloccata fermeranno il processo fino al
rilascio del blocco, a meno che il file non sia stato aperto con
\const{O\_NONBLOCK}, nel qual caso di nuovo si otterrà un ritorno immediato
con l'errore di \errcode{EAGAIN}.

Infine occorre ricordare che le funzioni di lettura e scrittura non sono le
sole ad operare sui contenuti di un file, e che sia \func{creat} che
\func{open} (quando chiamata con \const{O\_TRUNC}) effettuano dei cambiamenti,
così come \func{truncate}, riducendone le dimensioni (a zero nei primi due
casi, a quanto specificato nel secondo). Queste operazioni sono assimilate a
degli accessi in scrittura e pertanto non potranno essere eseguite (fallendo
con un errore di \errcode{EAGAIN}) su un file su cui sia presente un qualunque
blocco (le prime due sempre, la terza solo nel caso che la riduzione delle
dimensioni del file vada a sovrapporsi ad una regione bloccata).

L'ultimo aspetto della interazione del \textit{mandatory locking} con le
funzioni di accesso ai file è quello relativo ai file mappati in memoria (vedi
sez.~\ref{sec:file_memory_map}); anche in tal caso infatti, quando si esegue
la mappatura con l'opzione \const{MAP\_SHARED}, si ha un accesso al contenuto
del file. Lo standard SVID prevede che sia impossibile eseguire il
\textit{memory mapping} di un file su cui sono presenti dei
blocchi\footnote{alcuni sistemi, come HP-UX, sono ancora più restrittivi e lo
  impediscono anche in caso di \textit{advisory locking}, anche se questo
  comportamento non ha molto senso, dato che comunque qualunque accesso
  diretto al file è consentito.} in Linux è stata però fatta la scelta
implementativa\footnote{per i dettagli si possono leggere le note relative
  all'implementazione, mantenute insieme ai sorgenti del kernel nel file
  \file{Documentation/mandatory.txt}.}  di seguire questo comportamento
soltanto quando si chiama \func{mmap} con l'opzione \const{MAP\_SHARED} (nel
qual caso la funzione fallisce con il solito \errcode{EAGAIN}) che comporta la
possibilità di modificare il file.

Si tenga conto infine che su Linux l'implementazione corrente del
\textit{mandatory locking} è difettosa e soffre di una \textit{race
  condition}, per cui una scrittura con \func{write} che si sovrapponga alla
richiesta di un \textit{read lock} può modificare i dati anche dopo che questo
è stato ottenuto, ed una lettura con \func{read} può restituire dati scritti
dopo l'ottenimento di un \textit{write lock}. Lo stesso tipo di problema si
può presentare anche con l'uso di file mappati in memoria; pertanto allo stato
attuale delle cose è sconsigliabile fare affidamento sul \textit{mandatory
  locking}.

% TODO il supporto è stato reso opzionale nel 4.5, verrà eliminato nel futuro
% (vedi http://lwn.net/Articles/667210/)

\itindend{file~locking}

\itindend{mandatory~locking}


\section{L'\textit{I/O multiplexing}}
\label{sec:file_multiplexing}


Uno dei problemi che si presentano quando si deve operare contemporaneamente
su molti file usando le funzioni illustrate in
sez.~\ref{sec:file_unix_interface} e sez.~\ref{sec:files_std_interface} è che
si può essere bloccati nelle operazioni su un file mentre un altro potrebbe
essere disponibile. L'\textit{I/O multiplexing} nasce risposta a questo
problema. In questa sezione forniremo una introduzione a questa problematica
ed analizzeremo le varie funzioni usate per implementare questa modalità di
I/O.


\subsection{La problematica dell'\textit{I/O multiplexing}}
\label{sec:file_noblocking}

Abbiamo visto in sez.~\ref{sec:sig_gen_beha}, affrontando la suddivisione fra
\textit{fast} e \textit{slow} \textit{system call}, che in certi casi le
funzioni di I/O eseguite su un file descriptor possono bloccarsi
indefinitamente. Questo non avviene mai per i file normali, per i quali le
funzioni di lettura e scrittura ritornano sempre subito, ma può avvenire per
alcuni file di dispositivo, come ad esempio una seriale o un terminale, o con
l'uso di file descriptor collegati a meccanismi di intercomunicazione come le
\textit{pipe} (vedi sez.~\ref{sec:ipc_unix}) ed i socket (vedi
sez.~\ref{sec:sock_socket_def}). In casi come questi ad esempio una operazione
di lettura potrebbe bloccarsi se non ci sono dati disponibili sul descrittore
su cui la si sta effettuando.

Questo comportamento è alla radice di una delle problematiche più comuni che
ci si trova ad affrontare nella gestione delle operazioni di I/O: la necessità
di operare su più file descriptor eseguendo funzioni che possono bloccarsi
indefinitamente senza che sia possibile prevedere quando questo può
avvenire. Un caso classico è quello di un server di rete (tratteremo la
problematica in dettaglio nella seconda parte della guida) in attesa di dati
in ingresso prevenienti da vari client.

In un caso di questo tipo, se si andasse ad operare sui vari file descriptor
aperti uno dopo l'altro, potrebbe accadere di restare bloccati nell'eseguire
una lettura su uno di quelli che non è ``\textsl{pronto}'', quando ce ne
potrebbe essere un altro con dati disponibili. Questo comporta nel migliore
dei casi una operazione ritardata inutilmente nell'attesa del completamento di
quella bloccata, mentre nel peggiore dei casi, quando la conclusione
dell'operazione bloccata dipende da quanto si otterrebbe dal file descriptor
``\textsl{disponibile}'', si potrebbe addirittura arrivare ad un
\textit{deadlock}.

\itindbeg{polling}

Abbiamo già accennato in sez.~\ref{sec:file_open_close} che è possibile
prevenire questo tipo di comportamento delle funzioni di I/O aprendo un file
in \textsl{modalità non-bloccante}, attraverso l'uso del flag
\const{O\_NONBLOCK} nella chiamata di \func{open}. In questo caso le funzioni
di lettura o scrittura eseguite sul file che si sarebbero bloccate ritornano
immediatamente, restituendo l'errore \errcode{EAGAIN}.  L'utilizzo di questa
modalità di I/O permette di risolvere il problema controllando a turno i vari
file descriptor, in un ciclo in cui si ripete l'accesso fintanto che esso non
viene garantito. Ovviamente questa tecnica, detta \textit{polling}, è
estremamente inefficiente: si tiene costantemente impiegata la CPU solo per
eseguire in continuazione delle \textit{system call} che nella gran parte dei
casi falliranno.

\itindend{polling}

É appunto per superare questo problema è stato introdotto il concetto di
\textit{I/O multiplexing}, una nuova modalità per la gestione dell'I/O che
consente di tenere sotto controllo più file descriptor in contemporanea,
permettendo di bloccare un processo quando le operazioni di lettura o
scrittura non sono immediatamente effettuabili, e di riprenderne l'esecuzione
una volta che almeno una di quelle che erano state richieste diventi
possibile, in modo da poterla eseguire con la sicurezza di non restare
bloccati.

Dato che, come abbiamo già accennato, per i normali file su disco non si ha
mai un accesso bloccante, l'uso più comune delle funzioni che esamineremo nei
prossimi paragrafi è per i server di rete, in cui esse vengono utilizzate per
tenere sotto controllo dei socket; pertanto ritorneremo su di esse con
ulteriori dettagli e qualche esempio di utilizzo concreto in
sez.~\ref{sec:TCP_sock_multiplexing}.


\subsection{Le funzioni \func{select} e \func{pselect}}
\label{sec:file_select}

Il primo kernel unix-like ad introdurre una interfaccia per l'\textit{I/O
  multiplexing} è stato BSD, con la funzione \funcd{select} che è apparsa in
BSD4.2 ed è stata standardizzata in BSD4.4, in seguito è stata portata su
tutti i sistemi che supportano i socket, compreso le varianti di System V ed
inserita in POSIX.1-2001; il suo prototipo è:\footnote{l'header
  \texttt{sys/select.h} è stato introdotto con POSIX.1-2001, è ed presente con
  la \acr{glibc} a partire dalla versione 2.0, in precedenza, con le
  \acr{libc4} e \acr{libc5}, occorreva includere \texttt{sys/time.h},
  \texttt{sys/types.h} e \texttt{unistd.h}.}

\begin{funcproto}{
\fhead{sys/select.h}
\fdecl{int select(int ndfs, fd\_set *readfds, fd\_set *writefds, fd\_set
    *exceptfds, \\
\phantom{int select(}struct timeval *timeout)}
\fdesc{Attende che uno fra i file descriptor degli insiemi specificati diventi
  attivo.} 
}
{La funzione ritorna $0$ in caso di successo e $-1$ per un errore, nel qual
  caso \var{errno} assumerà uno dei valori: 
  \begin{errlist}
  \item[\errcode{EBADF}] si è specificato un file descriptor non valido
    (chiuso o con errori) in uno degli insiemi.
  \item[\errcode{EINTR}] la funzione è stata interrotta da un segnale.
  \item[\errcode{EINVAL}] si è specificato per \param{ndfs} un valore negativo
    o un valore non valido per \param{timeout}.
  \end{errlist}
  ed inoltre \errval{ENOMEM} nel suo significato generico.}
\end{funcproto}

La funzione mette il processo in stato di \textit{sleep} (vedi
tab.~\ref{tab:proc_proc_states}) fintanto che almeno uno dei file descriptor
degli insiemi specificati (\param{readfds}, \param{writefds} e
\param{exceptfds}), non diventa attivo, per un tempo massimo specificato da
\param{timeout}.

\itindbeg{file~descriptor~set} 

Per specificare quali file descriptor si intende selezionare la funzione usa
un particolare oggetto, il \textit{file descriptor set}, identificato dal tipo
\typed{fd\_set}, che serve ad identificare un insieme di file descriptor, in
maniera analoga a come un \textit{signal set} (vedi sez.~\ref{sec:sig_sigset})
identifica un insieme di segnali. Per la manipolazione di questi \textit{file
  descriptor set} si possono usare delle opportune macro di preprocessore:

{\centering
\vspace{3pt}
\begin{funcbox}{
\fhead{sys/select.h}
\fdecl{void \macrod{FD\_ZERO}(fd\_set *set)}
\fdesc{Inizializza l'insieme (vuoto).} 
\fdecl{void \macrod{FD\_SET}(int fd, fd\_set *set)}
\fdesc{Inserisce il file descriptor \param{fd} nell'insieme.} 
\fdecl{void \macrod{FD\_CLR}(int fd, fd\_set *set)}
\fdesc{Rimuove il file descriptor \param{fd} dall'insieme.} 
\fdecl{int \macrod{FD\_ISSET}(int fd, fd\_set *set)}
\fdesc{Controlla se il file descriptor \param{fd} è nell'insieme.} 
}
\end{funcbox}}


In genere un \textit{file descriptor set} può contenere fino ad un massimo di
\macrod{FD\_SETSIZE} file descriptor.  Questo valore in origine corrispondeva
al limite per il numero massimo di file aperti (ad esempio in Linux, fino alla
serie 2.0.x, c'era un limite di 256 file per processo), ma da quando, nelle
versioni più recenti del kernel, questo limite è stato rimosso, esso indica le
dimensioni massime dei numeri usati nei \textit{file descriptor set}, ed il
suo valore, secondo lo standard POSIX 1003.1-2001, è definito in
\headfile{sys/select.h}, ed è pari a 1024.

Si tenga presente che i \textit{file descriptor set} devono sempre essere
inizializzati con \macro{FD\_ZERO}; passare a \func{select} un valore non
inizializzato può dar luogo a comportamenti non prevedibili. Allo stesso modo
usare \macro{FD\_SET} o \macro{FD\_CLR} con un file descriptor il cui valore
eccede \macro{FD\_SETSIZE} può dare luogo ad un comportamento indefinito.

La funzione richiede di specificare tre insiemi distinti di file descriptor;
il primo, \param{readfds}, verrà osservato per rilevare la disponibilità di
effettuare una lettura,\footnote{per essere precisi la funzione ritornerà in
  tutti i casi in cui la successiva esecuzione di \func{read} risulti non
  bloccante, quindi anche in caso di \textit{end-of-file}.} il secondo,
\param{writefds}, per verificare la possibilità di effettuare una scrittura ed
il terzo, \param{exceptfds}, per verificare l'esistenza di eccezioni come i
dati urgenti su un socket, (vedi sez.~\ref{sec:TCP_urgent_data}).

Dato che in genere non si tengono mai sotto controllo fino a
\macro{FD\_SETSIZE} file contemporaneamente, la funzione richiede di
specificare qual è il valore più alto fra i file descriptor indicati nei tre
insiemi precedenti. Questo viene fatto per efficienza, per evitare di passare
e far controllare al kernel una quantità di memoria superiore a quella
necessaria. Questo limite viene indicato tramite l'argomento \param{ndfs}, che
deve corrispondere al valore massimo aumentato di uno. Si ricordi infatti che
i file descriptor sono numerati progressivamente a partire da zero, ed il
valore indica il numero più alto fra quelli da tenere sotto controllo,
dimenticarsi di aumentare di uno il valore di \param{ndfs} è un errore comune.

Infine l'argomento \param{timeout}, espresso con il puntatore ad una struttura
di tipo \struct{timeval} (vedi fig.~\ref{fig:sys_timeval_struct}) specifica un
tempo massimo di attesa prima che la funzione ritorni; se impostato a
\val{NULL} la funzione attende indefinitamente. Si può specificare anche un
tempo nullo (cioè una struttura \struct{timeval} con i campi impostati a
zero), qualora si voglia semplicemente controllare lo stato corrente dei file
descriptor, e così può essere utilizzata eseguire il \textit{polling} su un
gruppo di file descriptor. Usare questo argomento con tutti i \textit{file
  descriptor set} vuoti è un modo portabile, disponibile anche su sistemi in
cui non sono disponibili le funzioni avanzate di sez.~\ref{sec:sig_timer_adv},
per tenere un processo in stato di \textit{sleep} con precisioni inferiori al
secondo.

In caso di successo la funzione restituisce il numero di file descriptor
pronti, seguendo il comportamento previsto dallo standard
POSIX.1-2001,\footnote{si tenga però presente che esistono alcune versioni di
  Unix che non si comportano in questo modo, restituendo un valore positivo
  generico.}  e ciascun insieme viene sovrascritto per indicare quali sono i
file descriptor pronti per le operazioni ad esso relative, in modo da poterli
controllare con \macro{FD\_ISSET}.  Se invece scade il tempo indicato
da \param{timout} viene restituito un valore nullo e i \textit{file descriptor
  set} non vengono modificati. In caso di errore la funzione restituisce $-1$, i
valori dei tre insiemi e di \param{timeout} sono indefiniti e non si può fare
nessun affidamento sul loro contenuto; nelle versioni più recenti della
funzione invece i \textit{file descriptor set} non vengono modificati anche in
caso di errore.

Si tenga presente infine che su Linux, in caso di programmazione
\textit{multi-thread} se un file descriptor viene chiuso in un altro
\textit{thread} rispetto a quello in cui si sta usando \func{select}, questa
non subisce nessun effetto. In altre varianti di sistemi unix-like invece
\func{select} ritorna indicando che il file descriptor è pronto, con
conseguente possibile errore nel caso lo si usi senza che sia stato
riaperto. Lo standard non prevede niente al riguardo e non si deve dare per
assunto nessuno dei due comportamenti se si vogliono scrivere programmi
portabili.

\itindend{file~descriptor~set}

Una volta ritornata la funzione, si potrà controllare quali sono i file
descriptor pronti, ed operare su di essi. Si tenga presente però che
\func{select} fornisce solo di un suggerimento, esistono infatti condizioni in
cui \func{select} può riportare in maniera spuria che un file descriptor è
pronto, ma l'esecuzione di una operazione di I/O si bloccherebbe: ad esempio
con Linux questo avviene quando su un socket arrivano dei dati che poi vengono
scartati perché corrotti (ma sono possibili pure altri casi); in tal caso pur
risultando il relativo file descriptor pronto in lettura una successiva
esecuzione di una \func{read} si bloccherebbe. Per questo motivo quando si usa
l'\textit{I/O multiplexing} è sempre raccomandato l'uso delle funzioni di
lettura e scrittura in modalità non bloccante.

Su Linux quando la \textit{system call} \func{select} viene interrotta da un
segnale modifica il valore nella struttura puntata da \param{timeout},
impostandolo al tempo restante. In tal caso infatti si ha un errore di
\errcode{EINTR} ed occorre rilanciare la funzione per proseguire l'attesa, ed
in questo modo non è necessario ricalcolare tutte le volte il tempo
rimanente. Questo può causare problemi di portabilità sia quando si usa codice
scritto su Linux che legge questo valore, sia quando si usano programmi
scritti per altri sistemi che non dispongono di questa caratteristica e
ricalcolano \param{timeout} tutte le volte. In genere questa caratteristica è
disponibile nei sistemi che derivano da System V e non è disponibile per
quelli che derivano da BSD; lo standard POSIX.1-2001 non permette questo
comportamento e per questo motivo la \acr{glibc} nasconde il comportamento
passando alla \textit{system call} una copia dell'argomento \param{timeout}.

Uno dei problemi che si presentano con l'uso di \func{select} è che il suo
comportamento dipende dal valore del file descriptor che si vuole tenere sotto
controllo.  Infatti il kernel riceve con \param{ndfs} un limite massimo per
tale valore, e per capire quali sono i file descriptor da tenere sotto
controllo dovrà effettuare una scansione su tutto l'intervallo, che può anche
essere molto ampio anche se i file descriptor sono solo poche unità; tutto ciò
ha ovviamente delle conseguenze ampiamente negative per le prestazioni.

Inoltre c'è anche il problema che il numero massimo dei file che si possono
tenere sotto controllo, la funzione è nata quando il kernel consentiva un
numero massimo di 1024 file descriptor per processo, adesso che il numero può
essere arbitrario si viene a creare una dipendenza del tutto artificiale dalle
dimensioni della struttura \type{fd\_set}, che può necessitare di essere
estesa, con ulteriori perdite di prestazioni. 

Lo standard POSIX è rimasto a lungo senza primitive per l'\textit{I/O
  multiplexing}, introdotto solo con le ultime revisioni dello standard (POSIX
1003.1g-2000 e POSIX 1003.1-2001). La scelta è stata quella di seguire
l'interfaccia creata da BSD, ma prevede che tutte le funzioni ad esso relative
vengano dichiarate nell'header \headfiled{sys/select.h}, che sostituisce i
precedenti, ed inoltre aggiunge a \func{select} una nuova funzione
\funcd{pselect},\footnote{il supporto per lo standard POSIX 1003.1-2001, ed
  l'header \headfile{sys/select.h}, compaiono in Linux a partire dalla
  \acr{glibc} 2.1. Le \acr{libc4} e \acr{libc5} non contengono questo header,
  la \acr{glibc} 2.0 contiene una definizione sbagliata di \func{psignal},
  senza l'argomento \param{sigmask}, la definizione corretta è presente dalle
  \acr{glibc} 2.1-2.2.1 se si è definito \macro{\_GNU\_SOURCE} e nelle
  \acr{glibc} 2.2.2-2.2.4 se si è definito \macro{\_XOPEN\_SOURCE} con valore
  maggiore di 600.} il cui prototipo è:

\begin{funcproto}{
\fhead{sys/select.h}
\fdecl{int pselect(int n, fd\_set *readfds, fd\_set *writefds, 
  fd\_set *exceptfds, \\ 
\phantom{int pselect(}struct timespec *timeout, sigset\_t *sigmask)}
\fdesc{Attende che uno dei file descriptor degli insiemi specificati diventi
  attivo.} 
}
{La funzione ritorna il numero (anche nullo) di file descriptor che sono
  attivi in caso di successo e $-1$ per un errore, nel qual caso \var{errno}
  assumerà uno dei valori:
  \begin{errlist}
  \item[\errcode{EBADF}] si è specificato un file descriptor sbagliato in uno
    degli insiemi.
  \item[\errcode{EINTR}] la funzione è stata interrotta da un segnale.
  \item[\errcode{EINVAL}] si è specificato per \param{ndfs} un valore negativo
    o un valore non valido per \param{timeout}.
   \end{errlist}
   ed inoltre \errval{ENOMEM} nel suo significato generico.
}
\end{funcproto}

La funzione è sostanzialmente identica a \func{select}, solo che usa una
struttura \struct{timespec} (vedi fig.~\ref{fig:sys_timespec_struct}) per
indicare con maggiore precisione il timeout e non ne aggiorna il valore in
caso di interruzione. In realtà anche in questo caso la \textit{system call}
di Linux aggiorna il valore al tempo rimanente, ma la funzione fornita dalla
\acr{glibc} modifica questo comportamento passando alla \textit{system call}
una variabile locale, in modo da mantenere l'aderenza allo standard POSIX che
richiede che il valore di \param{timeout} non sia modificato. 

Rispetto a \func{select} la nuova funzione prende un argomento
aggiuntivo \param{sigmask}, un puntatore ad una maschera di segnali (si veda
sez.~\ref{sec:sig_sigmask}).  Nell'esecuzione la maschera dei segnali corrente
viene sostituita da quella così indicata immediatamente prima di eseguire
l'attesa, e viene poi ripristinata al ritorno della funzione. L'uso
di \param{sigmask} è stato introdotto allo scopo di prevenire possibili
\textit{race condition} quando oltre alla presenza di dati sui file descriptor
come nella \func{select} ordinaria, ci si deve porre in attesa anche
dell'arrivo di un segnale.

Come abbiamo visto in sez.~\ref{sec:sig_example} la tecnica classica per
rilevare l'arrivo di un segnale è quella di utilizzare il gestore per
impostare una variabile globale e controllare questa nel corpo principale del
programma; abbiamo visto in quell'occasione come questo lasci spazio a
possibili \textit{race condition}, per cui diventa essenziale utilizzare
\func{sigprocmask} per disabilitare la ricezione del segnale prima di eseguire
il controllo e riabilitarlo dopo l'esecuzione delle relative operazioni, onde
evitare l'arrivo di un segnale immediatamente dopo il controllo, che andrebbe
perso.

Nel nostro caso il problema si pone quando, oltre al segnale, si devono tenere
sotto controllo anche dei file descriptor con \func{select}, in questo caso si
può fare conto sul fatto che all'arrivo di un segnale essa verrebbe interrotta
e si potrebbero eseguire di conseguenza le operazioni relative al segnale e
alla gestione dati con un ciclo del tipo:
\includecodesnip{listati/select_race.c} 
qui però emerge una \textit{race condition}, perché se il segnale arriva prima
della chiamata a \func{select}, questa non verrà interrotta, e la ricezione
del segnale non sarà rilevata.

Per questo è stata introdotta \func{pselect} che attraverso l'argomento
\param{sigmask} permette di riabilitare la ricezione il segnale
contestualmente all'esecuzione della funzione,\footnote{in Linux però, fino al
  kernel 2.6.16, non era presente la relativa \textit{system call}, e la
  funzione era implementata nella \acr{glibc} attraverso \func{select} (vedi
  \texttt{man select\_tut}) per cui la possibilità di \textit{race condition}
  permaneva; in tale situazione si può ricorrere ad una soluzione alternativa,
  chiamata \itindex{self-pipe~trick} \textit{self-pipe trick}, che consiste
  nell'aprire una \textit{pipe} (vedi sez.~\ref{sec:ipc_pipes}) ed usare
  \func{select} sul capo in lettura della stessa; si può indicare l'arrivo di
  un segnale scrivendo sul capo in scrittura all'interno del gestore dello
  stesso; in questo modo anche se il segnale va perso prima della chiamata di
  \func{select} questa lo riconoscerà comunque dalla presenza di dati sulla
  \textit{pipe}.} ribloccandolo non appena essa ritorna, così che il
precedente codice potrebbe essere riscritto nel seguente modo:
\includecodesnip{listati/pselect_norace.c} 
in questo caso utilizzando \var{oldmask} durante l'esecuzione di
\func{pselect} la ricezione del segnale sarà abilitata, ed in caso di
interruzione si potranno eseguire le relative operazioni.


\subsection{Le funzioni \func{poll} e \func{ppoll}}
\label{sec:file_poll}

Nello sviluppo di System V, invece di utilizzare l'interfaccia di
\func{select}, che è una estensione tipica di BSD, è stata introdotta una
interfaccia completamente diversa, basata sulla funzione di sistema
\funcd{poll},\footnote{la funzione è prevista dallo standard XPG4, ed è stata
  introdotta in Linux come \textit{system call} a partire dal kernel 2.1.23 ed
  inserita nelle \acr{libc} 5.4.28, originariamente l'argomento \param{nfds}
  era di tipo \ctyp{unsigned int}, la funzione è stata inserita nello standard
  POSIX.1-2001 in cui è stato introdotto il tipo nativo \typed{nfds\_t}.} il
cui prototipo è:

\begin{funcproto}{
\fhead{sys/poll.h}
\fdecl{int poll(struct pollfd *ufds, nfds\_t nfds, int timeout)}
\fdesc{Attende un cambiamento di stato su un insieme di file
  descriptor.} 
}

{La funzione ritorna $0$ in caso di successo e $-1$ per un errore, nel qual
  caso \var{errno} assumerà uno dei valori: 
  \begin{errlist}
  \item[\errcode{EBADF}] si è specificato un file descriptor sbagliato in uno
    degli insiemi.
  \item[\errcode{EINTR}] la funzione è stata interrotta da un segnale.
  \item[\errcode{EINVAL}] il valore di \param{nfds} eccede il limite
    \const{RLIMIT\_NOFILE}.
  \end{errlist}
  ed inoltre \errval{EFAULT} e \errval{ENOMEM} nel loro significato generico.}
\end{funcproto}

La funzione permette di tenere sotto controllo contemporaneamente \param{ndfs}
file descriptor, specificati attraverso il puntatore \param{ufds} ad un
vettore di strutture \struct{pollfd}.  Come con \func{select} si può
interrompere l'attesa dopo un certo tempo, questo deve essere specificato con
l'argomento \param{timeout} in numero di millisecondi: un valore negativo
indica un'attesa indefinita, mentre un valore nullo comporta il ritorno
immediato, e può essere utilizzato per impiegare \func{poll} in modalità
\textsl{non-bloccante}.

\begin{figure}[!htb]
  \footnotesize \centering
  \begin{minipage}[c]{0.90\textwidth}
    \includestruct{listati/pollfd.h}
  \end{minipage} 
  \normalsize 
  \caption{La struttura \structd{pollfd}, utilizzata per specificare le
    modalità di controllo di un file descriptor alla funzione \func{poll}.}
  \label{fig:file_pollfd}
\end{figure}

Per ciascun file da controllare deve essere inizializzata una struttura
\struct{pollfd} nel vettore indicato dall'argomento \param{ufds}.  La
struttura, la cui definizione è riportata in fig.~\ref{fig:file_pollfd},
prevede tre campi: in \var{fd} deve essere indicato il numero del file
descriptor da controllare, in \var{events} deve essere specificata una
maschera binaria di flag che indichino il tipo di evento che si vuole
controllare, mentre in \var{revents} il kernel restituirà il relativo
risultato. 

Usando un valore negativo per \param{fd} la corrispondente struttura sarà
ignorata da \func{poll} ed il campo \var{revents} verrà azzerato, questo
consente di eliminare temporaneamente un file descriptor dalla lista senza
dover modificare il vettore \param{ufds}. Dato che i dati in ingresso sono del
tutto indipendenti da quelli in uscita (che vengono restituiti in
\var{revents}) non è necessario reinizializzare tutte le volte il valore delle
strutture \struct{pollfd} a meno di non voler cambiare qualche condizione.

Le costanti che definiscono i valori relativi ai bit usati nelle maschere
binarie dei campi \var{events} e \var{revents} sono riportate in
tab.~\ref{tab:file_pollfd_flags}, insieme al loro significato. Le si sono
suddivise in tre gruppi principali, nel primo gruppo si sono indicati i bit
utilizzati per controllare l'attività in ingresso, nel secondo quelli per
l'attività in uscita, infine il terzo gruppo contiene dei valori che vengono
utilizzati solo nel campo \var{revents} per notificare delle condizioni di
errore.

\begin{table}[htb]
  \centering
  \footnotesize
  \begin{tabular}[c]{|l|l|}
    \hline
    \textbf{Flag}  & \textbf{Significato} \\
    \hline
    \hline
    \constd{POLLIN}    & È possibile la lettura.\\
    \constd{POLLRDNORM}& Sono disponibili in lettura dati normali.\\ 
    \constd{POLLRDBAND}& Sono disponibili in lettura dati prioritari.\\
    \constd{POLLPRI}   & È possibile la lettura di dati urgenti.\\ 
    \hline
    \constd{POLLOUT}   & È possibile la scrittura immediata.\\
    \constd{POLLWRNORM}& È possibile la scrittura di dati normali.\\ 
    \constd{POLLWRBAND}& È possibile la scrittura di dati prioritari.\\
    \hline
    \constd{POLLERR}   & C'è una condizione di errore.\\
    \constd{POLLHUP}   & Si è verificato un hung-up.\\
    \constd{POLLRDHUP} & Si è avuta una \textsl{half-close} su un
                        socket.\footnotemark\\ 
    \constd{POLLNVAL}  & Il file descriptor non è aperto.\\
    \hline
    \constd{POLLMSG}   & Definito per compatibilità con SysV.\\
    \hline    
  \end{tabular}
  \caption{Costanti per l'identificazione dei vari bit dei campi
    \var{events} e \var{revents} di \struct{pollfd}.}
  \label{tab:file_pollfd_flags}
\end{table}

\footnotetext{si tratta di una estensione specifica di Linux, disponibile a
  partire dal kernel 2.6.17 definendo la marco \macro{\_GNU\_SOURCE}, che
  consente di riconoscere la chiusura in scrittura dell'altro capo di un
  socket, situazione che si viene chiamata appunto \textit{half-close}
  (\textsl{mezza chiusura}) su cui torneremo con maggiori dettagli in
  sez.~\ref{sec:TCP_shutdown}.}

Il valore \const{POLLMSG} non viene utilizzato ed è definito solo per
compatibilità con l'implementazione di System V che usa i cosiddetti
``\textit{stream}''. Si tratta di una interfaccia specifica di SysV non
presente in Linux, che non ha nulla a che fare con gli \textit{stream} delle
librerie standard del C visti in sez.~\ref{sec:file_stream}. Da essa derivano
i nomi di alcune costanti poiché per quegli \textit{stream} sono definite tre
classi di dati: \textsl{normali}, \textit{prioritari} ed \textit{urgenti}.  In
Linux la distinzione ha senso solo per i dati urgenti dei socket (vedi
sez.~\ref{sec:TCP_urgent_data}), ma su questo e su come \func{poll} reagisce
alle varie condizioni dei socket torneremo in sez.~\ref{sec:TCP_serv_poll},
dove vedremo anche un esempio del suo utilizzo.

Le costanti relative ai diversi tipi di dati normali e prioritari che fanno
riferimento alle implementazioni in stile System V sono \const{POLLRDNORM},
\const{POLLWRNORM}, \const{POLLRDBAND} e \const{POLLWRBAND}. Le prime due sono
equivalenti rispettivamente a \const{POLLIN} e \const{POLLOUT},
\const{POLLRDBAND} non viene praticamente mai usata su Linux mentre
\const{POLLWRBAND} ha senso solo sui socket. In ogni caso queste costanti sono
utilizzabili soltanto qualora si sia definita la macro
\macro{\_XOPEN\_SOURCE}.

In caso di successo \func{poll} ritorna restituendo il numero di file (un
valore positivo) per i quali si è verificata una delle condizioni di attesa
richieste o per i quali si è verificato un errore, avvalorando i relativi bit
di \var{revents}. In caso di errori sui file vengono utilizzati i valori della
terza sezione di tab.~\ref{tab:file_pollfd_flags} che hanno significato solo
per \var{revents} (se specificati in \var{events} vengono ignorati). Un valore
di ritorno nullo indica che si è raggiunto il timeout, mentre un valore
negativo indica un errore nella chiamata, il cui codice viene riportato al
solito tramite \var{errno}.

L'uso di \func{poll} consente di superare alcuni dei problemi illustrati in
precedenza per \func{select}; anzitutto, dato che in questo caso si usa un
vettore di strutture \struct{pollfd} di dimensione arbitraria, non esiste il
limite introdotto dalle dimensioni massime di un \textit{file descriptor set}
e la dimensione dei dati passati al kernel dipende solo dal numero dei file
descriptor che si vogliono controllare, non dal loro valore. Infatti, anche se
usando dei bit un \textit{file descriptor set} può essere più efficiente di un
vettore di strutture \struct{pollfd}, qualora si debba osservare un solo file
descriptor con un valore molto alto ci si troverà ad utilizzare inutilmente un
maggiore quantitativo di memoria.

Inoltre con \func{select} lo stesso \textit{file descriptor set} è usato sia
in ingresso che in uscita, e questo significa che tutte le volte che si vuole
ripetere l'operazione occorre reinizializzarlo da capo. Questa operazione, che
può essere molto onerosa se i file descriptor da tenere sotto osservazione
sono molti, non è invece necessaria con \func{poll}.

Abbiamo visto in sez.~\ref{sec:file_select} come lo standard POSIX preveda una
variante di \func{select} che consente di gestire correttamente la ricezione
dei segnali nell'attesa su un file descriptor.  Con l'introduzione di una
implementazione reale di \func{pselect} nel kernel 2.6.16, è stata aggiunta
anche una analoga funzione che svolga lo stesso ruolo per \func{poll}.

In questo caso si tratta di una estensione che è specifica di Linux e non è
prevista da nessuno standard; essa può essere utilizzata esclusivamente se si
definisce la macro \macro{\_GNU\_SOURCE} ed ovviamente non deve essere usata
se si ha a cuore la portabilità. La funzione è \funcd{ppoll}, ed il suo
prototipo è:

\begin{funcproto}{
\fhead{sys/poll.h}
\fdecl{int ppoll(struct pollfd *fds, nfds\_t nfds, 
  const struct timespec *timeout, \\
\phantom{int ppoll(}const sigset\_t *sigmask)} 

\fdesc{Attende un cambiamento di stato su un insieme di file descriptor.}
}

{La funzione ritorna il numero di file descriptor con attività in caso di
  successo, $0$ se c'è stato un timeout e $-1$ per un errore, nel qual caso
  \var{errno} assumerà uno dei valori:
  \begin{errlist}
  \item[\errcode{EBADF}] si è specificato un file descriptor sbagliato in uno
    degli insiemi.
  \item[\errcode{EINTR}] la funzione è stata interrotta da un segnale.
  \item[\errcode{EINVAL}] il valore di \param{nfds} eccede il limite
    \const{RLIMIT\_NOFILE}.
  \end{errlist}
ed inoltre \errval{EFAULT} e \errval{ENOMEM} nel loro significato generico.
}  
\end{funcproto}

La funzione ha lo stesso comportamento di \func{poll}, solo che si può
specificare, con l'argomento \param{sigmask}, il puntatore ad una maschera di
segnali; questa sarà la maschera utilizzata per tutto il tempo che la funzione
resterà in attesa, all'uscita viene ripristinata la maschera originale.  L'uso
di questa funzione è cioè equivalente, come illustrato nella pagina di
manuale, all'esecuzione atomica del seguente codice:
\includecodesnip{listati/ppoll_means.c} 

Eccetto per \param{timeout}, che come per \func{pselect} deve essere un
puntatore ad una struttura \struct{timespec}, gli altri argomenti comuni con
\func{poll} hanno lo stesso significato, e la funzione restituisce gli stessi
risultati illustrati in precedenza. Come nel caso di \func{pselect} la
\textit{system call} che implementa \func{ppoll} restituisce, se la funzione
viene interrotta da un segnale, il tempo mancante in \param{timeout}, e come
per \func{pselect} la funzione di libreria fornita dalla \acr{glibc} maschera
questo comportamento non modificando mai il valore di \param{timeout} anche se
in questo caso non esiste nessuno standard che richieda questo comportamento.

Infine anche per \func{poll} e \func{ppoll} valgono le considerazioni relative
alla possibilità di avere delle notificazione spurie della disponibilità di
accesso ai file descriptor illustrate per \func{select} in
sez.~\ref{sec:file_select}, che non staremo a ripetere qui.

\subsection{L'interfaccia di \textit{epoll}}
\label{sec:file_epoll}

\itindbeg{epoll}

Nonostante \func{poll} presenti alcuni vantaggi rispetto a \func{select},
anche questa funzione non è molto efficiente quando deve essere utilizzata con
un gran numero di file descriptor,\footnote{in casi del genere \func{select}
  viene scartata a priori, perché può avvenire che il numero di file
  descriptor ecceda le dimensioni massime di un \textit{file descriptor set}.}
in particolare nel caso in cui solo pochi di questi diventano attivi. Il
problema in questo caso è che il tempo impiegato da \func{poll} a trasferire i
dati da e verso il kernel è proporzionale al numero di file descriptor
osservati, non a quelli che presentano attività.

Quando ci sono decine di migliaia di file descriptor osservati e migliaia di
eventi al secondo (il caso classico è quello di un server web di un sito con
molti accessi) l'uso di \func{poll} comporta la necessità di trasferire avanti
ed indietro da \textit{user space} a \textit{kernel space} una lunga lista di
strutture \struct{pollfd} migliaia di volte al secondo. A questo poi si
aggiunge il fatto che la maggior parte del tempo di esecuzione sarà impegnato
ad eseguire una scansione su tutti i file descriptor tenuti sotto controllo
per determinare quali di essi (in genere una piccola percentuale) sono
diventati attivi. In una situazione come questa l'uso delle funzioni classiche
dell'interfaccia dell'\textit{I/O multiplexing} viene a costituire un collo di
bottiglia che degrada irrimediabilmente le prestazioni.

Per risolvere questo tipo di situazioni sono state ideate delle interfacce
specialistiche (come \texttt{/dev/poll} in Solaris, o \texttt{kqueue} in BSD)
il cui scopo fondamentale è quello di restituire solamente le informazioni
relative ai file descriptor osservati che presentano una attività, evitando
così le problematiche appena illustrate. In genere queste prevedono che si
registrino una sola volta i file descriptor da tenere sotto osservazione, e
forniscono un meccanismo che notifica quali di questi presentano attività.

Le modalità con cui avviene la notifica sono due, la prima è quella classica
(quella usata da \func{poll} e \func{select}) che viene chiamata \textit{level
  triggered}.\footnote{la nomenclatura è stata introdotta da Jonathan Lemon in
  un articolo su \texttt{kqueue} al BSDCON 2000, e deriva da quella usata
  nell'elettronica digitale.} In questa modalità vengono notificati i file
descriptor che sono \textsl{pronti} per l'operazione richiesta, e questo
avviene indipendentemente dalle operazioni che possono essere state fatte su
di essi a partire dalla precedente notifica.  Per chiarire meglio il concetto
ricorriamo ad un esempio: se su un file descriptor sono diventati disponibili
in lettura 2000 byte ma dopo la notifica ne sono letti solo 1000 (ed è quindi
possibile eseguire una ulteriore lettura dei restanti 1000), in modalità
\textit{level triggered} questo sarà nuovamente notificato come
\textsl{pronto}.

La seconda modalità, è detta \textit{edge triggered}, e prevede che invece
vengano notificati solo i file descriptor che hanno subito una transizione da
\textsl{non pronti} a \textsl{pronti}. Questo significa che in modalità
\textit{edge triggered} nel caso del precedente esempio il file descriptor
diventato pronto da cui si sono letti solo 1000 byte non verrà nuovamente
notificato come pronto, nonostante siano ancora disponibili in lettura 1000
byte. Solo una volta che si saranno esauriti tutti i dati disponibili, e che
il file descriptor sia tornato non essere pronto, si potrà ricevere una
ulteriore notifica qualora ritornasse pronto.

Nel caso di Linux al momento la sola interfaccia che fornisce questo tipo di
servizio è chiamata \textit{epoll},\footnote{l'interfaccia è stata creata da
  Davide Libenzi, ed è stata introdotta per la prima volta nel kernel 2.5.44,
  ma la sua forma definitiva è stata raggiunta nel kernel 2.5.66, il supporto
  è stato aggiunto nella \acr{glibc} a partire dalla versione 2.3.2.} anche se
sono state in discussione altre interfacce con le quali effettuare lo stesso
tipo di operazioni; \textit{epoll} è in grado di operare sia in modalità
\textit{level triggered} che \textit{edge triggered}.

La prima versione di \textit{epoll} prevedeva l'uso di uno speciale file di
dispositivo, \texttt{/dev/epoll}, per ottenere un file descriptor da
utilizzare con le funzioni dell'interfaccia ma poi si è passati all'uso di
apposite \textit{system call}.  Il primo passo per usare l'interfaccia di
\textit{epoll} è pertanto quello ottenere detto file descriptor chiamando una
delle due funzioni di sistema \funcd{epoll\_create} e \funcd{epoll\_create1},
i cui prototipi sono:

\begin{funcproto}{
\fhead{sys/epoll.h}
\fdecl{int epoll\_create(int size)}
\fdecl{int epoll\_create1(int flags)}

\fdesc{Apre un file descriptor per \textit{epoll}.}
}
{Le funzioni ritornano un file descriptor per \textit{epoll} in caso di
  successo e $-1$ per un errore, nel qual caso \var{errno} assumerà uno dei
  valori:
  \begin{errlist}
  \item[\errcode{EINVAL}] si è specificato un valore di \param{size} non
    positivo o non valido per \param{flags}.
  \item[\errcode{EMFILE}] si è raggiunto il limite sul numero massimo di
    istanze di \textit{epoll} per utente stabilito da
    \sysctlfiled{fs/epoll/max\_user\_instances}.
  \item[\errcode{ENFILE}] si è raggiunto il massimo di file descriptor aperti
    nel sistema.
  \item[\errcode{ENOMEM}] non c'è sufficiente memoria nel kernel per creare
    l'istanza.
  \end{errlist}
}  
\end{funcproto}

Entrambe le funzioni restituiscono un file descriptor, detto anche
\textit{epoll descriptor}; si tratta di un file descriptor speciale (per cui
\func{read} e \func{write} non sono supportate) che viene associato alla
infrastruttura utilizzata dal kernel per gestire la notifica degli eventi, e
che può a sua volta essere messo sotto osservazione con una chiamata a
\func{select}, \func{poll} o \func{epoll\_ctl}; in tal caso risulterà pronto
quando saranno disponibili eventi da notificare riguardo i file descriptor da
lui osservati.\footnote{è anche possibile inviarlo ad un altro processo
  attraverso un socket locale (vedi sez.~\ref{sec:sock_fd_passing}) ma
  l'operazione non ha alcun senso dato che il nuovo processo non avrà a
  disposizione le copie dei file descriptor messe sotto osservazione tramite
  esso.} Una volta che se ne sia terminato l'uso si potranno rilasciare tutte
le risorse allocate chiudendolo semplicemente con \func{close}.

Nel caso di \func{epoll\_create} l'argomento \param{size} serviva a dare
l'indicazione del numero di file descriptor che si vorranno tenere sotto
controllo, e costituiva solo un suggerimento per semplificare l'allocazione di
risorse sufficienti, non un valore massimo, ma a partire dal kernel 2.6.8 esso
viene totalmente ignorato e l'allocazione è sempre dinamica.

La seconda versione della funzione, \func{epoll\_create1} è stata introdotta
come estensione della precedente (è disponibile solo a partire dal kernel
2.6.27) per poter passare dei flag di controllo come maschera binaria in fase
di creazione del file descriptor. Al momento l'unico valore legale
per \param{flags} (a parte lo zero) è \constd{EPOLL\_CLOEXEC}, che consente di
impostare in maniera atomica sul file descriptor il flag di
\textit{close-on-exec} (si è trattato il significato di \const{O\_CLOEXEC} in
sez.~\ref{sec:file_open_close}), senza che sia necessaria una successiva
chiamata a \func{fcntl}.

Una volta ottenuto un file descriptor per \textit{epoll} il passo successivo è
indicare quali file descriptor mettere sotto osservazione e quali operazioni
controllare, per questo si deve usare la seconda funzione di sistema
dell'interfaccia, \funcd{epoll\_ctl}, il cui prototipo è:

\begin{funcproto}{
\fhead{sys/epoll.h}
\fdecl{int epoll\_ctl(int epfd, int op, int fd, struct epoll\_event *event)}

\fdesc{Esegue le operazioni di controllo di \textit{epoll}.}
}

{La funzione ritorna $0$ in caso di successo e $-1$ per un errore, nel qual
  caso \var{errno} assumerà uno dei valori:
  \begin{errlist}
  \item[\errcode{EBADF}] i file descriptor \param{epfd} o \param{fd} non sono
    validi.
  \item[\errcode{EEXIST}] l'operazione richiesta è \const{EPOLL\_CTL\_ADD} ma
    \param{fd} è già stato inserito in \param{epfd}.
  \item[\errcode{EINVAL}] il file descriptor \param{epfd} non è stato ottenuto
    con \func{epoll\_create}, o \param{fd} è lo stesso \param{epfd} o
    l'operazione richiesta con \param{op} non è supportata.
  \item[\errcode{ENOENT}] l'operazione richiesta è \const{EPOLL\_CTL\_MOD} o
    \const{EPOLL\_CTL\_DEL} ma \param{fd} non è inserito in \param{epfd}.
  \item[\errcode{ENOMEM}] non c'è sufficiente memoria nel kernel gestire
    l'operazione richiesta.
  \item[\errcode{ENOSPC}] si è raggiunto il limite massimo di registrazioni
    per utente di file descriptor da osservare imposto da
    \sysctlfiled{fs/epoll/max\_user\_watches}.
  \item[\errcode{EPERM}] il file associato a \param{fd} non supporta l'uso di
    \textit{epoll}.
  \end{errlist}
  }  
\end{funcproto}

La funzione prende sempre come primo argomento un file descriptor di
\textit{epoll}, \param{epfd}, che indica quale istanza di \textit{epoll} usare
e deve pertanto essere stato ottenuto in precedenza con una chiamata a
\func{epoll\_create} o \func{epoll\_create1}. L'argomento \param{fd} indica
invece il file descriptor che si vuole tenere sotto controllo, quest'ultimo
può essere un qualunque file descriptor utilizzabile con \func{poll}, ed anche
un altro file descriptor di \textit{epoll}, ma non lo stesso \param{epfd}.

Il comportamento della funzione viene controllato dal valore dall'argomento
\param{op} che consente di specificare quale operazione deve essere eseguita.
Le costanti che definiscono i valori utilizzabili per \param{op}
sono riportate in tab.~\ref{tab:epoll_ctl_operation}, assieme al significato
delle operazioni cui fanno riferimento.

\begin{table}[htb]
  \centering
  \footnotesize
  \begin{tabular}[c]{|l|p{8cm}|}
    \hline
    \textbf{Valore}  & \textbf{Significato} \\
    \hline
    \hline
    \constd{EPOLL\_CTL\_ADD}& Aggiunge un nuovo file descriptor da osservare
                              \param{fd} alla lista dei file descriptor
                              controllati tramite \param{epfd}, in
                              \param{event} devono essere specificate le
                              modalità di osservazione.\\
    \constd{EPOLL\_CTL\_MOD}& Modifica le modalità di osservazione del file
                              descriptor \param{fd} secondo il contenuto di
                              \param{event}.\\
    \constd{EPOLL\_CTL\_DEL}& Rimuove il file descriptor \param{fd} dalla lista
                              dei file controllati tramite \param{epfd}.\\
   \hline    
  \end{tabular}
  \caption{Valori dell'argomento \param{op} che consentono di scegliere quale
    operazione di controllo effettuare con la funzione \func{epoll\_ctl}.} 
  \label{tab:epoll_ctl_operation}
\end{table}

% era stata aggiunta EPOLL_CTL_DISABLE in previsione del kernel 3.7, vedi
% http://lwn.net/Articles/520012/ e http://lwn.net/Articles/520198/
% ma non è mai stata inserita.

Le modalità di utilizzo di \textit{epoll} prevedono che si definisca qual'è
l'insieme dei file descriptor da tenere sotto controllo utilizzando una serie
di chiamate a \const{EPOLL\_CTL\_ADD}.\footnote{un difetto dell'interfaccia è
  che queste chiamate devono essere ripetute per ciascun file descriptor,
  incorrendo in una perdita di prestazioni qualora il numero di file
  descriptor sia molto grande; per questo è stato proposto di introdurre come
  estensione una funzione \code{epoll\_ctlv} che consenta di effettuare con
  una sola chiamata le impostazioni per un blocco di file descriptor.} L'uso
di \const{EPOLL\_CTL\_MOD} consente in seguito di modificare le modalità di
osservazione di un file descriptor che sia già stato aggiunto alla lista di
osservazione. Qualora non si abbia più interesse nell'osservazione di un file
descriptor lo si può rimuovere dalla lista associata a \param{epfd} con
\const{EPOLL\_CTL\_DEL}.

Anche se è possibile tenere sotto controllo lo stesso file descriptor in due
istanze distinte di \textit{epoll} in genere questo è sconsigliato in quanto
entrambe riceveranno le notifiche, e gestire correttamente le notifiche
multiple richiede molta attenzione. Se invece si cerca di inserire due volte
lo stesso file descriptor nella stessa istanza di \textit{epoll} la funzione
fallirà con un errore di \errval{EEXIST}.  Tuttavia è possibile inserire nella
stessa istanza file descriptor duplicati (si ricordi quanto visto in
sez.~\ref{sec:file_dup}), una tecnica che può essere usata per registrarli con
un valore diverso per \param{events} e classificare così diversi tipi di
eventi.

Si tenga presente che quando si chiude un file descriptor questo, se era stato
posto sotto osservazione da una istanza di \textit{epoll}, viene rimosso
automaticamente solo nel caso esso sia l'unico riferimento al file aperto
sottostante (più precisamente alla struttura \kstruct{file}, si ricordi
fig.~\ref{fig:file_dup}) e non è necessario usare
\const{EPOLL\_CTL\_DEL}. Questo non avviene qualora esso sia stato duplicato
(perché la suddetta struttura non viene disallocata) e si potranno ricevere
eventi ad esso relativi anche dopo che lo si è chiuso; per evitare
l'inconveniente è necessario rimuoverlo esplicitamente con
\const{EPOLL\_CTL\_DEL}.

L'ultimo argomento, \param{event}, deve essere un puntatore ad una struttura
di tipo \struct{epoll\_event}, ed ha significato solo con le operazioni
\const{EPOLL\_CTL\_MOD} e \const{EPOLL\_CTL\_ADD}, per le quali serve ad
indicare quale tipo di evento relativo ad \param{fd} si vuole che sia tenuto
sotto controllo.  L'argomento viene ignorato con l'operazione
\const{EPOLL\_CTL\_DEL}.\footnote{fino al kernel 2.6.9 era comunque richiesto
  che questo fosse un puntatore valido, anche se poi veniva ignorato; a
  partire dal 2.6.9 si può specificare anche un valore \val{NULL} ma se si
  vuole mantenere la compatibilità con le versioni precedenti occorre usare un
  puntatore valido.}

\begin{figure}[!htb]
  \footnotesize \centering
  \begin{minipage}[c]{0.90\textwidth}
    \includestruct{listati/epoll_event.h}
  \end{minipage} 
  \normalsize 
  \caption{La struttura \structd{epoll\_event}, che consente di specificare
    gli eventi associati ad un file descriptor controllato con
    \textit{epoll}.}
  \label{fig:epoll_event}
\end{figure}

La struttura \struct{epoll\_event} è l'analoga di \struct{pollfd} e come
quest'ultima serve sia in ingresso (quando usata con \func{epoll\_ctl}) ad
impostare quali eventi osservare, che in uscita (nei risultati ottenuti con
\func{epoll\_wait}) per ricevere le notifiche degli eventi avvenuti.  La sua
definizione è riportata in fig.~\ref{fig:epoll_event}. 

Il primo campo, \var{events}, è una maschera binaria in cui ciascun bit
corrisponde o ad un tipo di evento, o una modalità di notifica; detto campo
deve essere specificato come OR aritmetico delle costanti riportate in
tab.~\ref{tab:epoll_events}. Nella prima parte della tabella si sono indicate
le costanti che permettono di indicare il tipo di evento, che sono le
equivalenti delle analoghe di tab.~\ref{tab:file_pollfd_flags} per
\func{poll}. Queste sono anche quelle riportate nella struttura
\struct{epoll\_event} restituita da \func{epoll\_wait} per indicare il tipo di
evento presentatosi, insieme a quelle della seconda parte della tabella, che
vengono comunque riportate anche se non le si sono impostate con
\func{epoll\_ctl}. La terza parte della tabella contiene le costanti che
modificano le modalità di notifica.

\begin{table}[htb]
  \centering
  \footnotesize
  \begin{tabular}[c]{|l|p{10cm}|}
    \hline
    \textbf{Valore}  & \textbf{Significato} \\
    \hline
    \hline
    \constd{EPOLLIN}     & Il file è pronto per le operazioni di lettura
                          (analogo di \const{POLLIN}).\\
    \constd{EPOLLOUT}    & Il file è pronto per le operazioni di scrittura
                          (analogo di \const{POLLOUT}).\\
    \constd{EPOLLRDHUP}  & L'altro capo di un socket di tipo
                          \const{SOCK\_STREAM} (vedi sez.~\ref{sec:sock_type})
                          ha chiuso la connessione o il capo in scrittura
                          della stessa (vedi
                          sez.~\ref{sec:TCP_shutdown}).\footnotemark\\
    \constd{EPOLLPRI}    & Ci sono dati urgenti disponibili in lettura (analogo
                          di \const{POLLPRI}); questa condizione viene comunque
                          riportata in uscita, e non è necessaria impostarla
                          in ingresso.\\ 
    \hline
    \constd{EPOLLERR}    & Si è verificata una condizione di errore 
                          (analogo di \const{POLLERR}); questa condizione
                          viene comunque riportata in uscita, e non è
                          necessaria impostarla in ingresso.\\
    \constd{EPOLLHUP}    & Si è verificata una condizione di hung-up; questa
                          condizione viene comunque riportata in uscita, e non
                          è necessaria impostarla in ingresso.\\
    \hline
    \constd{EPOLLET}     & Imposta la notifica in modalità \textit{edge
                            triggered} per il file descriptor associato.\\ 
    \constd{EPOLLONESHOT}& Imposta la modalità \textit{one-shot} per il file
                          descriptor associato (questa modalità è disponibile
                          solo a partire dal kernel 2.6.2).\\
    \constd{EPOLLWAKEUP} & Attiva la prevenzione della sospensione del sistema
                          se il file descriptor che si è marcato con esso
                          diventa pronto (aggiunto a partire dal kernel 3.5),
                          può essere impostato solo dall'amministratore (o da
                          un processo con la capacità
                          \const{CAP\_BLOCK\_SUSPEND}).\\ 
    \hline
  \end{tabular}
  \caption{Costanti che identificano i bit del campo \param{events} di
    \struct{epoll\_event}.}
  \label{tab:epoll_events}
\end{table}

\footnotetext{questa modalità è disponibile solo a partire dal kernel 2.6.17,
  ed è utile per riconoscere la chiusura di una connessione dall'altro capo di
  un socket quando si lavora in modalità \textit{edge triggered}.}

% TODO aggiunto con il kernel 4.5  EPOLLEXCLUSIVE, vedi
% http://lwn.net/Articles/633422/#excl 

Il secondo campo, \var{data}, è una \dirct{union} che serve a identificare il
file descriptor a cui si intende fare riferimento, ed in astratto può
contenere un valore qualsiasi (specificabile in diverse forme) che ne permetta
una indicazione univoca. Il modo più comune di usarlo però è quello in cui si
specifica il terzo argomento di \func{epoll\_ctl} nella forma
\var{event.data.fd}, assegnando come valore di questo campo lo stesso valore
dell'argomento \param{fd}, cosa che permette una immediata identificazione del
file descriptor.

% TODO verificare se prima o poi epoll_ctlv verrà introdotta

Le impostazioni di default prevedono che la notifica degli eventi richiesti
sia effettuata in modalità \textit{level triggered}, a meno che sul file
descriptor non si sia impostata la modalità \textit{edge triggered},
registrandolo con \const{EPOLLET} attivo nel campo \var{events}.  

Infine una particolare modalità di notifica è quella impostata con
\const{EPOLLONESHOT}: a causa dell'implementazione di \textit{epoll} infatti
quando si è in modalità \textit{edge triggered} l'arrivo in rapida successione
di dati in blocchi separati (questo è tipico con i socket di rete, in quanto i
dati arrivano a pacchetti) può causare una generazione di eventi (ad esempio
segnalazioni di dati in lettura disponibili) anche se la condizione è già
stata rilevata (si avrebbe cioè una rottura della logica \textit{edge
  triggered}).

Anche se la situazione è facile da gestire, la si può evitare utilizzando
\const{EPOLLONESHOT} per impostare la modalità \textit{one-shot}, in cui la
notifica di un evento viene effettuata una sola volta, dopo di che il file
descriptor osservato, pur restando nella lista di osservazione, viene
automaticamente disattivato (la cosa avviene contestualmente al ritorno di
\func{epoll\_wait} a causa dell'evento in questione) e per essere riutilizzato
dovrà essere riabilitato esplicitamente con una successiva chiamata con
\const{EPOLL\_CTL\_MOD}.

Una volta impostato l'insieme di file descriptor che si vogliono osservare con
i relativi eventi, la funzione di sistema che consente di attendere
l'occorrenza di uno di tali eventi è \funcd{epoll\_wait}, il cui prototipo è:

\begin{funcproto}{
\fhead{sys/epoll.h}
\fdecl{int epoll\_wait(int epfd, struct epoll\_event * events, int maxevents,
  int timeout)}

\fdesc{Attende che uno dei file descriptor osservati sia pronto.}
}

{La funzione ritorna il numero di file descriptor pronti in caso di successo e
  $-1$ per un errore, nel qual caso \var{errno} assumerà uno dei valori:
  \begin{errlist}
  \item[\errcode{EBADF}] il file descriptor \param{epfd} non è valido.
  \item[\errcode{EFAULT}] il puntatore \param{events} non è valido.
  \item[\errcode{EINTR}] la funzione è stata interrotta da un segnale prima
    della scadenza di \param{timeout}.
  \item[\errcode{EINVAL}] il file descriptor \param{epfd} non è stato ottenuto
    con \func{epoll\_create}, o \param{maxevents} non è maggiore di zero.
  \end{errlist}
}  
\end{funcproto}

La funzione si blocca in attesa di un evento per i file descriptor registrati
nella lista di osservazione di \param{epfd} fino ad un tempo massimo
specificato in millisecondi tramite l'argomento \param{timeout}. Gli eventi
registrati vengono riportati in un vettore di strutture \struct{epoll\_event}
(che deve essere stato allocato in precedenza) all'indirizzo indicato
dall'argomento \param{events}, fino ad un numero massimo di eventi impostato
con l'argomento \param{maxevents}.

La funzione ritorna il numero di eventi rilevati, o un valore nullo qualora
sia scaduto il tempo massimo impostato con \param{timeout}. Per quest'ultimo,
oltre ad un numero di millisecondi, si può utilizzare il valore nullo, che
indica di non attendere e ritornare immediatamente (anche in questo caso il
valore di ritorno sarà nullo) o il valore $-1$, che indica un'attesa
indefinita. L'argomento \param{maxevents} dovrà invece essere sempre un intero
positivo.

Come accennato la funzione restituisce i suoi risultati nel vettore di
strutture \struct{epoll\_event} puntato da \param{events}; in tal caso nel
campo \param{events} di ciascuna di esse saranno attivi i flag relativi agli
eventi accaduti, mentre nel campo \var{data} sarà restituito il valore che era
stato impostato per il file descriptor per cui si è verificato l'evento quando
questo era stato registrato con le operazioni \const{EPOLL\_CTL\_MOD} o
\const{EPOLL\_CTL\_ADD}, in questo modo il campo \var{data} consente di
identificare il file descriptor, ed è per questo che, come accennato, è
consuetudine usare per \var{data} il valore del file descriptor stesso.

Si ricordi che le occasioni per cui \func{epoll\_wait} ritorna dipendono da
come si è impostata la modalità di osservazione (se \textit{level triggered} o
\textit{edge triggered}) del singolo file descriptor. L'interfaccia assicura
che se arrivano più eventi fra due chiamate successive ad \func{epoll\_wait}
questi vengano combinati. Inoltre qualora su un file descriptor fossero
presenti eventi non ancora notificati, e si effettuasse una modifica
dell'osservazione con \const{EPOLL\_CTL\_MOD}, questi verrebbero riletti alla
luce delle modifiche.

Si tenga presente infine che con l'uso della modalità \textit{edge triggered}
il ritorno di \func{epoll\_wait} avviene solo quando il file descriptor ha
cambiato stato diventando pronto. Esso non sarà riportato nuovamente fino ad
un altro cambiamento di stato, per cui occorre assicurarsi di aver
completamente esaurito le operazioni su di esso.  Questa condizione viene
generalmente rilevata dall'occorrere di un errore di \errcode{EAGAIN} al
ritorno di una \func{read} o una \func{write}, (è opportuno ricordare ancora
una volta che l'uso dell'\textit{I/O multiplexing} richiede di operare sui
file in modalità non bloccante) ma questa non è la sola modalità possibile, ad
esempio la condizione può essere riconosciuta anche per il fatto che sono
stati restituiti meno dati di quelli richiesti.

Si tenga presente che in modalità \textit{edge triggered}, dovendo esaurire le
attività di I/O dei file descriptor risultati pronti per poter essere
rinotificati, la gestione elementare per cui li si trattano uno per uno in
sequenza può portare ad un effetto denominato \textit{starvation}
(``\textsl{carestia}'').  Si rischia cioè di concentrare le operazioni sul
primo file descriptor che dispone di molti dati, prolungandole per tempi molto
lunghi con un ritardo che può risultare eccessivo nei confronti di quelle da
eseguire sugli altri che verrebbero dopo.  Per evitare questo tipo di
problematiche viene consigliato di usare \func{epoll\_wait} per registrare un
elenco dei file descriptor da gestire, e di trattarli a turno in maniera più
equa.

Come già per \func{select} e \func{poll} anche per l'interfaccia di
\textit{epoll} si pone il problema di gestire l'attesa di segnali e di dati
contemporaneamente.  Valgono le osservazioni fatte in
sez.~\ref{sec:file_select}, e per poterlo fare di nuovo è necessaria una
variante della funzione di attesa che consenta di reimpostare all'uscita una
maschera di segnali, analoga alle estensioni \func{pselect} e \func{ppoll} che
abbiamo visto in precedenza per \func{select} e \func{poll}. In questo caso la
funzione di sistema si chiama \funcd{epoll\_pwait}\footnote{la funzione è
  stata introdotta a partire dal kernel 2.6.19, ed è, come tutta l'interfaccia
  di \textit{epoll}, specifica di Linux.} ed il suo prototipo è:

\begin{funcproto}{
\fhead{sys/epoll.h}
\fdecl{int epoll\_pwait(int epfd, struct epoll\_event * events, int maxevents, 
    int timeout, \\
\phantom{int epoll\_pwait(}const sigset\_t *sigmask)}

\fdesc{Attende che uno dei file descriptor osservati sia pronto, mascherando
    i segnali.}  }

{La funzione ritorna il numero di file descriptor pronti in caso di successo e
  $-1$ per un errore, nel qual caso \var{errno} assumerà uno dei valori già
  visti con \func{epoll\_wait}.

}  
\end{funcproto}

La funzione è del tutto analoga \func{epoll\_wait}, soltanto che alla sua
uscita viene ripristinata la maschera di segnali originale, sostituita durante
l'esecuzione da quella impostata con l'argomento \param{sigmask}; in sostanza
la chiamata a questa funzione è equivalente al seguente codice, eseguito però
in maniera atomica:
\includecodesnip{listati/epoll_pwait_means.c} 

Si tenga presente che come le precedenti funzioni di \textit{I/O multiplexing}
anche le funzioni dell'interfaccia di \textit{epoll} vengono utilizzate
prevalentemente con i server di rete, quando si devono tenere sotto
osservazione un gran numero di socket; per questo motivo rimandiamo anche in
questo caso la trattazione di un esempio concreto a quando avremo esaminato in
dettaglio le caratteristiche dei socket; in particolare si potrà trovare un
programma che utilizza questa interfaccia in sez.~\ref{sec:TCP_serv_epoll}.

\itindend{epoll}


\subsection{La notifica di eventi tramite file descriptor}
\label{sec:sig_signalfd_eventfd}

Abbiamo visto in sez.~\ref{sec:file_select} come il meccanismo classico delle
notifiche di eventi tramite i segnali, presente da sempre nei sistemi
unix-like, porti a notevoli problemi nell'interazione con le funzioni per
l'\textit{I/O multiplexing}, tanto che per evitare possibili \textit{race
  condition} sono state introdotte estensioni dello standard POSIX e funzioni
apposite come \func{pselect}, \func{ppoll} e \func{epoll\_pwait}.

Benché i segnali siano il meccanismo più usato per effettuare notifiche ai
processi, la loro interfaccia di programmazione, che comporta l'esecuzione di
una funzione di gestione in maniera asincrona e totalmente scorrelata
dall'ordinario flusso di esecuzione del processo, si è però dimostrata quasi
subito assai problematica. Oltre ai limiti relativi ai limiti al cosa si può
fare all'interno della funzione del gestore di segnali (quelli illustrati in
sez.~\ref{sec:sig_signal_handler}), c'è il problema più generale consistente
nel fatto che questa modalità di funzionamento cozza con altre interfacce di
programmazione previste dal sistema in cui si opera in maniera
\textsl{sincrona}, come quelle dell'\textit{I/O multiplexing} appena
illustrate.

In questo tipo di interfacce infatti ci si aspetta che il processo gestisca
gli eventi a cui deve reagire in maniera sincrona generando le opportune
risposte, mentre con l'arrivo di un segnale si possono avere interruzioni
asincrone in qualunque momento.  Questo comporta la necessità di dover
gestire, quando si deve tener conto di entrambi i tipi di eventi, le
interruzioni delle funzioni di attesa sincrone, ed evitare possibili
\textit{race conditions}. In sostanza se non ci fossero i segnali non ci
sarebbe da preoccuparsi, fintanto che si effettuano operazioni all'interno di
un processo, della non atomicità delle \textit{system call} lente che vengono
interrotte e devono essere riavviate.

Abbiamo visto però in sez.~\ref{sec:sig_real_time} che insieme ai segnali
\textit{real-time} sono state introdotte anche delle interfacce di gestione
sincrona dei segnali, con la funzione \func{sigwait} e le sue affini. Queste
funzioni consentono di gestire i segnali bloccando un processo fino alla
avvenuta ricezione e disabilitando l'esecuzione asincrona rispetto al resto
del programma del gestore del segnale. Questo consente di risolvere i problemi
di atomicità nella gestione degli eventi associati ai segnali, avendo tutto il
controllo nel flusso principale del programma, ottenendo così una gestione
simile a quella dell'\textit{I/O multiplexing}, ma non risolve i problemi
delle interazioni con quest'ultimo, perché o si aspetta la ricezione di un
segnale o si aspetta che un file descriptor sia accessibile e nessuna delle
rispettive funzioni consente di fare contemporaneamente entrambe le cose.

Per risolvere questo problema nello sviluppo del kernel si è pensato di
introdurre un meccanismo alternativo per la notifica dei segnali (esteso anche
ad altri eventi generici) che, ispirandosi di nuovo alla filosofia di Unix per
cui tutto è un file, consentisse di eseguire la notifica con l'uso di
opportuni file descriptor. Ovviamente si tratta di una funzionalità specifica
di Linux, non presente in altri sistemi unix-like, e non prevista da nessuno
standard, per cui va evitata se si ha a cuore la portabilità.

In sostanza, come per \func{sigwait}, si può disabilitare l'esecuzione di un
gestore in occasione dell'arrivo di un segnale, e rilevarne l'avvenuta
ricezione leggendone la notifica tramite l'uso di uno speciale file
descriptor. Trattandosi di un file descriptor questo potrà essere tenuto sotto
osservazione con le ordinarie funzioni dell'\textit{I/O multiplexing} (vale a
dire con le solite \func{select}, \func{poll} e \func{epoll\_wait}) allo
stesso modo di quelli associati a file o socket, per cui alla fine si potrà
attendere in contemporanea sia l'arrivo del segnale che la disponibilità di
accesso ai dati relativi a questi ultimi.

La funzione di sistema che permette di abilitare la ricezione dei segnali
tramite file descriptor è \funcd{signalfd},\footnote{in realtà quella
  riportata è l'interfaccia alla funzione fornita dalla \acr{glibc}, esistono
  infatti due versioni diverse della \textit{system call}; una prima versione,
  \func{signalfd}, introdotta nel kernel 2.6.22 e disponibile con la
  \acr{glibc} 2.8 che non supporta l'argomento \texttt{flags}, ed una seconda
  versione, \funcm{signalfd4}, introdotta con il kernel 2.6.27 e che è quella
  che viene sempre usata a partire dalla \acr{glibc} 2.9, che prende un
  argomento aggiuntivo \code{size\_t sizemask} che indica la dimensione della
  maschera dei segnali, il cui valore viene impostato automaticamente dalla
  \acr{glibc}.}  il cui prototipo è:

\begin{funcproto}{
\fhead{sys/signalfd.h}
\fdecl{int signalfd(int fd, const sigset\_t *mask, int flags)}

\fdesc{Crea o modifica un file descriptor per la ricezione dei segnali.}
}

{La funzione ritorna un numero di file descriptor in caso di successo e $-1$
  per un errore, nel qual caso \var{errno} assumerà uno dei valori:
  \begin{errlist}
  \item[\errcode{EBADF}] il valore \param{fd} non indica un file descriptor.
  \item[\errcode{EINVAL}] il file descriptor \param{fd} non è stato ottenuto
    con \func{signalfd} o il valore di \param{flags} non è valido.
  \item[\errcode{ENODEV}] il kernel non può montare internamente il
    dispositivo per la gestione anonima degli \textit{inode}
    associati al file descriptor.
  \item[\errcode{ENOMEM}] non c'è memoria sufficiente per creare un nuovo file
    descriptor di \func{signalfd}.
  \end{errlist}
  ed inoltre \errval{EMFILE} e \errval{ENFILE} nel loro significato generico.
  
}  
\end{funcproto}

La funzione consente di creare o modificare le caratteristiche di un file
descriptor speciale su cui ricevere le notifiche della ricezione di
segnali. Per creare ex-novo uno di questi file descriptor è necessario passare
$-1$ come valore per l'argomento \param{fd}, ogni altro valore positivo verrà
invece interpretato come il numero del file descriptor (che deve esser stato
precedentemente creato sempre con \func{signalfd}) di cui si vogliono
modificare le caratteristiche. Nel primo caso la funzione ritornerà il valore
del nuovo file descriptor e nel secondo caso il valore indicato
con \param{fd}, in caso di errore invece verrà restituito $-1$.

L'elenco dei segnali che si vogliono gestire con \func{signalfd} deve essere
specificato tramite l'argomento \param{mask}. Questo deve essere passato come
puntatore ad una maschera di segnali creata con l'uso delle apposite macro già
illustrate in sez.~\ref{sec:sig_sigset}. La maschera deve indicare su quali
segnali si intende operare con \func{signalfd}; l'elenco può essere modificato
con una successiva chiamata a \func{signalfd}. Dato che \signal{SIGKILL} e
\signal{SIGSTOP} non possono essere intercettati (e non prevedono neanche la
possibilità di un gestore) un loro inserimento nella maschera verrà ignorato
senza generare errori.

L'argomento \param{flags} consente di impostare direttamente in fase di
creazione due flag per il file descriptor analoghi a quelli che si possono
impostare con una creazione ordinaria con \func{open}, evitando una
impostazione successiva con \func{fcntl} (si ricordi che questo è un argomento
aggiuntivo, introdotto con la versione fornita a partire dal kernel 2.6.27,
per kernel precedenti il valore deve essere nullo).  L'argomento deve essere
specificato come maschera binaria dei valori riportati in
tab.~\ref{tab:signalfd_flags}.

\begin{table}[htb]
  \centering
  \footnotesize
  \begin{tabular}[c]{|l|p{8cm}|}
    \hline
    \textbf{Valore}  & \textbf{Significato} \\
    \hline
    \hline
    \constd{SFD\_NONBLOCK}&imposta sul file descriptor il flag di
                           \const{O\_NONBLOCK} per renderlo non bloccante.\\ 
    \constd{SFD\_CLOEXEC}& imposta il flag di \const{O\_CLOEXEC} per la
                           chiusura automatica del file descriptor nella
                           esecuzione di \func{exec}.\\
    \hline    
  \end{tabular}
  \caption{Valori dell'argomento \param{flags} per la funzione \func{signalfd}
    che consentono di impostare i flag del file descriptor.} 
  \label{tab:signalfd_flags}
\end{table}

Si tenga presente che la chiamata a \func{signalfd} non disabilita la gestione
ordinaria dei segnali indicati da \param{mask}; questa, se si vuole effettuare
la ricezione tramite il file descriptor, dovrà essere disabilitata
esplicitamente bloccando gli stessi segnali con \func{sigprocmask}, altrimenti
verranno comunque eseguite le azioni di default (o un eventuale gestore
installato in precedenza). Il blocco non ha invece nessun effetto sul file
descriptor restituito da \func{signalfd}, dal quale sarà possibile pertanto
ricevere qualunque segnale, anche se questo risultasse bloccato.

Si tenga presente inoltre che la lettura di una struttura
\struct{signalfd\_siginfo} relativa ad un segnale pendente è equivalente alla
esecuzione di un gestore, vale a dire che una volta letta il segnale non sarà
più pendente e non potrà essere ricevuto, qualora si ripristino le normali
condizioni di gestione, né da un gestore, né dalla funzione \func{sigwaitinfo}.

Come anticipato, essendo questo lo scopo principale della nuova interfaccia,
il file descriptor può essere tenuto sotto osservazione tramite le funzioni
dell'\textit{I/O multiplexing} (vale a dire con le solite \func{select},
\func{poll} e \func{epoll\_wait}), e risulterà accessibile in lettura quando
uno o più dei segnali indicati tramite \param{mask} sarà pendente.

La funzione può essere chiamata più volte dallo stesso processo, consentendo
così di tenere sotto osservazione segnali diversi tramite file descriptor
diversi. Inoltre è anche possibile tenere sotto osservazione lo stesso segnale
con più file descriptor, anche se la pratica è sconsigliata; in tal caso la
ricezione del segnale potrà essere effettuata con una lettura da uno qualunque
dei file descriptor a cui è associato, ma questa potrà essere eseguita
soltanto una volta. Questo significa che tutti i file descriptor su cui è
presente lo stesso segnale risulteranno pronti in lettura per le funzioni di
\textit{I/O multiplexing}, ma una volta eseguita la lettura su uno di essi il
segnale sarà considerato ricevuto ed i relativi dati non saranno più
disponibili sugli altri file descriptor, che (a meno di una ulteriore
occorrenza del segnale nel frattempo) di non saranno più pronti.

Quando il file descriptor per la ricezione dei segnali non serve più potrà
essere chiuso con \func{close} liberando tutte le risorse da esso allocate. In
tal caso qualora vi fossero segnali pendenti questi resteranno tali, e
potranno essere ricevuti normalmente una volta che si rimuova il blocco
imposto con \func{sigprocmask}.

Oltre a poter essere usato con le funzioni dell'\textit{I/O multiplexing}, il
file descriptor restituito da \func{signalfd} cerca di seguire la semantica di
un sistema unix-like anche con altre \textit{system call}; in particolare esso
resta aperto (come ogni altro file descriptor) attraverso una chiamata ad
\func{exec}, a meno che non lo si sia creato con il flag di
\const{SFD\_CLOEXEC} o si sia successivamente impostato il
\textit{close-on-exec} con \func{fcntl}. Questo comportamento corrisponde
anche alla ordinaria semantica relativa ai segnali bloccati, che restano
pendenti attraverso una \func{exec}.

Analogamente il file descriptor resta sempre disponibile attraverso una
\func{fork} per il processo figlio, che ne riceve una copia; in tal caso però
il figlio potrà leggere dallo stesso soltanto i dati relativi ai segnali
ricevuti da lui stesso. Nel caso di \textit{thread} viene nuovamente seguita
la semantica ordinaria dei segnali, che prevede che un singolo \textit{thread}
possa ricevere dal file descriptor solo le notifiche di segnali inviati
direttamente a lui o al processo in generale, e non quelli relativi ad altri
\textit{thread} appartenenti allo stesso processo.

L'interfaccia fornita da \func{signalfd} prevede che la ricezione dei segnali
sia eseguita leggendo i dati relativi ai segnali pendenti dal file descriptor
restituito dalla funzione con una normalissima \func{read}.  Qualora non vi
siano segnali pendenti la \func{read} si bloccherà a meno di non aver
impostato la modalità di I/O non bloccante sul file descriptor, o direttamente
in fase di creazione con il flag \const{SFD\_NONBLOCK}, o in un momento
successivo con \func{fcntl}.  

\begin{figure}[!htb]
  \footnotesize \centering
  \begin{minipage}[c]{0.95\textwidth}
    \includestruct{listati/signalfd_siginfo.h}
  \end{minipage} 
  \normalsize 
  \caption{La struttura \structd{signalfd\_siginfo}, restituita in lettura da
    un file descriptor creato con \func{signalfd}.}
  \label{fig:signalfd_siginfo}
\end{figure}

I dati letti dal file descriptor vengono scritti sul buffer indicato come
secondo argomento di \func{read} nella forma di una sequenza di una o più
strutture \struct{signalfd\_siginfo} (la cui definizione si è riportata in
fig.~\ref{fig:signalfd_siginfo}) a seconda sia della dimensione del buffer che
del numero di segnali pendenti. Per questo motivo il buffer deve essere almeno
di dimensione pari a quella di \struct{signalfd\_siginfo}, qualora sia di
dimensione maggiore potranno essere letti in unica soluzione i dati relativi
ad eventuali più segnali pendenti, fino al numero massimo di strutture
\struct{signalfd\_siginfo} che possono rientrare nel buffer.

\begin{figure}[!htb]
  \footnotesize \centering
  \begin{minipage}[c]{\codesamplewidth}
    \includecodesample{listati/FifoReporter-init.c}
  \end{minipage} 
  \normalsize 
  \caption{Sezione di inizializzazione del codice del programma
    \file{FifoReporter.c}.}
  \label{fig:fiforeporter_code_init}
\end{figure}

Il contenuto di \struct{signalfd\_siginfo} ricalca da vicino quella
dell'analoga struttura \struct{siginfo\_t} (illustrata in
fig.~\ref{fig:sig_siginfo_t}) usata dall'interfaccia ordinaria dei segnali, e
restituisce dati simili. Come per \struct{siginfo\_t} i campi che vengono
avvalorati dipendono dal tipo di segnale e ricalcano i valori che abbiamo già
illustrato in sez.~\ref{sec:sig_sigaction}.\footnote{si tenga presente però
  che per un bug i kernel fino al 2.6.25 non avvalorano correttamente i campi
  \var{ssi\_ptr} e \var{ssi\_int} per segnali inviati con \func{sigqueue}.}

Come esempio di questa nuova interfaccia ed anche come esempio di applicazione
della interfaccia di \textit{epoll}, si è scritto un programma elementare che
stampi sullo \textit{standard output} sia quanto viene scritto da terzi su una
\textit{named fifo}, che l'avvenuta ricezione di alcuni segnali.  Il codice
completo si trova al solito nei sorgenti allegati alla guida (nel file
\texttt{FifoReporter.c}).

In fig.~\ref{fig:fiforeporter_code_init} si è riportata la parte iniziale del
programma in cui vengono effettuate le varie inizializzazioni necessarie per
l'uso di \textit{epoll} e \func{signalfd}, a partire (\texttt{\small 12-16})
dalla definizione delle varie variabili e strutture necessarie. Al solito si è
tralasciata la parte dedicata alla decodifica delle opzioni che consentono ad
esempio di cambiare il nome del file associato alla \textit{fifo}.

Il primo passo (\texttt{\small 19-20}) è la creazione di un file descriptor
\texttt{epfd} di \textit{epoll} con \func{epoll\_create} che è quello che
useremo per il controllo degli altri.  É poi necessario disabilitare la
ricezione dei segnali (nel caso \signal{SIGINT}, \signal{SIGQUIT} e
\signal{SIGTERM}) per i quali si vuole la notifica tramite file
descriptor. Per questo prima li si inseriscono (\texttt{\small 22-25}) in una
maschera di segnali \texttt{sigmask} che useremo con (\texttt{\small 26})
\func{sigprocmask} per disabilitarli.  Con la stessa maschera si potrà per
passare all'uso (\texttt{\small 28-29}) di \func{signalfd} per abilitare la
notifica sul file descriptor \var{sigfd}. Questo poi (\texttt{\small 30-33})
dovrà essere aggiunto con \func{epoll\_ctl} all'elenco di file descriptor
controllati con \texttt{epfd}.

Occorrerà infine (\texttt{\small 35-38}) creare la \textit{named fifo} se
questa non esiste ed aprirla per la lettura (\texttt{\small 39-40}); una volta
fatto questo sarà necessario aggiungere il relativo file descriptor
(\var{fifofd}) a quelli osservati da \textit{epoll} in maniera del tutto
analoga a quanto fatto con quello relativo alla notifica dei segnali.

\begin{figure}[!htb]
  \footnotesize \centering
  \begin{minipage}[c]{\codesamplewidth}
    \includecodesample{listati/FifoReporter-main.c}
  \end{minipage} 
  \normalsize 
  \caption{Ciclo principale del codice del programma \file{FifoReporter.c}.}
  \label{fig:fiforeporter_code_body}
\end{figure}

Una volta completata l'inizializzazione verrà eseguito indefinitamente il
ciclo principale del programma (\texttt{\small 2-45}) che si è riportato in
fig.~\ref{fig:fiforeporter_code_body}, fintanto che questo non riceva un
segnale di \signal{SIGINT} (ad esempio con la pressione di \texttt{C-c}). Il
ciclo prevede che si attenda (\texttt{\small 2-3}) la presenza di un file
descriptor pronto in lettura con \func{epoll\_wait} (si ricordi che entrambi i
file descriptor \var{fifofd} e \var{sigfd} sono stati posti in osservazioni
per eventi di tipo \const{EPOLLIN}) che si bloccherà fintanto che non siano
stati scritti dati sulla \textit{fifo} o che non sia arrivato un
segnale.\footnote{per semplificare il codice non si è trattato il caso in cui
  \func{epoll\_wait} viene interrotta da un segnale, assumendo che tutti
  quelli che possano interessare siano stati predisposti per la notifica
  tramite file descriptor, per gli altri si otterrà semplicemente l'uscita dal
  programma.}

Anche se in questo caso i file descriptor pronti possono essere al più due, si
è comunque adottato un approccio generico in cui questi verranno letti
all'interno di un opportuno ciclo (\texttt{\small 5-44}) sul numero
restituito da \func{epoll\_wait}, esaminando i risultati presenti nel vettore
\var{events} all'interno di una catena di condizionali alternativi sul valore
del file descriptor riconosciuto come pronto, controllando cioè a quale dei
due file descriptor possibili corrisponde il campo relativo,
\var{events[i].data.fd}.

Il primo condizionale (\texttt{\small 6-24}) è relativo al caso che si sia
ricevuto un segnale e che il file descriptor pronto corrisponda
(\texttt{\small 6}) a \var{sigfd}. Dato che in generale si possono ricevere
anche notifiche relativi a più di un singolo segnale, si è scelto di leggere
una struttura \struct{signalfd\_siginfo} alla volta, eseguendo la lettura
all'interno di un ciclo (\texttt{\small 8-24}) che prosegue fintanto che vi
siano dati da leggere.

Per questo ad ogni lettura si esamina (\texttt{\small 9-14}) se il valore di
ritorno della funzione \func{read} è negativo, uscendo dal programma
(\texttt{\small 11}) in caso di errore reale, o terminando il ciclo
(\texttt{\small 13}) con un \texttt{break} qualora si ottenga un errore di
\errcode{EAGAIN} per via dell'esaurimento dei dati. Si ricordi infatti come
sia la \textit{fifo} che il file descriptor per i segnali siano stati aperti in
modalità non-bloccante, come previsto per l’\textit{I/O multiplexing},
pertanto ci si aspetta di ricevere un errore di \errcode{EAGAIN} quando non vi
saranno più dati da leggere.

In presenza di dati invece il programma proseguirà l'esecuzione stampando
(\texttt{\small 19-20}) il nome del segnale ottenuto all'interno della
struttura \struct{signalfd\_siginfo} letta in \var{siginf} ed il \textit{pid}
del processo da cui lo ha ricevuto;\footnote{per la stampa si è usato il
  vettore \var{sig\_names} a ciascun elemento del quale corrisponde il nome
  del segnale avente il numero corrispondente, la cui definizione si è omessa
  dal codice di fig.~\ref{fig:fiforeporter_code_init} per brevità.} inoltre
(\texttt{\small 21-24}) si controllerà anche se il segnale ricevuto è
\signal{SIGINT}, che si è preso come segnale da utilizzare per la terminazione
del programma, che verrà eseguita dopo aver rimosso il file della \textit{name
  fifo}.
 
Il secondo condizionale (\texttt{\small 26-39}) è invece relativo al caso in
cui ci siano dati pronti in lettura sulla \textit{fifo} e che il file
descriptor pronto corrisponda (\texttt{\small 26}) a \var{fifofd}. Di nuovo si
effettueranno le letture in un ciclo (\texttt{\small 28-39}) ripetendole fin
tanto che la funzione \func{read} non restituisce un errore di
\errcode{EAGAIN} (\texttt{\small 29-35}). Il procedimento è lo stesso adottato
per il file descriptor associato al segnale, in cui si esce dal programma in
caso di errore reale, in questo caso però alla fine dei dati prima di uscire
si stampa anche (\texttt{\small 32}) un messaggio di chiusura.

Se invece vi sono dati validi letti dalla \textit{fifo} si inserirà
(\texttt{\small 36}) una terminazione di stringa sul buffer e si stamperà il
tutto (\texttt{\small 37-38}) sullo \textit{standard output}. L'ultimo
condizionale (\texttt{\small 40-44}) è semplicemente una condizione di cattura
per una eventualità che comunque non dovrebbe mai verificarsi, e che porta
alla uscita dal programma con una opportuna segnalazione di errore.

A questo punto si potrà eseguire il comando lanciandolo su un terminale, ed
osservarne le reazioni agli eventi generati da un altro terminale; lanciando
il programma otterremo qualcosa del tipo:
\begin{Console}
piccardi@hain:~/gapil/sources$ \textbf{./a.out} 
FifoReporter starting, pid 4568
\end{Console}
%$
e scrivendo qualcosa sull'altro terminale con:
\begin{Console}
root@hain:~# \textbf{echo prova > /tmp/reporter.fifo}  
\end{Console}
si otterrà:
\begin{Console}
Message from fifo:
prova
end message
\end{Console}
mentre inviando un segnale:
\begin{Console}
root@hain:~# \textbf{kill 4568}
\end{Console}
si avrà:
\begin{Console}
Signal received:
Got SIGTERM       
From pid 3361
\end{Console}
ed infine premendo \texttt{C-\bslash} sul terminale in cui è in esecuzione si
vedrà:
\begin{Console}
^\\Signal received:
Got SIGQUIT       
From pid 0
\end{Console}
e si potrà far uscire il programma con \texttt{C-c} ottenendo:
\begin{Console}
^CSignal received:
Got SIGINT        
From pid 0
SIGINT means exit
\end{Console}

Lo stesso paradigma di notifica tramite file descriptor usato per i segnali è
stato adottato anche per i timer. In questo caso, rispetto a quanto visto in
sez.~\ref{sec:sig_timer_adv}, la scadenza di un timer potrà essere letta da un
file descriptor senza dover ricorrere ad altri meccanismi di notifica come un
segnale o un \textit{thread}. Di nuovo questo ha il vantaggio di poter
utilizzare le funzioni dell'\textit{I/O multiplexing} per attendere allo
stesso tempo la disponibilità di dati o la ricezione della scadenza di un
timer. In realtà per questo sarebbe già sufficiente \func{signalfd} per
ricevere i segnali associati ai timer, ma la nuova interfaccia semplifica
notevolmente la gestione e consente di fare tutto con una sola \textit{system
  call}.

Le funzioni di questa nuova interfaccia ricalcano da vicino la struttura delle
analoghe versioni ordinarie introdotte con lo standard POSIX.1-2001, che
abbiamo già illustrato in sez.~\ref{sec:sig_timer_adv}.\footnote{questa
  interfaccia è stata introdotta in forma considerata difettosa con il kernel
  2.6.22, per cui è stata immediatamente tolta nel successivo 2.6.23 e
  reintrodotta in una forma considerata adeguata nel kernel 2.6.25, il
  supporto nella \acr{glibc} è stato introdotto a partire dalla versione
  2.8.6, la versione del kernel 2.6.22, presente solo su questo kernel, non è
  supportata e non deve essere usata.} La prima funzione di sistema prevista,
quella che consente di creare un timer, è \funcd{timerfd\_create}, il cui
prototipo è:

\begin{funcproto}{
\fhead{sys/timerfd.h}
\fdecl{int timerfd\_create(int clockid, int flags)}

\fdesc{Crea un timer associato ad un file descriptor di notifica.}
}

{La funzione ritorna un numero di file descriptor in caso di successo e $-1$
  per un errore, nel qual caso \var{errno} assumerà uno dei valori:
  \begin{errlist}
  \item[\errcode{EINVAL}] l'argomento \param{clockid} non è
    \const{CLOCK\_MONOTONIC} o \const{CLOCK\_REALTIME}, o
    l'argomento \param{flag} non è valido, o è diverso da zero per kernel
    precedenti il 2.6.27.
  \item[\errcode{ENODEV}] il kernel non può montare internamente il
    dispositivo per la gestione anonima degli \textit{inode} associati al file
    descriptor.
  \item[\errcode{ENOMEM}] non c'è memoria sufficiente per creare un nuovo file
    descriptor di \func{signalfd}.
  \end{errlist}
  ed inoltre \errval{EMFILE} e \errval{ENFILE} nel loro significato generico.
}  
\end{funcproto}

La funzione prende come primo argomento un intero che indica il tipo di
orologio a cui il timer deve fare riferimento, i valori sono gli stessi delle
funzioni dello standard POSIX-1.2001 già illustrati in
tab.~\ref{tab:sig_timer_clockid_types}, ma al momento i soli utilizzabili sono
\const{CLOCK\_REALTIME} e \const{CLOCK\_MONOTONIC}. L'argomento \param{flags},
come l'analogo di \func{signalfd}, consente di impostare i flag per l'I/O non
bloccante ed il \textit{close-on-exec} sul file descriptor
restituito,\footnote{il flag è stato introdotto a partire dal kernel 2.6.27,
  per le versioni precedenti deve essere passato un valore nullo.} e deve
essere specificato come una maschera binaria delle costanti riportate in
tab.~\ref{tab:timerfd_flags}.

\begin{table}[htb]
  \centering
  \footnotesize
  \begin{tabular}[c]{|l|p{8cm}|}
    \hline
    \textbf{Valore}  & \textbf{Significato} \\
    \hline
    \hline
    \constd{TFD\_NONBLOCK}& imposta sul file descriptor il flag di
                            \const{O\_NONBLOCK} per renderlo non bloccante.\\ 
    \constd{TFD\_CLOEXEC} & imposta il flag di \const{O\_CLOEXEC} per la
                            chiusura automatica del file descriptor nella
                            esecuzione di \func{exec}.\\
    \hline    
  \end{tabular}
  \caption{Valori dell'argomento \param{flags} per la funzione
    \func{timerfd\_create} che consentono di impostare i flag del file
    descriptor.}  
  \label{tab:timerfd_flags}
\end{table}

In caso di successo la funzione restituisce un file descriptor sul quale
verranno notificate le scadenze dei timer. Come per quelli restituiti da
\func{signalfd} anche questo file descriptor segue la semantica dei sistemi
unix-like, in particolare resta aperto attraverso una \func{exec} (a meno che
non si sia impostato il flag di \textit{close-on exec} con
\const{TFD\_CLOEXEC}) e viene duplicato attraverso una \func{fork}; questa
ultima caratteristica comporta però che anche il figlio può utilizzare i dati
di un timer creato nel padre, a differenza di quanto avviene invece con i
timer impostati con le funzioni ordinarie. Si ricordi infatti che, come
illustrato in sez.~\ref{sec:proc_fork}, allarmi, timer e segnali pendenti nel
padre vengono cancellati per il figlio dopo una \func{fork}.

Una volta creato il timer con \func{timerfd\_create} per poterlo utilizzare
occorre \textsl{armarlo} impostandone un tempo di scadenza ed una eventuale
periodicità di ripetizione, per farlo si usa una funzione di sistema omologa
di \func{timer\_settime} per la nuova interfaccia; questa è
\funcd{timerfd\_settime} ed il suo prototipo è:

\begin{funcproto}{
\fhead{sys/timerfd.h}
\fdecl{int timerfd\_settime(int fd, int flags,
                           const struct itimerspec *new\_value,\\
\phantom{int timerfd\_settime(}struct itimerspec *old\_value)}

\fdesc{Arma un timer associato ad un file descriptor di notifica.}
}

{La funzione ritorna un numero di file descriptor in caso di successo e $-1$
  per un errore, nel qual caso \var{errno} assumerà uno dei valori:
  \begin{errlist}
  \item[\errcode{EBADF}] l'argomento \param{fd} non corrisponde ad un file
    descriptor. 
  \item[\errcode{EFAULT}] o \param{new\_value} o \param{old\_value} non sono
    puntatori validi.
  \item[\errcode{EINVAL}] il file descriptor \param{fd} non è stato ottenuto
    con \func{timerfd\_create}, o i valori di \param{flag} o dei campi
    \var{tv\_nsec} in \param{new\_value} non sono validi.
  \end{errlist}
}  
\end{funcproto}

In questo caso occorre indicare su quale timer si intende operare specificando
come primo argomento il file descriptor ad esso associato, che deve essere
stato ottenuto da una precedente chiamata a \func{timerfd\_create}. I restanti
argomenti sono del tutto analoghi a quelli della omologa funzione
\func{timer\_settime}, e prevedono l'uso di strutture \struct{itimerspec}
(vedi fig.~\ref{fig:struct_itimerspec}) per le indicazioni di temporizzazione.

I valori ed il significato di questi argomenti sono gli stessi che sono già
stati illustrati in dettaglio in sez.~\ref{sec:sig_timer_adv} e non staremo a
ripetere quanto detto in quell'occasione; per brevità si ricordi che
con \param{new\_value.it\_value} si indica la prima scadenza del timer e
con \param{new\_value.it\_interval} la sua periodicità.  L'unica differenza
riguarda l'argomento \param{flags} che serve sempre ad indicare se il tempo di
scadenza del timer è da considerarsi relativo o assoluto rispetto al valore
corrente dell'orologio associato al timer, ma che in questo caso ha come
valori possibili rispettivamente soltanto $0$ e \constd{TFD\_TIMER\_ABSTIME}
(l'analogo di \const{TIMER\_ABSTIME}).

L'ultima funzione di sistema prevista dalla nuova interfaccia è
\funcd{timerfd\_gettime}, che è l'analoga di \func{timer\_gettime}, il suo
prototipo è:

\begin{funcproto}{
\fhead{sys/timerfd.h}
\fdecl{int timerfd\_gettime(int fd, struct itimerspec *curr\_value)}

\fdesc{Legge l'impostazione di un timer associato ad un file descriptor di
  notifica.} 
}

{La funzione ritorna un numero di file descriptor in caso di successo e $-1$
  per un errore, nel qual caso \var{errno} assumerà uno dei valori:
  \begin{errlist}
  \item[\errcode{EBADF}] l'argomento \param{fd} non corrisponde ad un file
    descriptor. 
  \item[\errcode{EINVAL}] il file descriptor \param{fd} non è stato ottenuto
    con \func{timerfd\_create}.
  \item[\errcode{EFAULT}] o \param{curr\_value} non è un puntatore valido.
  \end{errlist}
}  
\end{funcproto}

La funzione consente di rileggere le impostazioni del timer associato al file
descriptor \param{fd} nella struttura \struct{itimerspec} puntata
da \param{curr\_value}. Il campo \var{it\_value} riporta il tempo rimanente
alla prossima scadenza del timer, che viene sempre espresso in forma relativa,
anche se lo si è armato specificando \const{TFD\_TIMER\_ABSTIME}. Un valore
nullo (di entrambi i campi di \var{it\_value}) indica invece che il timer non
è stato ancora armato. Il campo \var{it\_interval} riporta la durata
dell'intervallo di ripetizione del timer, ed un valore nullo (di entrambi i
campi) indica che il timer è stato impostato per scadere una sola volta.

Il timer creato con \func{timerfd\_create} notificherà la sua scadenza
rendendo pronto per la lettura il file descriptor ad esso associato, che
pertanto potrà essere messo sotto controllo con una qualunque delle varie
funzioni dell'I/O multiplexing viste in precedenza. Una volta che il file
descriptor risulta pronto sarà possibile leggere il numero di volte che il
timer è scaduto con una ordinaria \func{read}. 

La funzione legge il valore in un dato di tipo \typed{uint64\_t}, e necessita
pertanto che le si passi un buffer di almeno 8 byte, fallendo con
\errval{EINVAL} in caso contrario, in sostanza la lettura deve essere
effettuata con una istruzione del tipo:
\includecodesnip{listati/readtimerfd.c} 

Il valore viene restituito da \func{read} seguendo l'ordinamento dei bit
(\textit{big-endian} o \textit{little-endian}) nativo della macchina in uso,
ed indica il numero di volte che il timer è scaduto dall'ultima lettura
eseguita con successo, o, se lo si legge per la prima volta, da quando lo si è
impostato con \func{timerfd\_settime}. Se il timer non è scaduto la funzione
si blocca fino alla prima scadenza, a meno di non aver creato il file
descriptor in modalità non bloccante con \const{TFD\_NONBLOCK} o aver
impostato la stessa con \func{fcntl}, nel qual caso fallisce con l'errore di
\errval{EAGAIN}.


% TODO trattare qui eventfd introdotto con il 2.6.22 


\section{L'accesso \textsl{asincrono} ai file}
\label{sec:file_asyncronous_operation}

Benché l'\textit{I/O multiplexing} sia stata la prima, e sia tutt'ora una fra
le più diffuse modalità di gestire l'I/O in situazioni complesse in cui si
debba operare su più file contemporaneamente, esistono altre modalità di
gestione delle stesse problematiche. In particolare sono importanti in questo
contesto le modalità di accesso ai file eseguibili in maniera
\textsl{asincrona}, quelle cioè in cui un processo non deve bloccarsi in
attesa della disponibilità dell'accesso al file, ma può proseguire
nell'esecuzione utilizzando invece un meccanismo di notifica asincrono (di
norma un segnale, ma esistono anche altre interfacce, come \textit{inotify}),
per essere avvisato della possibilità di eseguire le operazioni di I/O volute.


\subsection{Il \textit{Signal driven I/O}}
\label{sec:signal_driven_io}

\itindbeg{signal~driven~I/O}

Abbiamo accennato in sez.~\ref{sec:file_open_close} che è definito un flag
\const{O\_ASYNC}, che consentirebbe di aprire un file in modalità asincrona,
anche se in realtà è opportuno attivare in un secondo tempo questa modalità
impostando questo flag attraverso l'uso di \func{fcntl} con il comando
\const{F\_SETFL} (vedi sez.~\ref{sec:file_fcntl_ioctl}).\footnote{l'uso del
  flag di \const{O\_ASYNC} e dei comandi \const{F\_SETOWN} e \const{F\_GETOWN}
  per \func{fcntl} è specifico di Linux e BSD.}  In realtà parlare di apertura
in modalità asincrona non significa che le operazioni di lettura o scrittura
del file vengono eseguite in modo asincrono (tratteremo questo, che è ciò che
più propriamente viene chiamato \textsl{I/O asincrono}, in
sez.~\ref{sec:file_asyncronous_io}), quanto dell'attivazione un meccanismo di
notifica asincrona delle variazione dello stato del file descriptor aperto in
questo modo.

Quello che succede è che per tutti i file posti in questa modalità il sistema
genera un apposito segnale, \signal{SIGIO}, tutte le volte che diventa
possibile leggere o scrivere dal file descriptor; si tenga presente però che
essa non è utilizzabile con i file ordinari ma solo con socket, file di
terminale o pseudo terminale, ed anche, a partire dal kernel 2.6, per
\textit{fifo} e \textit{pipe}. Inoltre è possibile, come illustrato in
sez.~\ref{sec:file_fcntl_ioctl}, selezionare con il comando \const{F\_SETOWN}
di \func{fcntl} quale processo o quale gruppo di processi dovrà ricevere il
segnale. In questo modo diventa possibile effettuare le operazioni di I/O in
risposta alla ricezione del segnale, e non ci sarà più la necessità di restare
bloccati in attesa della disponibilità di accesso ai file.

% TODO: per i thread l'uso di F_SETOWN ha un significato diverso

Per questo motivo Stevens, ed anche le pagine di manuale di Linux, chiamano
questa modalità ``\textit{Signal driven I/O}''.  Si tratta di un'altra
modalità di gestione dell'I/O, alternativa all'uso di
\textit{epoll},\footnote{anche se le prestazioni ottenute con questa tecnica
  sono inferiori, il vantaggio è che questa modalità è utilizzabile anche con
  kernel che non supportano \textit{epoll}, come quelli della serie 2.4,
  ottenendo comunque prestazioni superiori a quelle che si hanno con
  \func{poll} e \func{select}.} che consente di evitare l'uso delle funzioni
\func{poll} o \func{select} che, come illustrato in sez.~\ref{sec:file_epoll},
quando vengono usate con un numero molto grande di file descriptor, non hanno
buone prestazioni.

Tuttavia con l'implementazione classica dei segnali questa modalità di I/O
presenta notevoli problemi, dato che non è possibile determinare, quando i
file descriptor sono più di uno, qual è quello responsabile dell'emissione del
segnale. Inoltre dato che i segnali normali non si accodano (si ricordi quanto
illustrato in sez.~\ref{sec:sig_notification}), in presenza di più file
descriptor attivi contemporaneamente, più segnali emessi nello stesso momento
verrebbero notificati una volta sola.

Linux però supporta le estensioni POSIX.1b dei segnali \textit{real-time}, che
vengono accodati e che permettono di riconoscere il file descriptor che li ha
emessi.  In questo caso infatti si può fare ricorso alle informazioni
aggiuntive restituite attraverso la struttura \struct{siginfo\_t}, utilizzando
la forma estesa \var{sa\_sigaction} del gestore installata con il flag
\const{SA\_SIGINFO} (si riveda quanto illustrato in
sez.~\ref{sec:sig_sigaction}).

Per far questo però occorre utilizzare le funzionalità dei segnali
\textit{real-time} (vedi sez.~\ref{sec:sig_real_time}) impostando
esplicitamente con il comando \const{F\_SETSIG} di \func{fcntl} un segnale
\textit{real-time} da inviare in caso di I/O asincrono (il segnale predefinito
è \signal{SIGIO}). In questo caso il gestore, tutte le volte che riceverà
\const{SI\_SIGIO} come valore del campo \var{si\_code} di \struct{siginfo\_t},
troverà nel campo \var{si\_fd} il valore del file descriptor che ha generato
il segnale. Si noti che il valore di\var{si\_code} resta \const{SI\_SIGIO}
qualunque sia il segnale che si è associato all'I/O, in quanto indica che il
segnale è stato generato a causa di attività di I/O.

Un secondo vantaggio dell'uso dei segnali \textit{real-time} è che essendo
questi ultimi dotati di una coda di consegna ogni segnale sarà associato ad
uno solo file descriptor; inoltre sarà possibile stabilire delle priorità
nella risposta a seconda del segnale usato, dato che i segnali
\textit{real-time} supportano anche questa funzionalità. In questo modo si può
identificare immediatamente un file su cui l'accesso è diventato possibile
evitando completamente l'uso di funzioni come \func{poll} e \func{select},
almeno fintanto che non si satura la coda.

Se infatti si eccedono le dimensioni di quest'ultima, il kernel, non potendo
più assicurare il comportamento corretto per un segnale \textit{real-time},
invierà al suo posto un solo \signal{SIGIO}, su cui si saranno accumulati
tutti i segnali in eccesso, e si dovrà allora determinare con un ciclo quali
sono i file diventati attivi. L'unico modo per essere sicuri che questo non
avvenga è di impostare la lunghezza della coda dei segnali \textit{real-time}
ad una dimensione identica al valore massimo del numero di file descriptor
utilizzabili, vale a dire impostare il contenuto di
\sysctlfile{kernel/rtsig-max} allo stesso valore del contenuto di
\sysctlfile{fs/file-max}.

% TODO fare esempio che usa O_ASYNC

\itindend{signal~driven~I/O}



\subsection{I meccanismi di notifica asincrona.}
\label{sec:file_asyncronous_lease}

Una delle domande più frequenti nella programmazione in ambiente unix-like è
quella di come fare a sapere quando un file viene modificato. La risposta, o
meglio la non risposta, tanto che questa nelle Unix FAQ \cite{UnixFAQ} viene
anche chiamata una \textit{Frequently Unanswered Question}, è che
nell'architettura classica di Unix questo non è possibile. Al contrario di
altri sistemi operativi infatti un kernel unix-like classico non prevedeva
alcun meccanismo per cui un processo possa essere \textsl{notificato} di
eventuali modifiche avvenute su un file. 

Questo è il motivo per cui i demoni devono essere \textsl{avvisati} in qualche
modo se il loro file di configurazione è stato modificato, perché possano
rileggerlo e riconoscere le modifiche; in genere questo vien fatto inviandogli
un segnale di \signal{SIGHUP} che, per una convenzione adottata dalla gran
parte di detti programmi, causa la rilettura della configurazione.

Questa scelta è stata fatta perché provvedere un simile meccanismo a livello
generico per qualunque file comporterebbe un notevole aumento di complessità
dell'architettura della gestione dei file, il tutto per fornire una
funzionalità che serve soltanto in alcuni casi particolari. Dato che
all'origine di Unix i soli programmi che potevano avere una tale esigenza
erano i demoni, attenendosi a uno dei criteri base della progettazione, che
era di far fare al kernel solo le operazioni strettamente necessarie e
lasciare tutto il resto a processi in \textit{user space}, non era stata
prevista nessuna funzionalità di notifica.

Visto però il crescente interesse nei confronti di una funzionalità di questo
tipo, che è molto richiesta specialmente nello sviluppo dei programmi ad
interfaccia grafica quando si deve presentare all'utente lo stato del
filesystem, sono state successivamente introdotte delle estensioni che
permettessero la creazione di meccanismi di notifica più efficienti dell'unica
soluzione disponibile con l'interfaccia tradizionale, che è quella del
\textit{polling}.

Queste nuove funzionalità sono delle estensioni specifiche, non
standardizzate, che sono disponibili soltanto su Linux (anche se altri kernel
supportano meccanismi simili). Alcune di esse sono realizzate, e solo a
partire dalla versione 2.4 del kernel, attraverso l'uso di alcuni
\textsl{comandi} aggiuntivi per la funzione \func{fcntl} (vedi
sez.~\ref{sec:file_fcntl_ioctl}), che divengono disponibili soltanto se si è
definita la macro \macro{\_GNU\_SOURCE} prima di includere \headfile{fcntl.h}.

\itindbeg{file~lease} 

La prima di queste funzionalità è quella del cosiddetto \textit{file lease};
questo è un meccanismo che consente ad un processo, detto \textit{lease
  holder}, di essere notificato quando un altro processo, chiamato a sua volta
\textit{lease breaker}, cerca di eseguire una \func{open} o una
\func{truncate} sul file del quale l'\textit{holder} detiene il
\textit{lease}.  La notifica avviene in maniera analoga a come illustrato in
precedenza per l'uso di \const{O\_ASYNC}: di default viene inviato al
\textit{lease holder} il segnale \signal{SIGIO}, ma questo segnale può essere
modificato usando il comando \const{F\_SETSIG} di \func{fcntl} (anche in
questo caso si può rispecificare lo stesso \signal{SIGIO}).

Se si è fatto questo (ed in genere è opportuno farlo, come in precedenza, per
utilizzare segnali \textit{real-time}) e se inoltre si è installato il gestore
del segnale con \const{SA\_SIGINFO} si riceverà nel campo \var{si\_fd} della
struttura \struct{siginfo\_t} il valore del file descriptor del file sul quale
è stato compiuto l'accesso; in questo modo un processo può mantenere anche più
di un \textit{file lease}.

Esistono due tipi di \textit{file lease}: di lettura (\textit{read lease}) e
di scrittura (\textit{write lease}). Nel primo caso la notifica avviene quando
un altro processo esegue l'apertura del file in scrittura o usa
\func{truncate} per troncarlo. Nel secondo caso la notifica avviene anche se
il file viene aperto in lettura; in quest'ultimo caso però il \textit{lease}
può essere ottenuto solo se nessun altro processo ha aperto lo stesso file.

Come accennato in sez.~\ref{sec:file_fcntl_ioctl} il comando di \func{fcntl}
che consente di acquisire un \textit{file lease} è \const{F\_SETLEASE}, che
viene utilizzato anche per rilasciarlo. In tal caso il file
descriptor \param{fd} passato a \func{fcntl} servirà come riferimento per il
file su cui si vuole operare, mentre per indicare il tipo di operazione
(acquisizione o rilascio) occorrerà specificare come valore
dell'argomento \param{arg} di \func{fcntl} uno dei tre valori di
tab.~\ref{tab:file_lease_fctnl}.

\begin{table}[htb]
  \centering
  \footnotesize
  \begin{tabular}[c]{|l|l|}
    \hline
    \textbf{Valore}  & \textbf{Significato} \\
    \hline
    \hline
    \constd{F\_RDLCK} & Richiede un \textit{read lease}.\\
    \constd{F\_WRLCK} & Richiede un \textit{write lease}.\\
    \constd{F\_UNLCK} & Rilascia un \textit{file lease}.\\
    \hline    
  \end{tabular}
  \caption{Costanti per i tre possibili valori dell'argomento \param{arg} di
    \func{fcntl} quando usata con i comandi \const{F\_SETLEASE} e
    \const{F\_GETLEASE}.} 
  \label{tab:file_lease_fctnl}
\end{table}

Se invece si vuole conoscere lo stato di eventuali \textit{file lease}
occorrerà chiamare \func{fcntl} sul relativo file descriptor \param{fd} con il
comando \const{F\_GETLEASE}, e si otterrà indietro nell'argomento \param{arg}
uno dei valori di tab.~\ref{tab:file_lease_fctnl}, che indicheranno la
presenza del rispettivo tipo di \textit{lease}, o, nel caso di
\const{F\_UNLCK}, l'assenza di qualunque \textit{file lease}.

Si tenga presente che un processo può mantenere solo un tipo di \textit{lease}
su un file, e che un \textit{lease} può essere ottenuto solo su file di dati
(\textit{pipe} e dispositivi sono quindi esclusi). Inoltre un processo non
privilegiato può ottenere un \textit{lease} soltanto per un file appartenente
ad un \ids{UID} corrispondente a quello del processo. Soltanto un processo con
privilegi di amministratore (cioè con la capacità \const{CAP\_LEASE}, vedi
sez.~\ref{sec:proc_capabilities}) può acquisire \textit{lease} su qualunque
file.

Se su un file è presente un \textit{lease} quando il \textit{lease breaker}
esegue una \func{truncate} o una \func{open} che confligge con
esso,\footnote{in realtà \func{truncate} confligge sempre, mentre \func{open},
  se eseguita in sola lettura, non confligge se si tratta di un \textit{read
    lease}.} la funzione si blocca (a meno di non avere aperto il file con
\const{O\_NONBLOCK}, nel qual caso \func{open} fallirebbe con un errore di
\errcode{EWOULDBLOCK}) e viene eseguita la notifica al \textit{lease holder},
così che questo possa completare le sue operazioni sul file e rilasciare il
\textit{lease}.  In sostanza con un \textit{read lease} si rilevano i
tentativi di accedere al file per modificarne i dati da parte di un altro
processo, mentre con un \textit{write lease} si rilevano anche i tentativi di
accesso in lettura.  Si noti comunque che le operazioni di notifica avvengono
solo in fase di apertura del file e non sulle singole operazioni di lettura e
scrittura.

L'utilizzo dei \textit{file lease} consente al \textit{lease holder} di
assicurare la consistenza di un file, a seconda dei due casi, prima che un
altro processo inizi con le sue operazioni di scrittura o di lettura su di
esso. In genere un \textit{lease holder} che riceve una notifica deve
provvedere a completare le necessarie operazioni (ad esempio scaricare
eventuali buffer), per poi rilasciare il \textit{lease} così che il
\textit{lease breaker} possa eseguire le sue operazioni. Questo si fa con il
comando \const{F\_SETLEASE}, o rimuovendo il \textit{lease} con
\const{F\_UNLCK}, o, nel caso di \textit{write lease} che confligge con una
operazione di lettura, declassando il \textit{lease} a lettura con
\const{F\_RDLCK}.

Se il \textit{lease holder} non provvede a rilasciare il \textit{lease} entro
il numero di secondi specificato dal parametro di sistema mantenuto in
\sysctlfiled{fs/lease-break-time} sarà il kernel stesso a rimuoverlo o
declassarlo automaticamente (questa è una misura di sicurezza per evitare che
un processo blocchi indefinitamente l'accesso ad un file acquisendo un
\textit{lease}). Una volta che un \textit{lease} è stato rilasciato o
declassato (che questo sia fatto dal \textit{lease holder} o dal kernel è lo
stesso) le chiamate a \func{open} o \func{truncate} eseguite dal \textit{lease
  breaker} rimaste bloccate proseguono automaticamente.

Benché possa risultare utile per sincronizzare l'accesso ad uno stesso file da
parte di più processi, l'uso dei \textit{file lease} non consente comunque di
risolvere il problema di rilevare automaticamente quando un file o una
directory vengono modificati,\footnote{questa funzionalità venne aggiunta
  principalmente ad uso di Samba per poter facilitare l'emulazione del
  comportamento di Windows sui file, ma ad oggi viene considerata una
  interfaccia mal progettata ed il suo uso è fortemente sconsigliato a favore
  di \textit{inotify}.} che è quanto necessario ad esempio ai programma di
gestione dei file dei vari desktop grafici.

\itindbeg{dnotify}

Per risolvere questo problema a partire dal kernel 2.4 è stata allora creata
un'altra interfaccia,\footnote{si ricordi che anche questa è una interfaccia
  specifica di Linux che deve essere evitata se si vogliono scrivere programmi
  portabili, e che le funzionalità illustrate sono disponibili soltanto se è
  stata definita la macro \macro{\_GNU\_SOURCE}.} chiamata \textit{dnotify},
che consente di richiedere una notifica quando una directory, o uno qualunque
dei file in essa contenuti, viene modificato.  Come per i \textit{file lease}
la notifica avviene di default attraverso il segnale \signal{SIGIO}, ma se ne
può utilizzare un altro, e di nuovo, per le ragioni già esposte in precedenza,
è opportuno che si utilizzino dei segnali \textit{real-time}.  Inoltre, come
in precedenza, si potrà ottenere nel gestore del segnale il file descriptor
che è stato modificato tramite il contenuto della struttura
\struct{siginfo\_t}.

\itindend{file~lease}

\begin{table}[htb]
  \centering
  \footnotesize
  \begin{tabular}[c]{|l|p{8cm}|}
    \hline
    \textbf{Valore}  & \textbf{Significato} \\
    \hline
    \hline
    \constd{DN\_ACCESS} & Un file è stato acceduto, con l'esecuzione di una fra
                          \func{read}, \func{pread}, \func{readv}.\\ 
    \constd{DN\_MODIFY} & Un file è stato modificato, con l'esecuzione di una
                          fra \func{write}, \func{pwrite}, \func{writev}, 
                          \func{truncate}, \func{ftruncate}.\\ 
    \constd{DN\_CREATE} & È stato creato un file nella directory, con
                          l'esecuzione di una fra \func{open}, \func{creat},
                          \func{mknod}, \func{mkdir}, \func{link},
                          \func{symlink}, \func{rename} (da un'altra
                          directory).\\
    \constd{DN\_DELETE} & È stato cancellato un file dalla directory con
                          l'esecuzione di una fra \func{unlink}, \func{rename}
                          (su un'altra directory), \func{rmdir}.\\
    \constd{DN\_RENAME} & È stato rinominato un file all'interno della
                          directory (con \func{rename}).\\
    \constd{DN\_ATTRIB} & È stato modificato un attributo di un file con
                          l'esecuzione di una fra \func{chown}, \func{chmod},
                          \func{utime}.\\ 
    \constd{DN\_MULTISHOT}& Richiede una notifica permanente di tutti gli
                            eventi.\\ 
    \hline    
  \end{tabular}
  \caption{Le costanti che identificano le varie classi di eventi per i quali
    si richiede la notifica con il comando \const{F\_NOTIFY} di \func{fcntl}.} 
  \label{tab:file_notify}
\end{table}

Ci si può registrare per le notifiche dei cambiamenti al contenuto di una
certa directory eseguendo la funzione \func{fcntl} su un file descriptor
associato alla stessa con il comando \const{F\_NOTIFY}. In questo caso
l'argomento \param{arg} di \func{fcntl} serve ad indicare per quali classi
eventi si vuole ricevere la notifica, e prende come valore una maschera
binaria composta dall'OR aritmetico di una o più delle costanti riportate in
tab.~\ref{tab:file_notify}.

A meno di non impostare in maniera esplicita una notifica permanente usando il
valore \const{DN\_MULTISHOT}, la notifica è singola: viene cioè inviata una
sola volta quando si verifica uno qualunque fra gli eventi per i quali la si è
richiesta. Questo significa che un programma deve registrarsi un'altra volta
se desidera essere notificato di ulteriori cambiamenti. Se si eseguono diverse
chiamate con \const{F\_NOTIFY} e con valori diversi per \param{arg} questi
ultimi si \textsl{accumulano}; cioè eventuali nuovi classi di eventi
specificate in chiamate successive vengono aggiunte a quelle già impostate
nelle precedenti.  Se si vuole rimuovere la notifica si deve invece
specificare un valore nullo.

\itindbeg{inotify}

Il maggiore problema di \textit{dnotify} è quello della scalabilità: si deve
usare un file descriptor per ciascuna directory che si vuole tenere sotto
controllo, il che porta facilmente ad avere un eccesso di file aperti. Inoltre
quando la directory che si controlla è all'interno di un dispositivo
rimovibile, mantenere il relativo file descriptor aperto comporta
l'impossibilità di smontare il dispositivo e di rimuoverlo, il che in genere
complica notevolmente la gestione dell'uso di questi dispositivi.

Un altro problema è che l'interfaccia di \textit{dnotify} consente solo di
tenere sotto controllo il contenuto di una directory; la modifica di un file
viene segnalata, ma poi è necessario verificare di quale file si tratta
(operazione che può essere molto onerosa quando una directory contiene un gran
numero di file).  Infine l'uso dei segnali come interfaccia di notifica
comporta tutti i problemi di gestione visti in sez.~\ref{sec:sig_management} e
sez.~\ref{sec:sig_adv_control}.  Per tutta questa serie di motivi in generale
quella di \textit{dnotify} viene considerata una interfaccia di usabilità
problematica ed il suo uso oggi è fortemente sconsigliato.

\itindend{dnotify}

Per risolvere i problemi appena illustrati è stata introdotta una nuova
interfaccia per l'osservazione delle modifiche a file o directory, chiamata
\textit{inotify}.\footnote{l'interfaccia è disponibile a partire dal kernel
  2.6.13, le relative funzioni sono state introdotte nelle glibc 2.4.}  Anche
questa è una interfaccia specifica di Linux (pertanto non deve essere usata se
si devono scrivere programmi portabili), ed è basata sull'uso di una coda di
notifica degli eventi associata ad un singolo file descriptor, il che permette
di risolvere il principale problema di \textit{dnotify}.  La coda viene creata
attraverso la funzione di sistema \funcd{inotify\_init}, il cui prototipo è:

\begin{funcproto}{
\fhead{sys/inotify.h}
\fdecl{int inotify\_init(void)}
\fdesc{Inizializza una istanza di \textit{inotify}.}
}

{La funzione ritornaun file descriptor in caso di successo, o $-1$ in caso di
  errore, nel qual caso \var{errno} assumerà uno dei valori:
  \begin{errlist}
  \item[\errcode{EMFILE}] si è raggiunto il numero massimo di istanze di
    \textit{inotify} consentite all'utente.
  \item[\errcode{ENFILE}] si è raggiunto il massimo di file descriptor aperti
    nel sistema.
  \item[\errcode{ENOMEM}] non c'è sufficiente memoria nel kernel per creare
    l'istanza.
  \end{errlist}
}
\end{funcproto}

La funzione non prende alcun argomento; inizializza una istanza di
\textit{inotify} e restituisce un file descriptor attraverso il quale verranno
effettuate le operazioni di notifica; si tratta di un file descriptor speciale
che non è associato a nessun file su disco, e che viene utilizzato solo per
notificare gli eventi che sono stati posti in osservazione. Per evitare abusi
delle risorse di sistema è previsto che un utente possa utilizzare un numero
limitato di istanze di \textit{inotify}; il valore di default del limite è di
128, ma questo valore può essere cambiato con \func{sysctl} o usando il file
\sysctlfiled{fs/inotify/max\_user\_instances}.

Dato che questo file descriptor non è associato a nessun file o directory
reale, l'inconveniente di non poter smontare un filesystem i cui file sono
tenuti sotto osservazione viene completamente eliminato; anzi, una delle
capacità dell'interfaccia di \textit{inotify} è proprio quella di notificare
il fatto che il filesystem su cui si trova il file o la directory osservata è
stato smontato.

Inoltre trattandosi di un file descriptor a tutti gli effetti, esso potrà
essere utilizzato come argomento per le funzioni \func{select} e \func{poll} e
con l'interfaccia di \textit{epoll}, ed a partire dal kernel 2.6.25 è stato
introdotto anche il supporto per il \texttt{signal-driven I/O}.  Siccome gli
eventi vengono notificati come dati disponibili in lettura, dette funzioni
ritorneranno tutte le volte che si avrà un evento di notifica.

Così, invece di dover utilizzare i segnali, considerati una pessima scelta dal
punto di vista dell'interfaccia utente, si potrà gestire l'osservazione degli
eventi con una qualunque delle modalità di \textit{I/O multiplexing}
illustrate in sez.~\ref{sec:file_multiplexing}. Qualora si voglia cessare
l'osservazione, sarà sufficiente chiudere il file descriptor e tutte le
risorse allocate saranno automaticamente rilasciate. Infine l'interfaccia di
\textit{inotify} consente di mettere sotto osservazione, oltre che una
directory, anche singoli file.

Una volta creata la coda di notifica si devono definire gli eventi da tenere
sotto osservazione; questo viene fatto attraverso una \textsl{lista di
  osservazione} (o \textit{watch list}) che è associata alla coda. Per gestire
la lista di osservazione l'interfaccia fornisce due funzioni di sistema, la
prima di queste è \funcd{inotify\_add\_watch}, il cui prototipo è:

\begin{funcproto}{
\fhead{sys/inotify.h}
\fdecl{int inotify\_add\_watch(int fd, const char *pathname, uint32\_t mask)}
\fdesc{Aggiunge un evento di osservazione a una lista di osservazione.} 
}

{La funzione ritorna un valore positivo in caso di successo, o $-1$ per un
  errore, nel qual caso \var{errno} assumerà uno dei valori:
  \begin{errlist}
  \item[\errcode{EACCES}] non si ha accesso in lettura al file indicato.
  \item[\errcode{EINVAL}] \param{mask} non contiene eventi legali o \param{fd}
    non è un file descriptor di \textit{inotify}.
  \item[\errcode{ENOSPC}] si è raggiunto il numero massimo di voci di
    osservazione o il kernel non ha potuto allocare una risorsa necessaria.
  \end{errlist}
  ed inoltre \errval{EFAULT}, \errval{ENOMEM} e \errval{EBADF} nel loro
  significato generico.}
\end{funcproto}

La funzione consente di creare un ``\textsl{osservatore}'' (il cosiddetto
``\textit{watch}'') nella lista di osservazione di una coda di notifica, che
deve essere indicata specificando il file descriptor ad essa associato
nell'argomento \param{fd}, che ovviamente dovrà essere un file descriptor
creato con \func{inotify\_init}.  Il file o la directory da porre sotto
osservazione vengono invece indicati per nome, da passare
nell'argomento \param{pathname}.  Infine il terzo argomento, \param{mask},
indica che tipo di eventi devono essere tenuti sotto osservazione e le
modalità della stessa.  L'operazione può essere ripetuta per tutti i file e le
directory che si vogliono tenere sotto osservazione,\footnote{anche in questo
  caso c'è un limite massimo che di default è pari a 8192, ed anche questo
  valore può essere cambiato con \func{sysctl} o usando il file
  \sysctlfiled{fs/inotify/max\_user\_watches}.} e si utilizzerà sempre un solo
file descriptor.

Il tipo di evento che si vuole osservare deve essere specificato
nell'argomento \param{mask} come maschera binaria, combinando i valori delle
costanti riportate in tab.~\ref{tab:inotify_event_watch} che identificano i
singoli bit della maschera ed il relativo significato. In essa si sono marcati
con un ``$\bullet$'' gli eventi che, quando specificati per una directory,
vengono osservati anche su tutti i file che essa contiene.  Nella seconda
parte della tabella si sono poi indicate alcune combinazioni predefinite dei
flag della prima parte.

\begin{table}[htb]
  \centering
  \footnotesize
  \begin{tabular}[c]{|l|c|p{8cm}|}
    \hline
    \textbf{Valore}  & & \textbf{Significato} \\
    \hline
    \hline
    \constd{IN\_ACCESS}        &$\bullet$& C'è stato accesso al file in
                                           lettura.\\  
    \constd{IN\_ATTRIB}        &$\bullet$& Ci sono stati cambiamenti sui dati
                                           dell'\textit{inode}
                                           (o sugli attributi estesi, vedi
                                           sez.~\ref{sec:file_xattr}).\\ 
    \constd{IN\_CLOSE\_WRITE}  &$\bullet$& È stato chiuso un file aperto in
                                           scrittura.\\  
    \constd{IN\_CLOSE\_NOWRITE}&$\bullet$& È stato chiuso un file aperto in
                                           sola lettura.\\
    \constd{IN\_CREATE}        &$\bullet$& È stato creato un file o una
                                           directory in una directory sotto
                                           osservazione.\\  
    \constd{IN\_DELETE}        &$\bullet$& È stato cancellato un file o una
                                           directory in una directory sotto
                                           osservazione.\\ 
    \constd{IN\_DELETE\_SELF}  & --      & È stato cancellato il file (o la
                                          directory) sotto osservazione.\\ 
    \constd{IN\_MODIFY}        &$\bullet$& È stato modificato il file.\\ 
    \constd{IN\_MOVE\_SELF}    &         & È stato rinominato il file (o la
                                           directory) sotto osservazione.\\ 
    \constd{IN\_MOVED\_FROM}   &$\bullet$& Un file è stato spostato fuori dalla
                                           directory sotto osservazione.\\ 
    \constd{IN\_MOVED\_TO}     &$\bullet$& Un file è stato spostato nella
                                           directory sotto osservazione.\\ 
    \constd{IN\_OPEN}          &$\bullet$& Un file è stato aperto.\\ 
    \hline    
    \constd{IN\_CLOSE}         &         & Combinazione di
                                           \const{IN\_CLOSE\_WRITE} e
                                           \const{IN\_CLOSE\_NOWRITE}.\\  
    \constd{IN\_MOVE}          &         & Combinazione di
                                           \const{IN\_MOVED\_FROM} e
                                           \const{IN\_MOVED\_TO}.\\
    \constd{IN\_ALL\_EVENTS}   &         & Combinazione di tutti i flag
                                           possibili.\\
    \hline    
  \end{tabular}
  \caption{Le costanti che identificano i bit della maschera binaria
    dell'argomento \param{mask} di \func{inotify\_add\_watch} che indicano il
    tipo di evento da tenere sotto osservazione.} 
  \label{tab:inotify_event_watch}
\end{table}

Oltre ai flag di tab.~\ref{tab:inotify_event_watch}, che indicano il tipo di
evento da osservare e che vengono utilizzati anche in uscita per indicare il
tipo di evento avvenuto, \func{inotify\_add\_watch} supporta ulteriori
flag,\footnote{i flag \const{IN\_DONT\_FOLLOW}, \const{IN\_MASK\_ADD} e
  \const{IN\_ONLYDIR} sono stati introdotti a partire dalle glibc 2.5, se si
  usa la versione 2.4 è necessario definirli a mano.}  riportati in
tab.~\ref{tab:inotify_add_watch_flag}, che indicano le modalità di
osservazione (da passare sempre nell'argomento \param{mask}) e che al
contrario dei precedenti non vengono mai impostati nei risultati in uscita.

\begin{table}[htb]
  \centering
  \footnotesize
  \begin{tabular}[c]{|l|p{8cm}|}
    \hline
    \textbf{Valore}  & \textbf{Significato} \\
    \hline
    \hline
    \constd{IN\_DONT\_FOLLOW}& Non dereferenzia \param{pathname} se questo è un
                               link simbolico.\\
    \constd{IN\_MASK\_ADD}   & Aggiunge a quelli già impostati i flag indicati
                               nell'argomento \param{mask}, invece di
                               sovrascriverli.\\
    \constd{IN\_ONESHOT}     & Esegue l'osservazione su \param{pathname} per
                               una sola volta, rimuovendolo poi dalla
                               \textit{watch list}.\\ 
    \constd{IN\_ONLYDIR}     & Se \param{pathname} è una directory riporta
                               soltanto gli eventi ad essa relativi e non
                               quelli per i file che contiene.\\ 
    \hline    
  \end{tabular}
  \caption{Le costanti che identificano i bit della maschera binaria
    dell'argomento \param{mask} di \func{inotify\_add\_watch} che indicano le
    modalità di osservazione.} 
  \label{tab:inotify_add_watch_flag}
\end{table}

Se non esiste nessun \textit{watch} per il file o la directory specificata
questo verrà creato per gli eventi specificati dall'argomento \param{mask},
altrimenti la funzione sovrascriverà le impostazioni precedenti, a meno che
non si sia usato il flag \const{IN\_MASK\_ADD}, nel qual caso gli eventi
specificati saranno aggiunti a quelli già presenti.

Come accennato quando si tiene sotto osservazione una directory vengono
restituite le informazioni sia riguardo alla directory stessa che ai file che
essa contiene; questo comportamento può essere disabilitato utilizzando il
flag \const{IN\_ONLYDIR}, che richiede di riportare soltanto gli eventi
relativi alla directory stessa. Si tenga presente inoltre che quando si
osserva una directory vengono riportati solo gli eventi sui file che essa
contiene direttamente, non quelli relativi a file contenuti in eventuali
sottodirectory; se si vogliono osservare anche questi sarà necessario creare
ulteriori \textit{watch} per ciascuna sottodirectory.

Infine usando il flag \const{IN\_ONESHOT} è possibile richiedere una notifica
singola;\footnote{questa funzionalità però è disponibile soltanto a partire dal
  kernel 2.6.16.} una volta verificatosi uno qualunque fra gli eventi
richiesti con \func{inotify\_add\_watch} l'\textsl{osservatore} verrà
automaticamente rimosso dalla lista di osservazione e nessun ulteriore evento
sarà più notificato.

In caso di successo \func{inotify\_add\_watch} ritorna un intero positivo,
detto \textit{watch descriptor}, che identifica univocamente un
\textsl{osservatore} su una coda di notifica; esso viene usato per farvi
riferimento sia riguardo i risultati restituiti da \textit{inotify}, che per
la eventuale rimozione dello stesso. 

La seconda funzione di sistema per la gestione delle code di notifica, che
permette di rimuovere un \textsl{osservatore}, è \funcd{inotify\_rm\_watch},
ed il suo prototipo è:

\begin{funcproto}{
\fhead{sys/inotify.h}
\fdecl{int inotify\_rm\_watch(int fd, uint32\_t wd)}
\fdesc{Rimuove un \textsl{osservatore} da una coda di notifica.} 
}

{La funzione ritorna $0$ in caso di successo e $-1$ per un errore, nel qual
  caso \var{errno} assumerà uno dei valori: 
  \begin{errlist}
  \item[\errcode{EBADF}] non si è specificato in \param{fd} un file descriptor
    valido.
  \item[\errcode{EINVAL}] il valore di \param{wd} non è corretto, o \param{fd}
    non è associato ad una coda di notifica.
  \end{errlist}
}
\end{funcproto}

La funzione rimuove dalla coda di notifica identificata dall'argomento
\param{fd} l'osservatore identificato dal \textit{watch descriptor}
\param{wd}; ovviamente deve essere usato per questo argomento un valore
ritornato da \func{inotify\_add\_watch}, altrimenti si avrà un errore di
\errval{EINVAL}. In caso di successo della rimozione, contemporaneamente alla
cancellazione dell'osservatore, sulla coda di notifica verrà generato un
evento di tipo \const{IN\_IGNORED} (vedi
tab.~\ref{tab:inotify_read_event_flag}). Si tenga presente che se un file
viene cancellato o un filesystem viene smontato i relativi osservatori vengono
rimossi automaticamente e non è necessario utilizzare
\func{inotify\_rm\_watch}.

Come accennato l'interfaccia di \textit{inotify} prevede che gli eventi siano
notificati come dati presenti in lettura sul file descriptor associato alla
coda di notifica. Una applicazione pertanto dovrà leggere i dati da detto file
con una \func{read}, che ritornerà sul buffer i dati presenti nella forma di
una o più strutture di tipo \struct{inotify\_event} (la cui definizione è
riportata in fig.~\ref{fig:inotify_event}). Qualora non siano presenti dati la
\func{read} si bloccherà (a meno di non aver impostato il file descriptor in
modalità non bloccante) fino all'arrivo di almeno un evento.

\begin{figure}[!htb]
  \footnotesize \centering
  \begin{minipage}[c]{0.90\textwidth}
    \includestruct{listati/inotify_event.h}
  \end{minipage} 
  \normalsize 
  \caption{La struttura \structd{inotify\_event} usata dall'interfaccia di
    \textit{inotify} per riportare gli eventi.}
  \label{fig:inotify_event}
\end{figure}

Una ulteriore caratteristica dell'interfaccia di \textit{inotify} è che essa
permette di ottenere con \func{ioctl}, come per i file descriptor associati ai
socket (si veda sez.~\ref{sec:sock_ioctl_IP}), il numero di byte disponibili
in lettura sul file descriptor, utilizzando su di esso l'operazione
\const{FIONREAD}.\footnote{questa è una delle operazioni speciali per i file
  (vedi sez.~\ref{sec:file_fcntl_ioctl}), che è disponibile solo per i socket
  e per i file descriptor creati con \func{inotify\_init}.} Si può così
utilizzare questa operazione, oltre che per predisporre una operazione di
lettura con un buffer di dimensioni adeguate, anche per ottenere rapidamente
il numero di file che sono cambiati.

Una volta effettuata la lettura con \func{read} a ciascun evento sarà
associata una struttura \struct{inotify\_event} contenente i rispettivi dati.
Per identificare a quale file o directory l'evento corrisponde viene
restituito nel campo \var{wd} il \textit{watch descriptor} con cui il relativo
osservatore è stato registrato. Il campo \var{mask} contiene invece una
maschera di bit che identifica il tipo di evento verificatosi; in essa
compariranno sia i bit elencati nella prima parte di
tab.~\ref{tab:inotify_event_watch}, che gli eventuali valori aggiuntivi di
tab.~\ref{tab:inotify_read_event_flag} (questi compaiono solo nel campo
\var{mask} di \struct{inotify\_event}, e non sono utilizzabili in fase di
registrazione dell'osservatore).

\begin{table}[htb]
  \centering
  \footnotesize
  \begin{tabular}[c]{|l|p{10cm}|}
    \hline
    \textbf{Valore}  & \textbf{Significato} \\
    \hline
    \hline
    \constd{IN\_IGNORED}    & L'osservatore è stato rimosso, sia in maniera 
                              esplicita con l'uso di \func{inotify\_rm\_watch}, 
                              che in maniera implicita per la rimozione 
                              dell'oggetto osservato o per lo smontaggio del
                              filesystem su cui questo si trova.\\
    \constd{IN\_ISDIR}      & L'evento avvenuto fa riferimento ad una directory
                              (consente così di distinguere, quando si pone
                              sotto osservazione una directory, fra gli eventi
                              relativi ad essa e quelli relativi ai file che
                              essa contiene).\\
    \constd{IN\_Q\_OVERFLOW}& Si sono eccedute le dimensioni della coda degli
                              eventi (\textit{overflow} della coda); in questo
                              caso il valore di \var{wd} è $-1$.\footnotemark\\
    \constd{IN\_UNMOUNT}    & Il filesystem contenente l'oggetto posto sotto
                              osservazione è stato smontato.\\
    \hline    
  \end{tabular}
  \caption{Le costanti che identificano i bit aggiuntivi usati nella maschera
    binaria del campo \var{mask} di \struct{inotify\_event}.} 
  \label{tab:inotify_read_event_flag}
\end{table}

\footnotetext{la coda di notifica ha una dimensione massima che viene
  controllata dal parametro di sistema
  \sysctlfiled{fs/inotify/max\_queued\_events}, che indica il numero massimo di
  eventi che possono essere mantenuti sulla stessa; quando detto valore viene
  ecceduto gli ulteriori eventi vengono scartati, ma viene comunque generato
  un evento di tipo \const{IN\_Q\_OVERFLOW}.}

Il campo \var{cookie} contiene invece un intero univoco che permette di
identificare eventi correlati (per i quali avrà lo stesso valore), al momento
viene utilizzato soltanto per rilevare lo spostamento di un file, consentendo
così all'applicazione di collegare la corrispondente coppia di eventi
\const{IN\_MOVED\_TO} e \const{IN\_MOVED\_FROM}.

Infine due campi \var{name} e \var{len} sono utilizzati soltanto quando
l'evento è relativo ad un file presente in una directory posta sotto
osservazione, in tal caso essi contengono rispettivamente il nome del file
(come \textit{pathname} relativo alla directory osservata) e la relativa
dimensione in byte. Il campo \var{name} viene sempre restituito come stringa
terminata da NUL, con uno o più zeri di terminazione, a seconda di eventuali
necessità di allineamento del risultato, ed il valore di \var{len} corrisponde
al totale della dimensione di \var{name}, zeri aggiuntivi compresi. La stringa
con il nome del file viene restituita nella lettura subito dopo la struttura
\struct{inotify\_event}; questo significa che le dimensioni di ciascun evento
di \textit{inotify} saranno pari a \code{sizeof(\struct{inotify\_event}) +
  len}.

Vediamo allora un esempio dell'uso dell'interfaccia di \textit{inotify} con un
semplice programma che permette di mettere sotto osservazione uno o più file e
directory. Il programma si chiama \texttt{inotify\_monitor.c} ed il codice
completo è disponibile coi sorgenti allegati alla guida, il corpo principale
del programma, che non contiene la sezione di gestione delle opzioni e le
funzioni di ausilio è riportato in fig.~\ref{fig:inotify_monitor_example}.

\begin{figure}[!htbp]
  \footnotesize \centering
  \begin{minipage}[c]{\codesamplewidth}
    \includecodesample{listati/inotify_monitor.c}
  \end{minipage}
  \normalsize
  \caption{Esempio di codice che usa l'interfaccia di \textit{inotify}.}
  \label{fig:inotify_monitor_example}
\end{figure}

Una volta completata la scansione delle opzioni il corpo del programma inizia
controllando (\texttt{\small 11-15}) che sia rimasto almeno un argomento che
indichi quale file o directory mettere sotto osservazione (e qualora questo
non avvenga esce stampando la pagina di aiuto); dopo di che passa
(\texttt{\small 16-20}) all'inizializzazione di \textit{inotify} ottenendo con
\func{inotify\_init} il relativo file descriptor (o si esce in caso di
errore).

Il passo successivo è aggiungere (\texttt{\small 21-30}) alla coda di
notifica gli opportuni osservatori per ciascuno dei file o directory indicati
all'invocazione del comando; questo viene fatto eseguendo un ciclo
(\texttt{\small 22-29}) fintanto che la variabile \var{i}, inizializzata a
zero (\texttt{\small 21}) all'inizio del ciclo, è minore del numero totale di
argomenti rimasti. All'interno del ciclo si invoca (\texttt{\small 23})
\func{inotify\_add\_watch} per ciascuno degli argomenti, usando la maschera
degli eventi data dalla variabile \var{mask} (il cui valore viene impostato
nella scansione delle opzioni), in caso di errore si esce dal programma
altrimenti si incrementa l'indice (\texttt{\small 29}).

Completa l'inizializzazione di \textit{inotify} inizia il ciclo principale
(\texttt{\small 32-56}) del programma, nel quale si resta in attesa degli
eventi che si intendono osservare. Questo viene fatto eseguendo all'inizio del
ciclo (\texttt{\small 33}) una \func{read} che si bloccherà fintanto che non
si saranno verificati eventi.

Dato che l'interfaccia di \textit{inotify} può riportare anche più eventi in
una sola lettura, si è avuto cura di passare alla \func{read} un buffer di
dimensioni adeguate, inizializzato in (\texttt{\small 7}) ad un valore di
approssimativamente 512 eventi (si ricordi che la quantità di dati restituita
da \textit{inotify} è variabile a causa della diversa lunghezza del nome del
file restituito insieme a \struct{inotify\_event}). In caso di errore di
lettura (\texttt{\small 35-40}) il programma esce con un messaggio di errore
(\texttt{\small 37-39}), a meno che non si tratti di una interruzione della
\textit{system call}, nel qual caso (\texttt{\small 36}) si ripete la lettura.

Se la lettura è andata a buon fine invece si esegue un ciclo (\texttt{\small
  43-52}) per leggere tutti gli eventi restituiti, al solito si inizializza
l'indice \var{i} a zero (\texttt{\small 42}) e si ripetono le operazioni
(\texttt{\small 43}) fintanto che esso non supera il numero di byte restituiti
in lettura. Per ciascun evento all'interno del ciclo si assegna alla variabile
\var{event} (si noti come si sia eseguito un opportuno \textit{casting} del
puntatore) l'indirizzo nel buffer della corrispondente struttura
\struct{inotify\_event} (\texttt{\small 44}), e poi si stampano il numero di
\textit{watch descriptor} (\texttt{\small 45}) ed il file a cui questo fa
riferimento (\texttt{\small 46}), ricavato dagli argomenti passati a riga di
comando sfruttando il fatto che i \textit{watch descriptor} vengono assegnati
in ordine progressivo crescente a partire da 1.

Qualora sia presente il riferimento ad un nome di file associato all'evento lo
si stampa (\texttt{\small 47-49}); si noti come in questo caso si sia
controllato il valore del campo \var{event->len} e non il fatto che
\var{event->name} riporti o meno un puntatore nullo. L'interfaccia infatti,
qualora il nome non sia presente, non tocca il campo \var{event->name}, che
si troverà pertanto a contenere quello che era precedentemente presente nella
rispettiva locazione di memoria, nel caso più comune il puntatore al nome di
un file osservato in precedenza.

Si utilizza poi (\texttt{\small 50}) la funzione \code{printevent}, che
interpreta il valore del campo \var{event->mask}, per stampare il tipo di
eventi accaduti.\footnote{per il relativo codice, che non riportiamo in quanto
  non essenziale alla comprensione dell'esempio, si possono utilizzare
  direttamente i sorgenti allegati alla guida.} Infine (\texttt{\small 51}) si
provvede ad aggiornare l'indice \var{i} per farlo puntare all'evento
successivo.

Se adesso usiamo il programma per mettere sotto osservazione una directory, e
da un altro terminale eseguiamo il comando \texttt{ls} otterremo qualcosa del
tipo di:
\begin{Console}
piccardi@gethen:~/gapil/sources$ \textbf{./inotify_monitor -a /home/piccardi/gapil/}
Watch descriptor 1
Observed event on /home/piccardi/gapil/
IN_OPEN, 
Watch descriptor 1
Observed event on /home/piccardi/gapil/
IN_CLOSE_NOWRITE, 
\end{Console}
%$

I lettori più accorti si saranno resi conto che nel ciclo di lettura degli
eventi appena illustrato non viene trattato il caso particolare in cui la
funzione \func{read} restituisce in \var{nread} un valore nullo. Lo si è fatto
perché con \textit{inotify} il ritorno di una \func{read} con un valore nullo
avviene soltanto, come forma di avviso, quando si sia eseguita la funzione
specificando un buffer di dimensione insufficiente a contenere anche un solo
evento. Nel nostro caso le dimensioni erano senz'altro sufficienti, per cui
tale evenienza non si verificherà mai.

Ci si potrà però chiedere cosa succede se il buffer è sufficiente per un
evento, ma non per tutti gli eventi verificatisi. Come si potrà notare nel
codice illustrato in precedenza non si è presa nessuna precauzione per
verificare che non ci fossero stati troncamenti dei dati. Anche in questo caso
il comportamento scelto è corretto, perché l'interfaccia di \textit{inotify}
garantisce automaticamente, anche quando ne sono presenti in numero maggiore,
di restituire soltanto il numero di eventi che possono rientrare completamente
nelle dimensioni del buffer specificato.\footnote{si avrà cioè, facendo
  riferimento sempre al codice di fig.~\ref{fig:inotify_monitor_example}, che
  \var{read} sarà in genere minore delle dimensioni di \var{buffer} ed uguale
  soltanto qualora gli eventi corrispondano esattamente alle dimensioni di
  quest'ultimo.} Se gli eventi sono di più saranno restituiti solo quelli che
entrano interamente nel buffer e gli altri saranno restituiti alla successiva
chiamata di \func{read}.

Infine un'ultima caratteristica dell'interfaccia di \textit{inotify} è che gli
eventi restituiti nella lettura formano una sequenza ordinata, è cioè
garantito che se si esegue uno spostamento di un file gli eventi vengano
generati nella sequenza corretta. L'interfaccia garantisce anche che se si
verificano più eventi consecutivi identici (vale a dire con gli stessi valori
dei campi \var{wd}, \var{mask}, \var{cookie}, e \var{name}) questi vengono
raggruppati in un solo evento.

\itindend{inotify}

% TODO trattare fanotify, vedi http://lwn.net/Articles/339399/ e 
% http://lwn.net/Articles/343346/ (incluso nel 2.6.36)
% fanotify_mark() ha FAN_MARK_FILESYSTEM dal 4.20
% fanotify() ha FAN_OPEN_EXEC dal 4.21/5.0


\subsection{L'interfaccia POSIX per l'I/O asincrono}
\label{sec:file_asyncronous_io}

Una modalità alternativa all'uso dell'\textit{I/O multiplexing} per gestione
dell'I/O simultaneo su molti file è costituita dal cosiddetto \textsl{I/O
  asincrono} o ``AIO''. Il concetto base dell'\textsl{I/O asincrono} è che le
funzioni di I/O non attendono il completamento delle operazioni prima di
ritornare, così che il processo non viene bloccato.  In questo modo diventa ad
esempio possibile effettuare una richiesta preventiva di dati, in modo da
poter effettuare in contemporanea le operazioni di calcolo e quelle di I/O.

Benché la modalità di apertura asincrona di un file vista in
sez.~\ref{sec:signal_driven_io} possa risultare utile in varie occasioni (in
particolar modo con i socket e gli altri file per i quali le funzioni di I/O
sono \textit{system call} lente), essa è comunque limitata alla notifica della
disponibilità del file descriptor per le operazioni di I/O, e non ad uno
svolgimento asincrono delle medesime.  Lo standard POSIX.1b definisce una
interfaccia apposita per l'I/O asincrono vero e proprio,\footnote{questa è
  stata ulteriormente perfezionata nelle successive versioni POSIX.1-2001 e
  POSIX.1-2008.} che prevede un insieme di funzioni dedicate per la lettura e
la scrittura dei file, completamente separate rispetto a quelle usate
normalmente.

In generale questa interfaccia è completamente astratta e può essere
implementata sia direttamente nel kernel che in \textit{user space} attraverso
l'uso di \textit{thread}. Per le versioni del kernel meno recenti esiste una
implementazione di questa interfaccia fornita completamente dalla \acr{glibc}
a partire dalla versione 2.1, che è realizzata completamente in \textit{user
  space}, ed è accessibile linkando i programmi con la libreria
\file{librt}. A partire dalla versione 2.5.32 è stato introdotto nel kernel
una nuova infrastruttura per l'I/O asincrono, ma ancora il supporto è parziale
ed insufficiente ad implementare tutto l'AIO POSIX.

Lo standard POSIX prevede che tutte le operazioni di I/O asincrono siano
controllate attraverso l'uso di una apposita struttura \struct{aiocb} (il cui
nome sta per \textit{asyncronous I/O control block}), che viene passata come
argomento a tutte le funzioni dell'interfaccia. La sua definizione, come
effettuata in \headfiled{aio.h}, è riportata in
fig.~\ref{fig:file_aiocb}. Nello steso file è definita la macro
\macrod{\_POSIX\_ASYNCHRONOUS\_IO}, che dichiara la disponibilità
dell'interfaccia per l'I/O asincrono.

\begin{figure}[!htb]
  \footnotesize \centering
  \begin{minipage}[c]{0.90\textwidth}
    \includestruct{listati/aiocb.h}
  \end{minipage}
  \normalsize 
  \caption{La struttura \structd{aiocb}, usata per il controllo dell'I/O
    asincrono.}
  \label{fig:file_aiocb}
\end{figure}

Le operazioni di I/O asincrono possono essere effettuate solo su un file già
aperto; il file deve inoltre supportare la funzione \func{lseek}, pertanto
terminali e \textit{pipe} sono esclusi. Non c'è limite al numero di operazioni
contemporanee effettuabili su un singolo file.  Ogni operazione deve
inizializzare opportunamente un \textit{control block}.  Il file descriptor su
cui operare deve essere specificato tramite il campo \var{aio\_fildes}; dato
che più operazioni possono essere eseguita in maniera asincrona, il concetto
di posizione corrente sul file viene a mancare; pertanto si deve sempre
specificare nel campo \var{aio\_offset} la posizione sul file da cui i dati
saranno letti o scritti.  Nel campo \var{aio\_buf} deve essere specificato
l'indirizzo del buffer usato per l'I/O, ed in \var{aio\_nbytes} la lunghezza
del blocco di dati da trasferire.

Il campo \var{aio\_reqprio} permette di impostare la priorità delle operazioni
di I/O, in generale perché ciò sia possibile occorre che la piattaforma
supporti questa caratteristica, questo viene indicato dal fatto che le macro
\macrod{\_POSIX\_PRIORITIZED\_IO}, e \macrod{\_POSIX\_PRIORITY\_SCHEDULING}
sono definite. La priorità viene impostata a partire da quella del processo
chiamante (vedi sez.~\ref{sec:proc_priority}), cui viene sottratto il valore
di questo campo.  Il campo \var{aio\_lio\_opcode} è usato solo dalla funzione
\func{lio\_listio}, che, come vedremo, permette di eseguire con una sola
chiamata una serie di operazioni, usando un vettore di \textit{control
  block}. Tramite questo campo si specifica quale è la natura di ciascuna di
esse.

Infine il campo \var{aio\_sigevent} è una struttura di tipo \struct{sigevent}
(illustrata in fig.~\ref{fig:struct_sigevent}) che serve a specificare il modo
in cui si vuole che venga effettuata la notifica del completamento delle
operazioni richieste; per la trattazione delle modalità di utilizzo della
stessa si veda quanto già visto in proposito in sez.~\ref{sec:sig_timer_adv}.

Le due funzioni base dell'interfaccia per l'I/O asincrono sono
\funcd{aio\_read} ed \funcd{aio\_write}.  Esse permettono di richiedere una
lettura od una scrittura asincrona di dati usando la struttura \struct{aiocb}
appena descritta; i rispettivi prototipi sono:

\begin{funcproto}{
\fhead{aio.h}
\fdecl{int aio\_read(struct aiocb *aiocbp)}
\fdesc{Richiede una lettura asincrona.} 
\fdecl{int aio\_write(struct aiocb *aiocbp)}
\fdesc{Richiede una scrittura asincrona.} 
}

{Le funzioni ritornano $0$ in caso di successo e $-1$ per un errore, nel qual
  caso \var{errno} assumerà uno dei valori: 
  \begin{errlist}
  \item[\errcode{EAGAIN}] la coda delle richieste è momentaneamente piena.
  \item[\errcode{EBADF}] si è specificato un file descriptor sbagliato.
  \item[\errcode{EINVAL}] si è specificato un valore non valido per i campi
    \var{aio\_offset} o \var{aio\_reqprio} di \param{aiocbp}.
  \item[\errcode{ENOSYS}] la funzione non è implementata.
  \end{errlist}
}
\end{funcproto}


Entrambe le funzioni ritornano immediatamente dopo aver messo in coda la
richiesta, o in caso di errore. Non è detto che gli errori \errcode{EBADF} ed
\errcode{EINVAL} siano rilevati immediatamente al momento della chiamata,
potrebbero anche emergere nelle fasi successive delle operazioni. Lettura e
scrittura avvengono alla posizione indicata da \var{aio\_offset}, a meno che
il file non sia stato aperto in \textit{append mode} (vedi
sez.~\ref{sec:file_open_close}), nel qual caso le scritture vengono effettuate
comunque alla fine del file, nell'ordine delle chiamate a \func{aio\_write}.

Si tenga inoltre presente che deallocare la memoria indirizzata da
\param{aiocbp} o modificarne i valori prima della conclusione di una
operazione può dar luogo a risultati impredicibili, perché l'accesso ai vari
campi per eseguire l'operazione può avvenire in un momento qualsiasi dopo la
richiesta. Questo comporta che non si devono usare per \param{aiocbp}
variabili automatiche e che non si deve riutilizzare la stessa struttura per
un'altra operazione fintanto che la precedente non sia stata ultimata. In
generale per ogni operazione si deve utilizzare una diversa struttura
\struct{aiocb}.

Dato che si opera in modalità asincrona, il successo di \func{aio\_read} o
\func{aio\_write} non implica che le operazioni siano state effettivamente
eseguite in maniera corretta; per verificarne l'esito l'interfaccia prevede
altre due funzioni, che permettono di controllare lo stato di esecuzione. La
prima è \funcd{aio\_error}, che serve a determinare un eventuale stato di
errore; il suo prototipo è:

\begin{funcproto}{
\fhead{aio.h}
\fdecl{int aio\_error(const struct aiocb *aiocbp)} 
\fdesc{Determina lo stato di errore di una operazione di I/O asincrono.} 
}

{La funzione ritorna $0$ se le operazioni si sono concluse con successo,
  altrimenti restituisce \errval{EINPROGRESS} se non sono concluse,
  \errcode{ECANCELED} se sono state cancellate o il relativo codice di errore
  se sono fallite.}
\end{funcproto}

Se l'operazione non si è ancora completata viene sempre restituito l'errore di
\errcode{EINPROGRESS}, mentre se è stata cancellata ritorna
\errcode{ECANCELED}. La funzione ritorna zero quando l'operazione si è
conclusa con successo, altrimenti restituisce il codice dell'errore
verificatosi, ed esegue la corrispondente impostazione di \var{errno}. Il
codice può essere sia \errcode{EINVAL} ed \errcode{EBADF}, dovuti ad un valore
errato per \param{aiocbp}, che uno degli errori possibili durante l'esecuzione
dell'operazione di I/O richiesta, nel qual caso saranno restituiti, a seconda
del caso, i codici di errore delle \textit{system call} \func{read},
\func{write}, \func{fsync} e \func{fdatasync}.

Una volta che si sia certi che le operazioni siano state concluse (cioè dopo
che una chiamata ad \func{aio\_error} non ha restituito
\errcode{EINPROGRESS}), si potrà usare la funzione \funcd{aio\_return}, che
permette di verificare il completamento delle operazioni di I/O asincrono; il
suo prototipo è:

\begin{funcproto}{
\fhead{aio.h}
\fdecl{ssize\_t aio\_return(const struct aiocb *aiocbp)}
\fdesc{Ottiene lo stato dei risultati di una operazione di I/O asincrono.} 
}

{La funzione ritorna lo stato di uscita dell'operazione eseguita (il valore
  che avrebbero restituito le equivalenti funzioni eseguite in maniera
  sincrona).}
\end{funcproto}

La funzione recupera il valore dello stato di ritorno delle operazioni di I/O
associate a \param{aiocbp} e deve essere chiamata una sola volta per ciascuna
operazione asincrona, essa infatti fa sì che il sistema rilasci le risorse ad
essa associate. É per questo motivo che occorre chiamare la funzione solo dopo
che l'operazione cui \param{aiocbp} fa riferimento si è completata
verificandolo con \func{aio\_error}, ed usarla una sola volta. Una chiamata
precedente il completamento delle operazioni darebbe risultati indeterminati,
così come chiamarla più di una volta.

La funzione restituisce il valore di ritorno relativo all'operazione eseguita,
così come ricavato dalla sottostante \textit{system call} (il numero di byte
letti, scritti o il valore di ritorno di \func{fsync} o \func{fdatasync}).  É
importante chiamare sempre questa funzione, altrimenti le risorse disponibili
per le operazioni di I/O asincrono non verrebbero liberate, rischiando di
arrivare ad un loro esaurimento.

Oltre alle operazioni di lettura e scrittura l'interfaccia POSIX.1b mette a
disposizione un'altra operazione, quella di sincronizzazione dell'I/O,
compiuta dalla funzione \funcd{aio\_fsync}, che ha lo stesso effetto della
analoga \func{fsync}, ma viene eseguita in maniera asincrona; il suo prototipo
è:

\begin{funcproto}{
\fhead{aio.h}
\fdecl{int aio\_fsync(int op, struct aiocb *aiocbp)} 
\fdesc{Richiede la sincronizzazione dei dati su disco.} 
}

{La funzione ritorna $0$ in caso di successo e $-1$ per un errore, nel qual
  caso \var{errno} assumerà gli stessi valori visti \func{aio\_read} con lo
  stesso significato.
}
\end{funcproto}

La funzione richiede la sincronizzazione dei dati delle operazioni di I/O
relative al file descriptor indicato in \texttt{aiocbp->aio\_fildes},
ritornando immediatamente. Si tenga presente che la funzione mette
semplicemente in coda la richiesta, l'esecuzione effettiva della
sincronizzazione dovrà essere verificata con \func{aio\_error} e
\func{aio\_return} come per le operazioni di lettura e
scrittura. L'argomento \param{op} permette di indicare la modalità di
esecuzione, se si specifica il valore \const{O\_DSYNC} le operazioni saranno
completate con una chiamata a \func{fdatasync}, se si specifica
\const{O\_SYNC} con una chiamata a \func{fsync} (per i dettagli vedi
sez.~\ref{sec:file_sync}).

Il successo della chiamata assicura la richiesta di sincronizzazione dei dati
relativi operazioni di I/O asincrono richieste fino a quel momento, niente è
garantito riguardo la sincronizzazione dei dati relativi ad eventuali
operazioni richieste successivamente. Se si è specificato un meccanismo di
notifica questo sarà innescato una volta che le operazioni di sincronizzazione
dei dati saranno completate (\texttt{aio\_sigevent} è l'unico altro campo
di \param{aiocbp} che viene usato.

In alcuni casi può essere necessario interrompere le operazioni di I/O (in
genere quando viene richiesta un'uscita immediata dal programma), per questo
lo standard POSIX.1b prevede una funzione apposita, \funcd{aio\_cancel}, che
permette di cancellare una operazione richiesta in precedenza; il suo
prototipo è:

\begin{funcproto}{
\fhead{aio.h}
\fdecl{int aio\_cancel(int fd, struct aiocb *aiocbp)}
\fdesc{Richiede la cancellazione delle operazioni di I/O asincrono.} 
}

{La funzione ritorna un intero positivo che indica il risultato
  dell'operazione in caso di successo e $-1$ per un errore, nel qual caso
  \var{errno} assumerà uno dei valori:
  \begin{errlist}
  \item[\errcode{EBADF}] \param{fd} non è un file descriptor valido.
  \item[\errcode{ENOSYS}] la funzione non è implementata.
  \end{errlist}
}
\end{funcproto}

La funzione permette di cancellare una operazione specifica sul file
\param{fd}, idicata con \param{aiocbp}, o tutte le operazioni pendenti,
specificando \val{NULL} come valore di \param{aiocbp}. Quando una operazione
viene cancellata una successiva chiamata ad \func{aio\_error} riporterà
\errcode{ECANCELED} come codice di errore, ed mentre il valore di ritorno per
\func{aio\_return} sarà $-1$, inoltre il meccanismo di notifica non verrà
invocato. Se con \param{aiocbp} si specifica una operazione relativa ad un
file descriptor diverso da \param{fd} il risultato è indeterminato.  In caso
di successo, i possibili valori di ritorno per \func{aio\_cancel} (anch'essi
definiti in \headfile{aio.h}) sono tre:
\begin{basedescript}{\desclabelwidth{3.0cm}}
\item[\constd{AIO\_ALLDONE}] indica che le operazioni di cui si è richiesta la
  cancellazione sono state già completate,
  
\item[\constd{AIO\_CANCELED}] indica che tutte le operazioni richieste sono
  state cancellate,  
  
\item[\constd{AIO\_NOTCANCELED}] indica che alcune delle operazioni erano in
  corso e non sono state cancellate.
\end{basedescript}

Nel caso si abbia \const{AIO\_NOTCANCELED} occorrerà chiamare
\func{aio\_error} per determinare quali sono le operazioni effettivamente
cancellate. Le operazioni che non sono state cancellate proseguiranno il loro
corso normale, compreso quanto richiesto riguardo al meccanismo di notifica
del loro avvenuto completamento.

Benché l'I/O asincrono preveda un meccanismo di notifica, l'interfaccia
fornisce anche una apposita funzione, \funcd{aio\_suspend}, che permette di
sospendere l'esecuzione del processo chiamante fino al completamento di una
specifica operazione; il suo prototipo è:

\begin{funcproto}{
\fhead{aio.h}
\fdecl{int aio\_suspend(const struct aiocb * const list[], int nent, \\
\phantom{int aio\_suspend(}const struct timespec *timeout)}
\fdesc{Attende il completamento di una operazione di I/O asincrono.} 
}

{La funzione ritorna $0$ se una (o più) operazioni sono state completate e
  $-1$ per un errore, nel qual caso \var{errno} assumerà uno dei valori:
  \begin{errlist}
    \item[\errcode{EAGAIN}] nessuna operazione è stata completata entro
      \param{timeout}.
    \item[\errcode{EINTR}] la funzione è stata interrotta da un segnale.
    \item[\errcode{ENOSYS}] la funzione non è implementata.
  \end{errlist}
}
\end{funcproto}
  
La funzione permette di bloccare il processo fintanto che almeno una delle
\param{nent} operazioni specificate nella lista \param{list} è completata, per
un tempo massimo specificato dalla struttura \struct{timespec} puntata
da \param{timout}, o fintanto che non arrivi un segnale (si tenga conto che
questo segnale potrebbe essere anche quello utilizzato come meccanismo di
notifica). La lista deve essere inizializzata con delle strutture
\struct{aiocb} relative ad operazioni effettivamente richieste, ma può
contenere puntatori nulli, che saranno ignorati. In caso si siano specificati
valori non validi l'effetto è indefinito.  
Un valore \val{NULL} per \param{timout} comporta l'assenza di timeout, mentre
se si vuole effettuare un \textit{polling} sulle operazioni occorrerà
specificare un puntatore valido ad una struttura \texttt{timespec} (vedi
fig.~\ref{fig:sys_timespec_struct}) contenente valori nulli, e verificare poi
con \func{aio\_error} quale delle operazioni della lista \param{list} è stata
completata.

Lo standard POSIX.1b infine ha previsto pure una funzione, \funcd{lio\_listio},
che permette di effettuare la richiesta di una intera lista di operazioni di
lettura o scrittura; il suo prototipo è:


\begin{funcproto}{
\fhead{aio.h}
\fdecl{int lio\_listio(int mode, struct aiocb * const list[], int nent, struct
    sigevent *sig)}

\fdesc{Richiede l'esecuzione di una serie di operazioni di I/O.} 
}

{La funzione ritorna $0$ in caso di successo e $-1$ per un errore, nel qual
  caso \var{errno} assumerà uno dei valori: 
  \begin{errlist}
    \item[\errcode{EAGAIN}] nessuna operazione è stata completata entro
      \param{timeout}.
    \item[\errcode{EINTR}] la funzione è stata interrotta da un segnale.
    \item[\errcode{EINVAL}] si è passato un valore di \param{mode} non valido
      o un numero di operazioni \param{nent} maggiore di
      \const{AIO\_LISTIO\_MAX}.
    \item[\errcode{ENOSYS}] la funzione non è implementata.
  \end{errlist}
}
\end{funcproto}

La funzione esegue la richiesta delle \param{nent} operazioni indicate nella
lista \param{list} un vettore di puntatori a strutture \struct{aiocb}
indicanti le operazioni da compiere (che verranno eseguite senza un ordine
particolare). La lista può contenere anche puntatori nulli, che saranno
ignorati (si possono così eliminare facilmente componenti della lista senza
doverla rigenerare).

Ciascuna struttura \struct{aiocb} della lista deve contenere un
\textit{control block} opportunamente inizializzato; in particolare per
ognuna di esse dovrà essere specificato il tipo di operazione con il campo
\var{aio\_lio\_opcode}, che può prendere i valori:
\begin{basedescript}{\desclabelwidth{2.0cm}}
\item[\constd{LIO\_READ}]  si richiede una operazione di lettura.
\item[\constd{LIO\_WRITE}] si richiede una operazione di scrittura.
\item[\constd{LIO\_NOP}] non si effettua nessuna operazione.
\end{basedescript}
dove \const{LIO\_NOP} viene usato quando si ha a che fare con un vettore di
dimensione fissa, per poter specificare solo alcune operazioni, o quando si
sono dovute cancellare delle operazioni e si deve ripetere la richiesta per
quelle non completate. 

L'argomento \param{mode} controlla il comportamento della funzione, se viene
usato il valore \constd{LIO\_WAIT} la funzione si blocca fino al completamento
di tutte le operazioni richieste; se si usa \constd{LIO\_NOWAIT} la funzione
ritorna immediatamente dopo aver messo in coda tutte le richieste. In tal caso
il chiamante può richiedere la notifica del completamento di tutte le
richieste, impostando l'argomento \param{sig} in maniera analoga a come si fa
per il campo \var{aio\_sigevent} di \struct{aiocb}.

% TODO: trattare libaio e le system call del kernel per l'I/O asincrono, vedi
% http://lse.sourceforge.net/io/aio.html,
% http://webfiveoh.com/content/guides/2012/aug/mon-13th/linux-asynchronous-io-and-libaio.html, 
% https://code.google.com/p/kernel/wiki/AIOUserGuide,
% http://bert-hubert.blogspot.de/2012/05/on-linux-asynchronous-file-io.html 
% https://www.fsl.cs.sunysb.edu/~vass/linux-aio.txt

% TODO trattare la poll API basata sull'I/O asicrono, introdotta con il kernel
% 4.18, vedi  https://lwn.net/Articles/743714/,
% https://lwn.net/Articles/742978/, https://lwn.net/Articles/758324/
% http://git.infradead.org/users/hch/vfs.git/commit/d2d9e26c7cb6d95d521153897910080cf56c7fad
% Reverted

% TODO trattare la nuova API per l'I/O asincrono (io_uring), introdotta con il
% kernel 5.1, vedi https://lwn.net/Articles/776703/,
% https://lwn.net/ml/linux-fsdevel/20190112213011.1439-1-axboe@kernel.dk/ 

\section{Altre modalità di I/O avanzato}
\label{sec:file_advanced_io}

Oltre alle precedenti modalità di \textit{I/O multiplexing} e \textsl{I/O
  asincrono}, esistono altre funzioni che implementano delle modalità di
accesso ai file più evolute rispetto alle normali funzioni di lettura e
scrittura che abbiamo esaminato in sez.~\ref{sec:file_unix_interface}. In
questa sezione allora prenderemo in esame le interfacce per l'\textsl{I/O
  mappato in memoria}, per l'\textsl{I/O vettorizzato} e altre funzioni di I/O
avanzato.


\subsection{File mappati in memoria}
\label{sec:file_memory_map}

\itindbeg{memory~mapping}

Una modalità alternativa di I/O, che usa una interfaccia completamente diversa
rispetto a quella classica vista in sez.~\ref{sec:file_unix_interface}, è il
cosiddetto \textit{memory-mapped I/O}, che attraverso il meccanismo della
\textsl{paginazione}  usato dalla memoria virtuale (vedi
sez.~\ref{sec:proc_mem_gen}) permette di \textsl{mappare} il contenuto di un
file in una sezione dello spazio di indirizzi del processo che lo ha allocato.

\begin{figure}[htb]
  \centering
  \includegraphics[width=12cm]{img/mmap_layout}
  \caption{Disposizione della memoria di un processo quando si esegue la
  mappatura in memoria di un file.}
  \label{fig:file_mmap_layout}
\end{figure}

Il meccanismo è illustrato in fig.~\ref{fig:file_mmap_layout}, una sezione del
file viene \textsl{mappata} direttamente nello spazio degli indirizzi del
programma.  Tutte le operazioni di lettura e scrittura su variabili contenute
in questa zona di memoria verranno eseguite leggendo e scrivendo dal contenuto
del file attraverso il sistema della memoria virtuale illustrato in
sez.~\ref{sec:proc_mem_gen} che in maniera analoga a quanto avviene per le
pagine che vengono salvate e rilette nella \textit{swap}, si incaricherà di
sincronizzare il contenuto di quel segmento di memoria con quello del file
mappato su di esso.  Per questo motivo si può parlare tanto di \textsl{file
  mappato in memoria}, quanto di \textsl{memoria mappata su file}.

L'uso del \textit{memory-mapping} comporta una notevole semplificazione delle
operazioni di I/O, in quanto non sarà più necessario utilizzare dei buffer
intermedi su cui appoggiare i dati da traferire, poiché questi potranno essere
acceduti direttamente nella sezione di memoria mappata; inoltre questa
interfaccia è più efficiente delle usuali funzioni di I/O, in quanto permette
di caricare in memoria solo le parti del file che sono effettivamente usate ad
un dato istante.

Infatti, dato che l'accesso è fatto direttamente attraverso la memoria
virtuale, la sezione di memoria mappata su cui si opera sarà a sua volta letta
o scritta sul file una pagina alla volta e solo per le parti effettivamente
usate, il tutto in maniera completamente trasparente al processo; l'accesso
alle pagine non ancora caricate avverrà allo stesso modo con cui vengono
caricate in memoria le pagine che sono state salvate sullo \textit{swap}.

Infine in situazioni in cui la memoria è scarsa, le pagine che mappano un file
vengono salvate automaticamente, così come le pagine dei programmi vengono
scritte sulla \textit{swap}; questo consente di accedere ai file su dimensioni
il cui solo limite è quello dello spazio di indirizzi disponibile, e non della
memoria su cui possono esserne lette delle porzioni.

L'interfaccia POSIX implementata da Linux prevede varie funzioni di sistema
per la gestione del \textit{memory mapped I/O}, la prima di queste, che serve
ad eseguire la mappatura in memoria di un file, è \funcd{mmap}; il suo
prototipo è:

\begin{funcproto}{
%\fhead{unistd.h}
\fhead{sys/mman.h} 
\fdecl{void * mmap(void * start, size\_t length, int prot, int flags, int
    fd, off\_t offset)}
\fdesc{Esegue la mappatura in memoria di una sezione di un file.} 
}

{La funzione ritorna il puntatore alla zona di memoria mappata in caso di
  successo, e \const{MAP\_FAILED} (\texttt{(void *) -1}) per un errore, nel
  qual caso \var{errno} assumerà uno dei valori:
  \begin{errlist}
    \item[\errcode{EACCES}] o \param{fd} non si riferisce ad un file regolare,
      o si è usato \const{MAP\_PRIVATE} ma \param{fd} non è aperto in lettura,
      o si è usato \const{MAP\_SHARED} e impostato \const{PROT\_WRITE} ed
      \param{fd} non è aperto in lettura/scrittura, o si è impostato
      \const{PROT\_WRITE} ed \param{fd} è in \textit{append-only}.
    \item[\errcode{EAGAIN}] il file è bloccato, o si è bloccata troppa memoria
      rispetto a quanto consentito dai limiti di sistema (vedi
      sez.~\ref{sec:sys_resource_limit}).
    \item[\errcode{EBADF}] il file descriptor non è valido, e non si è usato
      \const{MAP\_ANONYMOUS}.
    \item[\errcode{EINVAL}] i valori di \param{start}, \param{length} o
      \param{offset} non sono validi (o troppo grandi o non allineati sulla
      dimensione delle pagine), o \param{lengh} è zero (solo dal 2.6.12)
      o \param{flags} contiene sia \const{MAP\_PRIVATE} che
      \const{MAP\_SHARED} o nessuno dei due.
    \item[\errcode{ENFILE}] si è superato il limite del sistema sul numero di
      file aperti (vedi sez.~\ref{sec:sys_resource_limit}).
    \item[\errcode{ENODEV}] il filesystem di \param{fd} non supporta il memory
      mapping.
    \item[\errcode{ENOMEM}] non c'è memoria o si è superato il limite sul
      numero di mappature possibili.
    \item[\errcode{EOVERFLOW}] su architettura a 32 bit con il supporto per i
      \textit{large file} (che hanno una dimensione a 64 bit) il numero di
      pagine usato per \param{lenght} aggiunto a quello usato
      per \param{offset} eccede i 32 bit (\texttt{unsigned long}).
    \item[\errcode{EPERM}] l'argomento \param{prot} ha richiesto
      \const{PROT\_EXEC}, ma il filesystem di \param{fd} è montato con
      l'opzione \texttt{noexec}.
    \item[\errcode{ETXTBSY}] si è impostato \const{MAP\_DENYWRITE} ma
      \param{fd} è aperto in scrittura.
  \end{errlist}
}
\end{funcproto}

La funzione richiede di mappare in memoria la sezione del file \param{fd} a
partire da \param{offset} per \param{length} byte, preferibilmente
all'indirizzo \param{start}. Il valore \param{start} viene normalmente
considerato come un suggerimento, ma l'uso di un qualunque valore diverso da
\val{NULL}, in cui si rimette completamente al kernel la scelta
dell'indirizzo, viene sconsigliato per ragioni di portabilità. Il valore
di \param{offset} deve essere un multiplo della dimensione di una pagina di
memoria.

\begin{table}[htb]
  \centering
  \footnotesize
  \begin{tabular}[c]{|l|l|}
    \hline
    \textbf{Valore} & \textbf{Significato} \\
    \hline
    \hline
    \constd{PROT\_EXEC}  & Le pagine possono essere eseguite.\\
    \constd{PROT\_READ}  & Le pagine possono essere lette.\\
    \constd{PROT\_WRITE} & Le pagine possono essere scritte.\\
    \constd{PROT\_NONE}  & L'accesso alle pagine è vietato.\\
    \hline    
  \end{tabular}
  \caption{Valori dell'argomento \param{prot} di \func{mmap}, relativi alla
    protezione applicate alle pagine del file mappate in memoria.}
  \label{tab:file_mmap_prot}
\end{table}

Il valore dell'argomento \param{prot} indica la protezione\footnote{come
  accennato in sez.~\ref{sec:proc_memory} in Linux la memoria reale è divisa
  in pagine, ogni processo vede la sua memoria attraverso uno o più segmenti
  lineari di memoria virtuale; per ciascuno di questi segmenti il kernel
  mantiene nella \textit{page table} la mappatura sulle pagine di memoria
  reale, ed le modalità di accesso (lettura, esecuzione, scrittura); una loro
  violazione causa quella una \textit{segment violation}, e la relativa
  emissione del segnale \signal{SIGSEGV}.} da applicare al segmento di memoria
e deve essere specificato come maschera binaria ottenuta dall'OR di uno o più
dei valori riportati in tab.~\ref{tab:file_mmap_prot}; il valore specificato
deve essere compatibile con la modalità di accesso con cui si è aperto il
file.

\begin{table}[!htb]
  \centering
  \footnotesize
  \begin{tabular}[c]{|l|p{11cm}|}
    \hline
    \textbf{Valore} & \textbf{Significato} \\
    \hline
    \hline
    \constd{MAP\_32BIT}    & Esegue la mappatura sui primi 2Gb dello spazio
                             degli indirizzi, viene supportato solo sulle
                             piattaforme \texttt{x86-64} per compatibilità con
                             le applicazioni a 32 bit. Viene ignorato se si è
                             richiesto \const{MAP\_FIXED} (dal kernel 2.4.20).\\
    \constd{MAP\_ANON}     & Sinonimo di \const{MAP\_ANONYMOUS}, deprecato.\\
    \constd{MAP\_ANONYMOUS}& La mappatura non è associata a nessun file. Gli
                             argomenti \param{fd} e \param{offset} sono
                             ignorati. L'uso di questo flag con
                             \const{MAP\_SHARED} è stato implementato in Linux
                             a partire dai kernel della serie 2.4.x.\\
    \constd{MAP\_DENYWRITE}& In Linux viene ignorato per evitare
                             \textit{DoS}
                             (veniva usato per segnalare che tentativi di
                             scrittura sul file dovevano fallire con
                             \errcode{ETXTBSY}).\\ 
    \constd{MAP\_EXECUTABLE}& Ignorato.\\
    \constd{MAP\_FILE}     & Valore di compatibilità, ignorato.\\
    \constd{MAP\_FIXED}    & Non permette di restituire un indirizzo diverso
                             da \param{start}, se questo non può essere usato
                             \func{mmap} fallisce. Se si imposta questo flag il
                             valore di \param{start} deve essere allineato
                             alle dimensioni di una pagina.\\
    \constd{MAP\_GROWSDOWN}& Usato per gli \textit{stack}. 
                             Indica che la mappatura deve essere effettuata 
                             con gli indirizzi crescenti verso il basso.\\
    \constd{MAP\_HUGETLB}  & Esegue la mappatura usando le cosiddette
                             ``\textit{huge pages}'' (dal kernel 2.6.32).\\
    \constd{MAP\_LOCKED}   & Se impostato impedisce lo \textit{swapping} delle
                             pagine mappate (dal kernel 2.5.37).\\
    \constd{MAP\_NONBLOCK} & Esegue un \textit{prefaulting} più limitato che
                             non causa I/O (dal kernel 2.5.46).\\
    \constd{MAP\_NORESERVE}& Si usa con \const{MAP\_PRIVATE}. Non riserva
                             delle pagine di \textit{swap} ad uso del meccanismo
                             del \textit{copy on write} 
                             per mantenere le modifiche fatte alla regione
                             mappata, in questo caso dopo una scrittura, se
                             non c'è più memoria disponibile, si ha
                             l'emissione di un \signal{SIGSEGV}.\\
    \constd{MAP\_POPULATE} & Esegue il \textit{prefaulting} delle pagine di
                             memoria necessarie alla mappatura (dal kernel
                             2.5.46).\\ 
    \constd{MAP\_PRIVATE}  & I cambiamenti sulla memoria mappata non vengono
                             riportati sul file. Ne viene fatta una copia
                             privata cui solo il processo chiamante ha
                             accesso.  Incompatibile con \const{MAP\_SHARED}.\\
    \constd{MAP\_SHARED}   & I cambiamenti sulla memoria mappata vengono
                             riportati sul file e saranno immediatamente
                             visibili agli altri processi che mappano lo stesso
                             file. Incompatibile
                             con \const{MAP\_PRIVATE}.\\ 
    \const{MAP\_STACK}     & Al momento è ignorato, è stato fornito (dal kernel
                             2.6.27) a supporto della implementazione dei
                             \textit{thread} nella \acr{glibc}, per allocare
                             memoria in uno spazio utilizzabile come
                             \textit{stack} per le architetture hardware che
                             richiedono un trattamento speciale di
                             quest'ultimo.\\ 
    \constd{MAP\_UNINITIALIZED}& Specifico per i sistemi embedded ed
                             utilizzabile dal kernel 2.6.33 solo se è stata
                             abilitata in fase di compilazione dello stesso
                             l'opzione
                             \texttt{CONFIG\_MMAP\_ALLOW\_UNINITIALIZED}. Se
                             usato le pagine di memoria usate nella mappatura
                             anonima non vengono cancellate; questo migliora
                             le prestazioni sui sistemi con risorse minime, ma
                             comporta la possibilità di rileggere i dati di
                             altri processi che han chiuso una mappatura, per
                             cui viene usato solo quando (come si suppone sia
                             per i sistemi embedded) si ha il completo
                             controllo dell'uso della memoria da parte degli
                             utenti.\\ 
%    \constd{MAP\_DONTEXPAND}& Non consente una successiva espansione dell'area
%                              mappata con \func{mremap}, proposto ma pare non
%                              implementato.\\
    \hline
  \end{tabular}
  \caption{Valori possibili dell'argomento \param{flag} di \func{mmap}.}
  \label{tab:file_mmap_flag}
\end{table}

% TODO trattare MAP_HUGETLB introdotto con il kernel 2.6.32, e modifiche
% introdotte con il 3.8 per le dimensioni variabili delle huge pages

% TODO trattare  MAP_FIXED_NOREPLACE vedi https://lwn.net/Articles/751651/ e
% https://lwn.net/Articles/741369/ 

% TODO: verificare MAP_SYNC e MAP_SHARED_VALIDATE, vedi
% https://lwn.net/Articles/731706/, https://lwn.net/Articles/758594/ incluse
% con il 4.15


L'argomento \param{flags} specifica infine qual è il tipo di oggetto mappato,
le opzioni relative alle modalità con cui è effettuata la mappatura e alle
modalità con cui le modifiche alla memoria mappata vengono condivise o
mantenute private al processo che le ha effettuate. Deve essere specificato
come maschera binaria ottenuta dall'OR di uno o più dei valori riportati in
tab.~\ref{tab:file_mmap_flag}. Fra questi comunque deve sempre essere
specificato o \const{MAP\_PRIVATE} o \const{MAP\_SHARED} per indicare la
modalità con cui viene effettuata la mappatura.

Esistono infatti due modalità alternative di eseguire la mappatura di un file;
la più comune è \const{MAP\_SHARED} in cui la memoria è condivisa e le
modifiche effettuate su di essa sono visibili a tutti i processi che hanno
mappato lo stesso file. In questo caso le modifiche vengono anche riportate su
disco, anche se questo può non essere immediato a causa della bufferizzazione:
si potrà essere sicuri dell'aggiornamento solo in seguito alla chiamata di
\func{msync} o \func{munmap}, e solo allora le modifiche saranno visibili sul
file con l'I/O convenzionale.

Con \const{MAP\_PRIVATE} invece viene creata una copia privata del file,
questo non viene mai modificato e solo il processo chiamante ha accesso alla
mappatura. Le modifiche eseguite dal processo sulla mappatura vengono
effettuate utilizzando il meccanismo del \textit{copy on write}, mentenute in
memoria e salvate su \textit{swap} in caso di necessità.  Non è specificato se
i cambiamenti sul file originale vengano riportati sulla regione mappata.

Gli altri valori di \func{flag} modificano le caratteristiche della
mappatura. Fra questi il più rilevante è probabilmente \const{MAP\_ANONYMOUS}
che consente di creare segmenti di memoria condivisa fra processi diversi
senza appoggiarsi a nessun file (torneremo sul suo utilizzo in
sez.~\ref{sec:ipc_mmap_anonymous}). In tal caso gli argomenti \param{fd}
e \param{offset} vangono ignorati, anche se alcune implementazioni richiedono
che invece \param{fd} sia $-1$, convenzione che è opportuno seguire se si ha a
cuore la portabilità dei programmi.

Gli effetti dell'accesso ad una zona di memoria mappata su file possono essere
piuttosto complessi, essi si possono comprendere solo tenendo presente che
tutto quanto è comunque basato sul meccanismo della memoria virtuale. Questo
comporta allora una serie di conseguenze. La più ovvia è che se si cerca di
scrivere su una zona mappata in sola lettura si avrà l'emissione di un segnale
di violazione di accesso (\signal{SIGSEGV}), dato che i permessi sul segmento
di memoria relativo non consentono questo tipo di accesso.

È invece assai diversa la questione relativa agli accessi al di fuori della
regione di cui si è richiesta la mappatura. A prima vista infatti si potrebbe
ritenere che anch'essi debbano generare un segnale di violazione di accesso;
questo però non tiene conto del fatto che, essendo basata sul meccanismo della
paginazione, la mappatura in memoria non può che essere eseguita su un
segmento di dimensioni rigorosamente multiple di quelle di una pagina, ed in
generale queste potranno non corrispondere alle dimensioni effettive del file
o della sezione che si vuole mappare.

\begin{figure}[!htb] 
  \centering
  \includegraphics[height=6cm]{img/mmap_boundary}
  \caption{Schema della mappatura in memoria di una sezione di file di
    dimensioni non corrispondenti al bordo di una pagina.}
  \label{fig:file_mmap_boundary}
\end{figure}

Il caso più comune è quello illustrato in fig.~\ref{fig:file_mmap_boundary},
in cui la sezione di file non rientra nei confini di una pagina: in tal caso
il file sarà mappato su un segmento di memoria che si estende fino al
bordo della pagina successiva.  In questo caso è possibile accedere a quella
zona di memoria che eccede le dimensioni specificate da \param{length}, senza
ottenere un \signal{SIGSEGV} poiché essa è presente nello spazio di indirizzi
del processo, anche se non è mappata sul file. Il comportamento del sistema è
quello di restituire un valore nullo per quanto viene letto, e di non
riportare su file quanto viene scritto.

Un caso più complesso è quello che si viene a creare quando le dimensioni del
file mappato sono più corte delle dimensioni della mappatura, oppure quando il
file è stato troncato, dopo che è stato mappato, ad una dimensione inferiore a
quella della mappatura in memoria.  In questa situazione, per la sezione di
pagina parzialmente coperta dal contenuto del file, vale esattamente quanto
visto in precedenza; invece per la parte che eccede, fino alle dimensioni date
da \param{length}, l'accesso non sarà più possibile, ma il segnale emesso non
sarà \signal{SIGSEGV}, ma \signal{SIGBUS}, come illustrato in
fig.~\ref{fig:file_mmap_exceed}.

Non tutti i file possono venire mappati in memoria, dato che, come illustrato
in fig.~\ref{fig:file_mmap_layout}, la mappatura introduce una corrispondenza
biunivoca fra una sezione di un file ed una sezione di memoria. Questo
comporta che ad esempio non è possibile mappare in memoria file descriptor
relativi a \textit{pipe}, socket e \textit{fifo}, per i quali non ha senso
parlare di \textsl{sezione}. Lo stesso vale anche per alcuni file di
dispositivo, che non dispongono della relativa operazione \func{mmap} (si
ricordi quanto esposto in sez.~\ref{sec:file_vfs_work}). Si tenga presente
però che esistono anche casi di dispositivi (un esempio è l'interfaccia al
ponte PCI-VME del chip Universe) che sono utilizzabili solo con questa
interfaccia.

\begin{figure}[htb]
  \centering
  \includegraphics[height=6cm]{img/mmap_exceed}
  \caption{Schema della mappatura in memoria di file di dimensioni inferiori
    alla lunghezza richiesta.}
  \label{fig:file_mmap_exceed}
\end{figure}

Dato che passando attraverso una \func{fork} lo spazio di indirizzi viene
copiato integralmente, i file mappati in memoria verranno ereditati in maniera
trasparente dal processo figlio, mantenendo gli stessi attributi avuti nel
padre; così se si è usato \const{MAP\_SHARED} padre e figlio accederanno allo
stesso file in maniera condivisa, mentre se si è usato \const{MAP\_PRIVATE}
ciascuno di essi manterrà una sua versione privata indipendente. Non c'è
invece nessun passaggio attraverso una \func{exec}, dato che quest'ultima
sostituisce tutto lo spazio degli indirizzi di un processo con quello di un
nuovo programma.

Quando si effettua la mappatura di un file vengono pure modificati i tempi ad
esso associati (di cui si è trattato in sez.~\ref{sec:file_file_times}). Il
valore di \var{st\_atime} può venir cambiato in qualunque istante a partire
dal momento in cui la mappatura è stata effettuata: il primo riferimento ad
una pagina mappata su un file aggiorna questo tempo.  I valori di
\var{st\_ctime} e \var{st\_mtime} possono venir cambiati solo quando si è
consentita la scrittura sul file (cioè per un file mappato con
\const{PROT\_WRITE} e \const{MAP\_SHARED}) e sono aggiornati dopo la scrittura
o in corrispondenza di una eventuale \func{msync}.

Dato per i file mappati in memoria le operazioni di I/O sono gestite
direttamente dalla memoria virtuale, occorre essere consapevoli delle
interazioni che possono esserci con operazioni effettuate con l'interfaccia
dei file ordinaria illustrata in sez.~\ref{sec:file_unix_interface}. Il
problema è che una volta che si è mappato un file, le operazioni di lettura e
scrittura saranno eseguite sulla memoria, e riportate su disco in maniera
autonoma dal sistema della memoria virtuale.

Pertanto se si modifica un file con l'interfaccia ordinaria queste modifiche
potranno essere visibili o meno a seconda del momento in cui la memoria
virtuale trasporterà dal disco in memoria quella sezione del file, perciò è
del tutto imprevedibile il risultato della modifica di un file nei confronti
del contenuto della memoria su cui è mappato.

Per questo è sempre sconsigliabile eseguire scritture su un file attraverso
l'interfaccia ordinaria quando lo si è mappato in memoria, è invece possibile
usare l'interfaccia ordinaria per leggere un file mappato in memoria, purché
si abbia una certa cura; infatti l'interfaccia dell'I/O mappato in memoria
mette a disposizione la funzione \funcd{msync} per sincronizzare il contenuto
della memoria mappata con il file su disco; il suo prototipo è:

\begin{funcproto}{
%\fhead{unistd.h}
\fhead{sys/mman.h}
\fdecl{int msync(const void *start, size\_t length, int flags)}
\fdesc{Sincronizza i contenuti di una sezione di un file mappato in memoria.} 
}

{La funzione ritorna $0$ in caso di successo e $-1$ per un errore, nel qual
  caso \var{errno} assumerà uno dei valori: 
  \begin{errlist}
    \item[\errcode{EBUSY}] si è indicato \const{MS\_INVALIDATE} ma
      nell'intervallo di memoria specificato è presente un \textit{memory lock}.
    \item[\errcode{EFAULT}] l'intervallo indicato, o parte di esso, non
      risulta mappato (prima del kernel 2.4.19).
    \item[\errcode{EINVAL}] o \param{start} non è multiplo di
      \const{PAGE\_SIZE}, o si è specificato un valore non valido per
      \param{flags}.
    \item[\errcode{ENOMEM}] l'intervallo indicato, o parte di esso, non
      risulta mappato (dal kernel 2.4.19).
  \end{errlist}
}
\end{funcproto}

La funzione esegue la sincronizzazione di quanto scritto nella sezione di
memoria indicata da \param{start} e \param{offset}, scrivendo le modifiche sul
file (qualora questo non sia già stato fatto).  Provvede anche ad aggiornare i
relativi tempi di modifica. In questo modo si è sicuri che dopo l'esecuzione
di \func{msync} le funzioni dell'interfaccia ordinaria troveranno un contenuto
del file aggiornato.

\begin{table}[htb]
  \centering
  \footnotesize
  \begin{tabular}[c]{|l|p{11cm}|}
    \hline
    \textbf{Valore} & \textbf{Significato} \\
    \hline
    \hline
    \constd{MS\_SYNC}      & richiede una sincronizzazione e ritorna soltanto
                             quando questa è stata completata.\\
    \constd{MS\_ASYNC}     & richiede una sincronizzazione, ma ritorna subito 
                             non attendendo che questa sia finita.\\
    \constd{MS\_INVALIDATE}& invalida le pagine per tutte le mappature
                             in memoria così da rendere necessaria una
                             rilettura immediata delle stesse.\\
    \hline
  \end{tabular}
  \caption{Valori possibili dell'argomento \param{flag} di \func{msync}.}
  \label{tab:file_mmap_msync}
\end{table}

L'argomento \param{flag} è specificato come maschera binaria composta da un OR
dei valori riportati in tab.~\ref{tab:file_mmap_msync}, di questi però
\const{MS\_ASYNC} e \const{MS\_SYNC} sono incompatibili; con il primo valore
infatti la funzione si limita ad inoltrare la richiesta di sincronizzazione al
meccanismo della memoria virtuale, ritornando subito, mentre con il secondo
attende che la sincronizzazione sia stata effettivamente eseguita. Il terzo
valore fa sì che vengano invalidate, per tutte le mappature dello stesso file,
le pagine di cui si è richiesta la sincronizzazione, così che esse possano
essere immediatamente aggiornate con i nuovi valori.

Una volta che si sono completate le operazioni di I/O si può eliminare la
mappatura della memoria usando la funzione \funcd{munmap}, il suo prototipo è:

\begin{funcproto}{
%\fhead{unistd.h}
\fhead{sys/mman.h}
\fdecl{int munmap(void *start, size\_t length)}
\fdesc{Rilascia la mappatura sulla sezione di memoria specificata.} 
}

{La funzione ritorna $0$ in caso di successo e $-1$ per un errore, nel qual
  caso \var{errno} assumerà uno dei valori: 
  \begin{errlist}
    \item[\errcode{EINVAL}] l'intervallo specificato non ricade in una zona
      precedentemente mappata.
  \end{errlist}
}
\end{funcproto}

La funzione cancella la mappatura per l'intervallo specificato con
\param{start} e \param{length}; ogni successivo accesso a tale regione causerà
un errore di accesso in memoria. L'argomento \param{start} deve essere
allineato alle dimensioni di una pagina, e la mappatura di tutte le pagine
contenute anche parzialmente nell'intervallo indicato, verrà rimossa.
Indicare un intervallo che non contiene mappature non è un errore.  Si tenga
presente inoltre che alla conclusione di un processo ogni pagina mappata verrà
automaticamente rilasciata, mentre la chiusura del file descriptor usato per
il \textit{memory mapping} non ha alcun effetto su di esso.

Lo standard POSIX prevede anche una funzione che permetta di cambiare le
protezioni delle pagine di memoria; lo standard prevede che essa si applichi
solo ai \textit{memory mapping} creati con \func{mmap}, ma nel caso di Linux
la funzione può essere usata con qualunque pagina valida nella memoria
virtuale. Questa funzione di sistema è \funcd{mprotect} ed il suo prototipo è:

\begin{funcproto}{
\fhead{sys/mman.h} 
\fdecl{int mprotect(const void *addr, size\_t len, int prot)}
\fdesc{Modifica le protezioni delle pagine di memoria.} 
}

{La funzione ritorna $0$ in caso di successo e $-1$ per un errore, nel qual
  caso \var{errno} assumerà uno dei valori: 
  \begin{errlist}
    \item[\errcode{EINVAL}] il valore di \param{addr} non è valido o non è un
      multiplo di \const{PAGE\_SIZE}.
    \item[\errcode{EACCES}] l'operazione non è consentita, ad esempio si è
      cercato di marcare con \const{PROT\_WRITE} un segmento di memoria cui si
      ha solo accesso in lettura.
    \item[\errcode{ENOMEM}] non è stato possibile allocare le risorse
      necessarie all'interno del kernel o si è specificato un indirizzo di
      memoria non valido del processo o non corrispondente a pagine mappate
      (negli ultimi due casi prima del kernel 2.4.19 veniva prodotto,
      erroneamente, \errcode{EFAULT}).
  \end{errlist}
}
\end{funcproto}

La funzione prende come argomenti un indirizzo di partenza in \param{addr},
allineato alle dimensioni delle pagine di memoria, ed una dimensione
\param{size}. La nuova protezione deve essere specificata in \param{prot} con
una combinazione dei valori di tab.~\ref{tab:file_mmap_prot}.  La nuova
protezione verrà applicata a tutte le pagine contenute, anche parzialmente,
dall'intervallo fra \param{addr} e \param{addr}+\param{size}-1.

Infine Linux supporta alcune operazioni specifiche non disponibili su altri
kernel unix-like per poter usare le quali occorre però dichiarare
\macro{\_GNU\_SOURCE} prima dell'inclusione di \texttt{sys/mman.h}. La prima
di queste è la possibilità di modificare un precedente \textit{memory
  mapping}, ad esempio per espanderlo o restringerlo.  Questo è realizzato
dalla funzione di sistema \funcd{mremap}, il cui prototipo è:

\begin{funcproto}{
\fhead{sys/mman.h} 
\fdecl{void * mremap(void *old\_address, size\_t old\_size , size\_t
    new\_size, unsigned long flags)}
\fdesc{Restringe o allarga una mappatura in memoria.} 
}

{La funzione ritorna l'indirizzo alla nuova area di memoria in caso di
  successo o il valore \const{MAP\_FAILED} (pari a \texttt{(void *) -1}), nel
  qual caso \var{errno} assumerà uno dei valori:
  \begin{errlist}
    \item[\errcode{EINVAL}] il valore di \param{old\_address} non è un
      puntatore valido.
    \item[\errcode{EFAULT}] ci sono indirizzi non validi nell'intervallo
      specificato da \param{old\_address} e \param{old\_size}, o ci sono altre
      mappature di tipo non corrispondente a quella richiesta.
    \item[\errcode{ENOMEM}] non c'è memoria sufficiente oppure l'area di
      memoria non può essere espansa all'indirizzo virtuale corrente, e non si
      è specificato \const{MREMAP\_MAYMOVE} nei flag.
    \item[\errcode{EAGAIN}] il segmento di memoria scelto è bloccato e non può
      essere rimappato.
  \end{errlist}
}
\end{funcproto}

La funzione richiede come argomenti \param{old\_address} (che deve essere
allineato alle dimensioni di una pagina di memoria) che specifica il
precedente indirizzo del \textit{memory mapping} e \param{old\_size}, che ne
indica la dimensione. Con \param{new\_size} si specifica invece la nuova
dimensione che si vuole ottenere. Infine l'argomento \param{flags} è una
maschera binaria per i flag che controllano il comportamento della funzione.
Il solo valore utilizzato è \constd{MREMAP\_MAYMOVE} che consente di eseguire
l'espansione anche quando non è possibile utilizzare il precedente
indirizzo. Per questo motivo, se si è usato questo flag, la funzione può
restituire un indirizzo della nuova zona di memoria che non è detto coincida
con \param{old\_address}.

La funzione si appoggia al sistema della memoria virtuale per modificare
l'associazione fra gli indirizzi virtuali del processo e le pagine di memoria,
modificando i dati direttamente nella \textit{page table} del processo. Come
per \func{mprotect} la funzione può essere usata in generale, anche per pagine
di memoria non corrispondenti ad un \textit{memory mapping}, e consente così
di implementare la funzione \func{realloc} in maniera molto efficiente.

Una caratteristica comune a tutti i sistemi unix-like è che la mappatura in
memoria di un file viene eseguita in maniera lineare, cioè parti successive di
un file vengono mappate linearmente su indirizzi successivi in memoria.
Esistono però delle applicazioni (in particolare la tecnica è usata dai
database o dai programmi che realizzano macchine virtuali) in cui è utile
poter mappare sezioni diverse di un file su diverse zone di memoria.

Questo è ovviamente sempre possibile eseguendo ripetutamente la funzione
\func{mmap} per ciascuna delle diverse aree del file che si vogliono mappare
in sequenza non lineare (ed in effetti è quello che veniva fatto anche con
Linux prima che fossero introdotte queste estensioni) ma questo approccio ha
delle conseguenze molto pesanti in termini di prestazioni.  Infatti per
ciascuna mappatura in memoria deve essere definita nella \textit{page table}
del processo una nuova area di memoria virtuale, quella che nel gergo del
kernel viene chiamata VMA (\textit{virtual memory area}, che corrisponda alla
mappatura, in modo che questa diventi visibile nello spazio degli indirizzi
come illustrato in fig.~\ref{fig:file_mmap_layout}.

Quando un processo esegue un gran numero di mappature diverse (si può arrivare
anche a centinaia di migliaia) per realizzare a mano una mappatura non-lineare
esso vedrà un accrescimento eccessivo della sua \textit{page table}, e lo
stesso accadrà per tutti gli altri processi che utilizzano questa tecnica. In
situazioni in cui le applicazioni hanno queste esigenze si avranno delle
prestazioni ridotte, dato che il kernel dovrà impiegare molte risorse per
mantenere i dati relativi al \textit{memory mapping}, sia in termini di
memoria interna per i dati delle \textit{page table}, che di CPU per il loro
aggiornamento.

Per questo motivo con il kernel 2.5.46 è stato introdotto, ad opera di Ingo
Molnar, un meccanismo che consente la mappatura non-lineare. Anche questa è
una caratteristica specifica di Linux, non presente in altri sistemi
unix-like.  Diventa così possibile utilizzare una sola mappatura iniziale, e
quindi una sola \textit{virtual memory area} nella \textit{page table} del
processo, e poi rimappare a piacere all'interno di questa i dati del file. Ciò
è possibile grazie ad una nuova \textit{system call},
\funcd{remap\_file\_pages}, il cui prototipo è:

\begin{funcproto}{
\fhead{sys/mman.h} 
\fdecl{int remap\_file\_pages(void *start, size\_t size, int prot,
    ssize\_t pgoff, int flags)}
\fdesc{Rimappa non linearmente un \textit{memory mapping}.} 
}

{La funzione ritorna $0$ in caso di successo e $-1$ per un errore, nel qual
  caso \var{errno} assumerà uno dei valori: 
  \begin{errlist}
    \item[\errcode{EINVAL}] si è usato un valore non valido per uno degli
      argomenti o \param{start} non fa riferimento ad un \textit{memory
        mapping} valido creato con \const{MAP\_SHARED}.
  \end{errlist}
  ed inoltre 
 nel loro significato generico.}
\end{funcproto}

Per poter utilizzare questa funzione occorre anzitutto effettuare
preliminarmente una chiamata a \func{mmap} con \const{MAP\_SHARED} per
definire l'area di memoria che poi sarà rimappata non linearmente. Poi si
chiamerà questa funzione per modificare le corrispondenze fra pagine di
memoria e pagine del file; si tenga presente che \func{remap\_file\_pages}
permette anche di mappare la stessa pagina di un file in più pagine della
regione mappata.

La funzione richiede che si identifichi la sezione del file che si vuole
riposizionare all'interno del \textit{memory mapping} con gli argomenti
\param{pgoff} e \param{size}; l'argomento \param{start} invece deve indicare
un indirizzo all'interno dell'area definita dall'\func{mmap} iniziale, a
partire dal quale la sezione di file indicata verrà rimappata. L'argomento
\param{prot} deve essere sempre nullo, mentre \param{flags} prende gli stessi
valori di \func{mmap} (quelli di tab.~\ref{tab:file_mmap_prot}) ma di tutti i
flag solo \const{MAP\_NONBLOCK} non viene ignorato.

\itindbeg{prefaulting} 

Insieme alla funzione \func{remap\_file\_pages} nel kernel 2.5.46 con sono
stati introdotti anche due nuovi flag per \func{mmap}: \const{MAP\_POPULATE} e
\const{MAP\_NONBLOCK}.  Il primo dei due consente di abilitare il meccanismo
del \textit{prefaulting}. Questo viene di nuovo in aiuto per migliorare le
prestazioni in certe condizioni di utilizzo del \textit{memory mapping}.

Il problema si pone tutte le volte che si vuole mappare in memoria un file di
grosse dimensioni. Il comportamento normale del sistema della memoria virtuale
è quello per cui la regione mappata viene aggiunta alla \textit{page table}
del processo, ma i dati verranno effettivamente utilizzati (si avrà cioè un
\textit{page fault} che li trasferisce dal disco alla memoria) soltanto in
corrispondenza dell'accesso a ciascuna delle pagine interessate dal
\textit{memory mapping}.

Questo vuol dire che il passaggio dei dati dal disco alla memoria avverrà una
pagina alla volta con un gran numero di \textit{page fault}, chiaramente se si
sa in anticipo che il file verrà utilizzato immediatamente, è molto più
efficiente eseguire un \textit{prefaulting} in cui tutte le pagine di memoria
interessate alla mappatura vengono ``\textsl{popolate}'' in una sola volta,
questo comportamento viene abilitato quando si usa con \func{mmap} il flag
\const{MAP\_POPULATE}.

Dato che l'uso di \const{MAP\_POPULATE} comporta dell'I/O su disco che può
rallentare l'esecuzione di \func{mmap} è stato introdotto anche un secondo
flag, \const{MAP\_NONBLOCK}, che esegue un \textit{prefaulting} più limitato
in cui vengono popolate solo le pagine della mappatura che già si trovano
nella cache del kernel.\footnote{questo può essere utile per il linker
  dinamico, in particolare quando viene effettuato il \textit{prelink} delle
  applicazioni.}

\itindend{prefaulting}

Per i vantaggi illustrati all'inizio del paragrafo l'interfaccia del
\textit{memory mapped I/O} viene usata da una grande varietà di programmi,
spesso con esigenze molto diverse fra di loro riguardo le modalità con cui
verranno eseguiti gli accessi ad un file; è ad esempio molto comune per i
database effettuare accessi ai dati in maniera pressoché casuale, mentre un
riproduttore audio o video eseguirà per lo più letture sequenziali.

\itindend{memory~mapping}

Per migliorare le prestazioni a seconda di queste modalità di accesso è
disponibile una apposita funzione, \funcd{madvise},\footnote{tratteremo in
  sez.~\ref{sec:file_fadvise} le funzioni che consentono di ottimizzare
  l'accesso ai file con l'interfaccia classica.} che consente di fornire al
kernel delle indicazioni su come un processo intende accedere ad un segmento
di memoria, anche al di là delle mappature dei file, così che possano essere
adottate le opportune strategie di ottimizzazione. Il suo prototipo è:

\begin{funcproto}{
\fhead{sys/mman.h}
\fdecl{int madvise(void *start, size\_t length, int advice)}
\fdesc{Fornisce indicazioni sull'uso previsto di un segmento di memoria.} 
}

{La funzione ritorna $0$ in caso di successo e $-1$ per un errore, nel qual
  caso \var{errno} assumerà uno dei valori: 
  \begin{errlist}
    \item[\errcode{EBADF}] la mappatura esiste ma non corrisponde ad un file.
    \item[\errcode{EINVAL}] \param{start} non è allineato alla dimensione di
      una pagina, \param{length} ha un valore negativo, o \param{advice} non è
      un valore valido, o si è richiesto il rilascio (con
      \const{MADV\_DONTNEED}) di pagine bloccate o condivise o si è usato
      \const{MADV\_MERGEABLE} o \const{MADV\_UNMERGEABLE} ma il kernel non è
      stato compilato per il relativo supporto.
    \item[\errcode{EIO}] la paginazione richiesta eccederebbe i limiti (vedi
      sez.~\ref{sec:sys_resource_limit}) sulle pagine residenti in memoria del
      processo (solo in caso di \const{MADV\_WILLNEED}).
    \item[\errcode{ENOMEM}] gli indirizzi specificati non sono mappati, o, in
      caso \const{MADV\_WILLNEED}, non c'è sufficiente memoria per soddisfare
      la richiesta.
  \end{errlist}
  ed inoltre \errval{EAGAIN} e \errval{ENOSYS} nel loro significato generico.}
\end{funcproto}

La sezione di memoria sulla quale si intendono fornire le indicazioni deve
essere indicata con l'indirizzo iniziale \param{start} e l'estensione
\param{length}, il valore di \param{start} deve essere allineato,
mentre \param{length} deve essere un numero positivo; la versione di Linux
consente anche un valore nullo per \param{length}, inoltre se una parte
dell'intervallo non è mappato in memoria l'indicazione viene comunque
applicata alle restanti parti, anche se la funzione ritorna un errore di
\errval{ENOMEM}.

L'indicazione viene espressa dall'argomento \param{advice} che deve essere
specificato con uno dei valori riportati in
tab.~\ref{tab:madvise_advice_values}; si tenga presente che i valori indicati
nella seconda parte della tabella sono specifici di Linux e non sono previsti
dallo standard POSIX.1b.  La funzione non ha, tranne il caso di
\const{MADV\_DONTFORK}, nessun effetto sul comportamento di un programma, ma
può influenzarne le prestazioni fornendo al kernel indicazioni sulle esigenze
dello stesso, così che sia possibile scegliere le opportune strategie per la
gestione del \textit{read-ahead} (vedi sez.~\ref{sec:file_fadvise}) e del
caching dei dati.

\begin{table}[!htb]
  \centering
  \footnotesize
  \begin{tabular}[c]{|l|p{10 cm}|}
    \hline
    \textbf{Valore} & \textbf{Significato} \\
    \hline
    \hline
    \constd{MADV\_DONTNEED}& non ci si aspetta nessun accesso nell'immediato
                             futuro, pertanto le pagine possono essere
                             liberate dal kernel non appena necessario; l'area
                             di memoria resterà accessibile, ma un accesso
                             richiederà che i dati vengano ricaricati dal file
                             a cui la mappatura fa riferimento.\\
    \constd{MADV\_NORMAL}  & nessuna indicazione specifica, questo è il valore
                             di default usato quando non si è chiamato
                             \func{madvise}.\\
    \constd{MADV\_RANDOM}  & ci si aspetta un accesso casuale all'area
                             indicata, pertanto l'applicazione di una lettura
                             anticipata con il meccanismo del
                             \textit{read-ahead} (vedi 
                             sez.~\ref{sec:file_fadvise}) è di
                             scarsa utilità e verrà disabilitata.\\
    \constd{MADV\_SEQUENTIAL}& ci si aspetta un accesso sequenziale al file,
                               quindi da una parte sarà opportuno eseguire una
                               lettura anticipata, e dall'altra si potranno
                               scartare immediatamente le pagine una volta che
                               queste siano state lette.\\
    \const{MADV\_WILLNEED}& ci si aspetta un accesso nell'immediato futuro,
                            pertanto l'applicazione del \textit{read-ahead}
                            deve essere incentivata.\\
    \hline
    \constd{MADV\_DONTDUMP}& esclude da un \textit{core dump} (vedi
                             sez.~\ref{sec:sig_standard}) le pagine 
                             specificate, viene usato per evitare di scrivere
                             su disco dati relativi a zone di memoria che si sa
                             non essere utili in un \textit{core dump}.\\
    \constd{MADV\_DODUMP}  & rimuove l'effetto della precedente
                             \const{MADV\_DONTDUMP} (dal kernel 3.4).\\ 
    \constd{MADV\_DONTFORK}& impedisce che l'intervallo specificato venga
                             ereditato dal processo figlio dopo una
                             \func{fork}; questo consente di evitare che il
                             meccanismo del \textit{copy on write} effettui la
                             rilocazione delle pagine quando il padre scrive
                             sull'area di memoria dopo la \func{fork}, cosa che
                             può causare problemi per l'hardware che esegue
                             operazioni in DMA su quelle pagine (dal kernel
                             2.6.16).\\
    \constd{MADV\_DOFORK}  & rimuove l'effetto della precedente
                             \const{MADV\_DONTFORK} (dal kernel 2.6.16).\\ 
    \constd{MADV\_HUGEPAGE}& abilita il meccanismo delle \textit{Transparent
                             Huge Page} (vedi sez.~\ref{sec:huge_pages})
                             sulla regione indicata; se questa è allineata
                             alle relative dimensioni il kernel alloca
                             direttamente delle \textit{huge page}; è
                             utilizzabile solo con mappature anomime private
                             (dal kernel 2.6.38).\\
    \constd{MADV\_NOHUGEPAGE}& impedisce che la regione indicata venga
                               collassata in eventuali \textit{huge page} (dal
                               kernel 2.6.38).\\
    \constd{MADV\_HWPOISON} &opzione ad uso di debug per verificare codice
                              che debba gestire errori nella gestione della
                              memoria; richiede una apposita opzione di
                              compilazione del kernel, privilegi amministrativi
                              (la capacità \const{CAP\_SYS\_ADMIN}) e provoca
                              l'emissione di un segnale di \const{SIGBUS} dal
                              programma chiamante e rimozione della mappatura
                              (dal kernel 2.6.32).\\
    \constd{MADV\_SOFT\_OFFLINE}&opzione utilizzata per il debug del
                              codice di verifica degli errori di gestione
                              memoria, richiede una apposita opzione di
                              compilazione (dal kernel 2.6.33).\\
    \constd{MADV\_MERGEABLE}& marca la pagina come accorpabile, indicazione
                              principalmente ad uso dei sistemi di
                              virtualizzazione\footnotemark (dal kernel
                              2.6.32).\\ 
    \constd{MADV\_REMOVE}  & libera un intervallo di pagine di memoria ed il
                             relativo supporto sottostante; è supportato
                             soltanto sui filesystem in RAM \textit{tmpfs} e
                             \textit{shmfs} se usato su altri tipi di
                             filesystem causa un errore di \errcode{ENOSYS}
                             (dal kernel 2.6.16).\\
    \constd{MADV\_UNMERGEABLE}& rimuove l'effetto della precedente
                                \const{MADV\_MERGEABLE} (dal kernel 2.6.32). \\
    \hline
  \end{tabular}
  \caption{Valori dell'argomento \param{advice} di \func{madvise}.}
  \label{tab:madvise_advice_values}
\end{table}

% TODO aggiunta MADV_FREE dal kernel 4.5 (vedi http://lwn.net/Articles/590991/)
% TODO aggiunta MADV_WIPEONFORK dal kernel 4.14 that causes the affected memory
% region to appear to be full of zeros in the child process after a fork. It
% differs from the existing MADV_DONTFORK in that the address range will
% remain valid in the child (dalla notizia in https://lwn.net/Articles/733256/).  

\footnotetext{a partire dal kernel 2.6.32 è stato introdotto un meccanismo che
  identifica pagine di memoria identiche e le accorpa in una unica pagina
  (soggetta al \textit{copy-on-write} per successive modifiche); per evitare
  di controllare tutte le pagine solo quelle marcate con questo flag vengono
  prese in considerazione per l'accorpamento; in questo modo si possono
  migliorare le prestazioni nella gestione delle macchine virtuali diminuendo
  la loro occupazione di memoria, ma il meccanismo può essere usato anche in
  altre applicazioni in cui sian presenti numerosi processi che usano gli
  stessi dati; per maggiori dettagli si veda
  \href{http://kernelnewbies.org/Linux_2_6_32\#head-d3f32e41df508090810388a57efce73f52660ccb}{\texttt{http://kernelnewbies.org/Linux\_2\_6\_32}}
  e la documentazione nei sorgenti del kernel
  (\texttt{Documentation/vm/ksm.txt}).} 


A differenza da quanto specificato nello standard POSIX.1b, per il quale l'uso
di \func{madvise} è a scopo puramente indicativo, Linux considera queste
richieste come imperative, per cui ritorna un errore qualora non possa
soddisfarle; questo comportamento differisce da quanto specificato nello
standard.

Nello standard POSIX.1-2001 è prevista una ulteriore funzione
\funcd{posix\_madvise} che su Linux viene reimplementata utilizzando
\func{madvise}; il suo prototipo è:

\begin{funcproto}{
\fhead{sys/mman.h}
\fdecl{int posix\_madvise(void *start, size\_t lenght, int advice)}
\fdesc{Fornisce indicazioni sull'uso previsto di un segmento di memoria.} 
}

{La funzione ritorna $0$ in caso di successo ed un valore positivo per un
  errore, nel qual caso \var{errno} assumerà uno dei valori:
  \begin{errlist}
    \item[\errcode{EINVAL}] \param{start} non è allineato alla dimensione di
      una pagina, \param{length} ha un valore negativo, o \param{advice} non è
      un valore valido.
    \item[\errcode{ENOMEM}] gli indirizzi specificati non sono nello spazio di
      indirizzi del processo.
  \end{errlist}
}
\end{funcproto}

Gli argomenti \param{start} e \param{lenght} hanno lo stesso identico
significato degli analoghi di \func{madvise}, a cui si rimanda per la loro
descrizione ma a differenza di quanto indicato dallo standard per questa
funzione, su Linux un valore nullo di \param{len} è consentito.

\begin{table}[!htb]
  \centering
  \footnotesize
  \begin{tabular}[c]{|l|l|}
    \hline
    \textbf{Valore} & \textbf{Significato} \\
    \hline
    \hline
    \constd{POSIX\_MADV\_DONTNEED}& analogo a \const{MADV\_DONTNEED}.\\
    \constd{POSIX\_MADV\_NORMAL}  & identico a \const{MADV\_NORMAL}.\\
    \constd{POSIX\_MADV\_RANDOM}  & identico a \const{MADV\_RANDOM}.\\
    \constd{POSIX\_MADV\_SEQUENTIAL}& identico a \const{MADV\_SEQUENTIAL}.\\
    \constd{POSIX\_MADV\_WILLNEED}& identico a \const{MADV\_WILLNEED}.\\
     \hline
  \end{tabular}
  \caption{Valori dell'argomento \param{advice} di \func{posix\_madvise}.}
  \label{tab:posix_madvise_advice_values}
\end{table}


L'argomento \param{advice} invece può assumere solo i valori indicati in
tab.~\ref{tab:posix_madvise_advice_values}, che riflettono gli analoghi di
\func{madvise}, con lo stesso effetto per tutti tranne
\const{POSIX\_MADV\_DONTNEED}.  Infatti a partire dalla \acr{glibc} 2.6
\const{POSIX\_MADV\_DONTNEED} viene ignorato, in quanto l'uso del
corrispondente \const{MADV\_DONTNEED} di \func{madvise} ha, per la semantica
imperativa, l'effetto immediato di far liberare le pagine da parte del kernel,
che viene considerato distruttivo.



\subsection{I/O vettorizzato: \func{readv} e \func{writev}}
\label{sec:file_multiple_io}

Una seconda modalità di I/O diversa da quella ordinaria è il cosiddetto
\textsl{I/O vettorizzato}, che nasce per rispondere al caso abbastanza comune
in cui ci si trova nell'esigenza di dover eseguire una serie multipla di
operazioni di I/O, come una serie di letture o scritture di vari buffer. Un
esempio tipico è quando i dati sono strutturati nei campi di una struttura ed
essi devono essere caricati o salvati su un file.  Benché l'operazione sia
facilmente eseguibile attraverso una serie multipla di chiamate a \func{read}
e \func{write}, ci sono casi in cui si vuole poter contare sulla atomicità
delle operazioni di lettura e scrittura rispetto all'esecuzione del programma.

Per questo motivo fino da BSD 4.2 vennero introdotte delle nuove
\textit{system call} che permettessero di effettuare con una sola chiamata una
serie di letture da, o scritture su, una serie di buffer, quello che poi venne
chiamato \textsl{I/O vettorizzato}. Queste funzioni di sistema sono
\funcd{readv} e \funcd{writev},\footnote{in Linux le due funzioni sono riprese
  da BSD4.4, esse sono previste anche dallo standard POSIX.1-2001.} ed i
relativi prototipi sono:


\begin{funcproto}{
\fhead{sys/uio.h}
\fdecl{int readv(int fd, const struct iovec *vector, int count)}
\fdecl{int writev(int fd, const struct iovec *vector, int count)}
\fdesc{Eseguono rispettivamente una lettura o una scrittura vettorizzata.} 
}

{Le funzioni ritornano il numero di byte letti o scritti in caso di successo e
  $-1$ per un errore, nel qual caso \var{errno} assumerà uno dei valori:
  \begin{errlist}
    \item[\errcode{EINVAL}] si è specificato un valore non valido per uno degli
    argomenti (ad esempio \param{count} è maggiore di \const{IOV\_MAX}).
  \end{errlist}
  più tutti i valori, con lo stesso significato, che possono risultare
  dalle condizioni di errore di \func{read} e \func{write}.
 }
\end{funcproto}


Entrambe le funzioni usano una struttura \struct{iovec}, la cui definizione è
riportata in fig.~\ref{fig:file_iovec}, che definisce dove i dati devono
essere letti o scritti ed in che quantità. Il primo campo della struttura,
\var{iov\_base}, contiene l'indirizzo del buffer ed il secondo,
\var{iov\_len}, la dimensione dello stesso.

\begin{figure}[!htb]
  \footnotesize \centering
  \begin{minipage}[c]{\textwidth}
    \includestruct{listati/iovec.h}
  \end{minipage} 
  \normalsize 
  \caption{La struttura \structd{iovec}, usata dalle operazioni di I/O
    vettorizzato.} 
  \label{fig:file_iovec}
\end{figure}

La lista dei buffer da utilizzare viene indicata attraverso l'argomento
\param{vector} che è un vettore di strutture \struct{iovec}, la cui lunghezza
è specificata dall'argomento \param{count}.\footnote{fino alle libc5, Linux
  usava \type{size\_t} come tipo dell'argomento \param{count}, una scelta
  logica, che però è stata dismessa per restare aderenti allo standard
  POSIX.1-2001.}  Ciascuna struttura dovrà essere inizializzata opportunamente
per indicare i vari buffer da e verso i quali verrà eseguito il trasferimento
dei dati. Essi verranno letti (o scritti) nell'ordine in cui li si sono
specificati nel vettore \param{vector}.

La standardizzazione delle due funzioni all'interno della revisione
POSIX.1-2001 prevede anche che sia possibile avere un limite al numero di
elementi del vettore \param{vector}. Qualora questo sussista, esso deve essere
indicato dal valore dalla costante \const{IOV\_MAX}, definita come le altre
costanti analoghe (vedi sez.~\ref{sec:sys_limits}) in \headfile{limits.h}; lo
stesso valore deve essere ottenibile in esecuzione tramite la funzione
\func{sysconf} richiedendo l'argomento \const{\_SC\_IOV\_MAX} (vedi
sez.~\ref{sec:sys_limits}).

Nel caso di Linux il limite di sistema è di 1024, però se si usa la
\acr{glibc} essa fornisce un \textit{wrapper} per le \textit{system call}
che si accorge se una operazione supererà il precedente limite, in tal caso i
dati verranno letti o scritti con le usuali \func{read} e \func{write} usando
un buffer di dimensioni sufficienti appositamente allocato in grado di
contenere tutti i dati indicati da \param{vector}. L'operazione avrà successo
ma si perderà l'atomicità del trasferimento da e verso la destinazione finale.

Si tenga presente infine che queste funzioni operano sui file con
l'interfaccia dei file descriptor, e non è consigliabile mescolarle con
l'interfaccia classica dei \textit{file stream} di
sez.~\ref{sec:files_std_interface}; a causa delle bufferizzazioni interne di
quest'ultima infatti si potrebbero avere risultati indefiniti e non
corrispondenti a quanto aspettato.

Come per le normali operazioni di lettura e scrittura, anche per l'\textsl{I/O
  vettorizzato} si pone il problema di poter effettuare le operazioni in
maniera atomica a partire da un certa posizione sul file. Per questo motivo a
partire dal kernel 2.6.30 sono state introdotte anche per l'\textsl{I/O
  vettorizzato} le analoghe delle funzioni \func{pread} e \func{pwrite} (vedi
sez.~\ref{sec:file_read} e \ref{sec:file_write}); le due funzioni sono
\funcd{preadv} e \funcd{pwritev} ed i rispettivi prototipi sono:\footnote{le
  due funzioni sono analoghe alle omonime presenti in BSD; le \textit{system
    call} usate da Linux (introdotte a partire dalla versione 2.6.30)
  utilizzano degli argomenti diversi per problemi collegati al formato a 64
  bit dell'argomento \param{offset}, che varia a seconda delle architetture,
  ma queste differenze vengono gestite dalle funzioni di librerie di libreria
  che mantengono l'interfaccia delle analoghe tratte da BSD.}


\begin{funcproto}{
\fhead{sys/uio.h}
\fdecl{int preadv(int fd, const struct iovec *vector, int count, off\_t
    offset)}
\fdecl{int pwritev(int fd, const struct iovec *vector, int count, off\_t
    offset)}
\fdesc{Eseguono una lettura o una scrittura vettorizzata a partire da una data
  posizione sul file.} 
}

{ Le funzioni hanno gli stessi valori di ritorno delle corrispondenti
  \func{readv} e \func{writev} ed anche gli eventuali errori sono gli stessi,
  con in più quelli che si possono ottenere dalle possibili condizioni di
  errore di \func{lseek}.
}
\end{funcproto}

Le due funzioni eseguono rispettivamente una lettura o una scrittura
vettorizzata a partire dalla posizione \param{offset} sul file indicato
da \param{fd}, la posizione corrente sul file, come vista da eventuali altri
processi che vi facciano riferimento, non viene alterata. A parte la presenza
dell'ulteriore argomento il comportamento delle funzioni è identico alle
precedenti \func{readv} e \func{writev}. 

Con l'uso di queste funzioni si possono evitare eventuali \textit{race
  condition} quando si deve eseguire la una operazione di lettura e scrittura
vettorizzata a partire da una certa posizione su un file, mentre al contempo
si possono avere in concorrenza processi che utilizzano lo stesso file
descriptor (si ricordi quanto visto in sez.~\ref{sec:file_adv_func}) con delle
chiamate a \func{lseek}.

% TODO trattare preadv2() e pwritev2(), introdotte con il kernel 4.6, vedi
% http://lwn.net/Articles/670231/ ed il flag RWF_HIPRI, anche l'aggiunta del
% flag RWF_APPEND a pwritev2 con il kernel 4.16, vedi
% https://lwn.net/Articles/746129/ 


\subsection{L'I/O diretto fra file descriptor: \func{sendfile} e
  \func{splice}} 
\label{sec:file_sendfile_splice}

Uno dei problemi che si presentano nella gestione dell'I/O è quello in cui si
devono trasferire grandi quantità di dati da un file descriptor ed un altro;
questo usualmente comporta la lettura dei dati dal primo file descriptor in un
buffer in memoria, da cui essi vengono poi scritti sul secondo.

Benché il kernel ottimizzi la gestione di questo processo quando si ha a che
fare con file normali, in generale quando i dati da trasferire sono molti si
pone il problema di effettuare trasferimenti di grandi quantità di dati da
\textit{kernel space} a \textit{user space} e all'indietro, quando in realtà
potrebbe essere più efficiente mantenere tutto in \textit{kernel
  space}. Tratteremo in questa sezione alcune funzioni specialistiche che
permettono di ottimizzare le prestazioni in questo tipo di situazioni.

La prima funzione che è stata ideata per ottimizzare il trasferimento dei dati
fra due file descriptor è \func{sendfile}.\footnote{la funzione è stata
  introdotta con i kernel della serie 2.2, e disponibile dalla \acr{glibc}
  2.1.} La funzione è presente in diverse versioni di Unix (la si ritrova ad
esempio in FreeBSD, HPUX ed altri Unix) ma non è presente né in POSIX.1-2001
né in altri standard (pertanto si eviti di utilizzarla se si devono scrivere
programmi portabili) per cui per essa vengono utilizzati prototipi e
semantiche differenti. Nel caso di Linux il prototipo di \funcd{sendfile} è:


\begin{funcproto}{
\fhead{sys/sendfile.h}
\fdecl{ssize\_t sendfile(int out\_fd, int in\_fd, off\_t *offset, size\_t
    count)}
\fdesc{Copia dei dati da un file descriptor ad un altro.} 
}

{La funzione ritorna il numero di byte trasferiti in caso di successo e $-1$
  per un errore, nel qual caso \var{errno} assumerà uno dei valori:
  \begin{errlist}
    \item[\errcode{EAGAIN}] si è impostata la modalità non bloccante su
      \param{out\_fd} e la scrittura si bloccherebbe.
    \item[\errcode{EINVAL}] i file descriptor non sono validi, o sono bloccati
      (vedi sez.~\ref{sec:file_locking}), o \func{mmap} non è disponibile per
      \param{in\_fd}.
    \item[\errcode{EIO}] si è avuto un errore di lettura da \param{in\_fd}.
    \item[\errcode{ENOMEM}] non c'è memoria sufficiente per la lettura da
      \param{in\_fd}.
  \end{errlist}
  ed inoltre \errcode{EBADF} e \errcode{EFAULT} nel loro significato
  generico.}
\end{funcproto}

La funzione copia direttamente \param{count} byte dal file descriptor
\param{in\_fd} al file descriptor \param{out\_fd}. In caso di successo la
funzione ritorna il numero di byte effettivamente copiati da \param{in\_fd} a
\param{out\_fd} e come per le ordinarie \func{read} e \func{write} questo
valore può essere inferiore a quanto richiesto con \param{count}.

Se il puntatore \param{offset} è nullo la funzione legge i dati a partire
dalla posizione corrente su \param{in\_fd}, altrimenti verrà usata la
posizione indicata dal valore puntato da \param{offset}; in questo caso detto
valore sarà aggiornato, come \textit{value result argument}, per indicare la
posizione del byte successivo all'ultimo che è stato letto, mentre la
posizione corrente sul file non sarà modificata. Se invece \param{offset} è
nullo la posizione corrente sul file sarà aggiornata tenendo conto dei byte
letti da \param{in\_fd}.

Fino ai kernel della serie 2.4 la funzione era utilizzabile su un qualunque
file descriptor, e permetteva di sostituire la invocazione successiva di una
\func{read} e una \func{write} (e l'allocazione del relativo buffer) con una
sola chiamata a \funcd{sendfile}. In questo modo si poteva diminuire il numero
di chiamate al sistema e risparmiare in trasferimenti di dati da
\textit{kernel space} a \textit{user space} e viceversa.  La massima utilità
della funzione si ottiene comunque per il trasferimento di dati da un file su
disco ad un socket di rete,\footnote{questo è il caso classico del lavoro
  eseguito da un server web, ed infatti Apache ha una opzione per il supporto
  esplicito di questa funzione.} dato che in questo caso diventa possibile
effettuare il trasferimento diretto via DMA dal controller del disco alla
scheda di rete, senza neanche allocare un buffer nel kernel (il meccanismo è
detto \textit{zerocopy} in quanto i dati non vengono mai copiati dal kernel,
che si limita a programmare solo le operazioni di lettura e scrittura via DMA)
ottenendo la massima efficienza possibile senza pesare neanche sul processore.

In seguito però ci si accorse che, fatta eccezione per il trasferimento
diretto da file a socket, non sempre \func{sendfile} comportava miglioramenti
significativi delle prestazioni rispetto all'uso in sequenza di \func{read} e
\func{write}. Nel caso generico infatti il kernel deve comunque allocare un
buffer ed effettuare la copia dei dati, e in tal caso spesso il guadagno
ottenibile nel ridurre il numero di chiamate al sistema non compensa le
ottimizzazioni che possono essere fatte da una applicazione in \textit{user
  space} che ha una conoscenza diretta su come questi sono strutturati, per
cui in certi casi si potevano avere anche dei peggioramenti.  Questo ha
portato, per i kernel della serie 2.6,\footnote{per alcune motivazioni di
  questa scelta si può fare riferimento a quanto illustrato da Linus Torvalds
  in \url{http://www.cs.helsinki.fi/linux/linux-kernel/2001-03/0200.html}.}
alla decisione di consentire l'uso della funzione soltanto quando il file da
cui si legge supporta le operazioni di \textit{memory mapping} (vale a dire
non è un socket) e quello su cui si scrive è un socket; in tutti gli altri
casi l'uso di \func{sendfile} da luogo ad un errore di \errcode{EINVAL}.

Nonostante ci possano essere casi in cui \func{sendfile} non migliora le
prestazioni, resta il dubbio se la scelta di disabilitarla sempre per il
trasferimento fra file di dati sia davvero corretta. Se ci sono peggioramenti
di prestazioni infatti si può sempre fare ricorso al metodo ordinario, ma
lasciare a disposizione la funzione consentirebbe se non altro di semplificare
la gestione della copia dei dati fra file, evitando di dover gestire
l'allocazione di un buffer temporaneo per il loro trasferimento. Comunque a
partire dal kernel 2.6.33 la restrizione su \param{out\_fd} è stata rimossa e
questo può essere un file qualunque, rimane però quella di non poter usare un
socket per \param{in\_fd}.

A partire dal kernel 2.6.17 come alternativa a \func{sendfile} è disponibile
la nuova \textit{system call} \func{splice}. Lo scopo di questa funzione è
quello di fornire un meccanismo generico per il trasferimento di dati da o
verso un file, utilizzando un buffer gestito internamente dal
kernel. Descritta in questi termini \func{splice} sembra semplicemente un
``\textsl{dimezzamento}'' di \func{sendfile}, nel senso che un trasferimento
di dati fra due file con \func{sendfile} non sarebbe altro che la lettura
degli stessi su un buffer seguita dalla relativa scrittura, cosa che in questo
caso si dovrebbe eseguire con due chiamate a \func{splice}.

In realtà le due \textit{system call} sono profondamente diverse nel loro
meccanismo di funzionamento;\footnote{questo fino al kernel 2.6.23, dove
  \func{sendfile} è stata reimplementata in termini di \func{splice}, pur
  mantenendo disponibile la stessa interfaccia verso l'\textit{user space}.}
\func{sendfile} infatti, come accennato, non necessita di avere a disposizione
un buffer interno, perché esegue un trasferimento diretto di dati; questo la
rende in generale più efficiente, ma anche limitata nelle sue applicazioni,
dato che questo tipo di trasferimento è possibile solo in casi specifici che
nel caso di Linux questi sono anche solo quelli in cui essa può essere
effettivamente utilizzata.

Il concetto che sta dietro a \func{splice} invece è diverso,\footnote{in
  realtà la proposta originale di Larry Mc Voy non differisce poi tanto negli
  scopi da \func{sendfile}, quello che rende \func{splice} davvero diversa è
  stata la reinterpretazione che ne è stata fatta nell'implementazione su
  Linux realizzata da Jens Anxboe, concetti che sono esposti sinteticamente
  dallo stesso Linus Torvalds in \url{http://kerneltrap.org/node/6505}.} si
tratta semplicemente di una funzione che consente di fare in maniera del tutto
generica delle operazioni di trasferimento di dati fra un file e un buffer
gestito interamente in \textit{kernel space}. In questo caso il cuore della
funzione (e delle affini \func{vmsplice} e \func{tee}, che tratteremo più
avanti) è appunto l'uso di un buffer in \textit{kernel space}, e questo è
anche quello che ne ha semplificato l'adozione, perché l'infrastruttura per la
gestione di un tale buffer è presente fin dagli albori di Unix per la
realizzazione delle \textit{pipe} (vedi sez.~\ref{sec:ipc_unix}). Dal punto di
vista concettuale allora \func{splice} non è altro che una diversa interfaccia
(rispetto alle \textit{pipe}) con cui utilizzare in \textit{user space}
l'oggetto ``\textsl{buffer in kernel space}''.

Così se per una \textit{pipe} o una \textit{fifo} il buffer viene utilizzato
come area di memoria (vedi fig.~\ref{fig:ipc_pipe_singular}) dove appoggiare i
dati che vengono trasferiti da un capo all'altro della stessa per creare un
meccanismo di comunicazione fra processi, nel caso di \func{splice} il buffer
viene usato o come fonte dei dati che saranno scritti su un file, o come
destinazione dei dati che vengono letti da un file. La funzione fornisce
quindi una interfaccia generica che consente di trasferire dati da un buffer
ad un file o viceversa; il prototipo di \funcd{splice}, accessibile solo dopo
aver definito la macro \macro{\_GNU\_SOURCE},\footnote{si ricordi che questa
  funzione non è contemplata da nessuno standard, è presente solo su Linux, e
  pertanto deve essere evitata se si vogliono scrivere programmi portabili.}
è il seguente:

\begin{funcproto}{
\fhead{fcntl.h} 
\fdecl{long splice(int fd\_in, off\_t *off\_in, int fd\_out, off\_t
    *off\_out, size\_t len, \\
\phantom{long splice(}unsigned int flags)}
\fdesc{Trasferisce dati da un file verso una \textit{pipe} o viceversa.} 
}

{La funzione ritorna il numero di byte trasferiti in caso di successo e $-1$
  per un errore, nel qual caso \var{errno} assumerà uno dei valori:
  \begin{errlist}
    \item[\errcode{EBADF}] uno o entrambi fra \param{fd\_in} e \param{fd\_out}
      non sono file descriptor validi o, rispettivamente, non sono stati
      aperti in lettura o scrittura.
    \item[\errcode{EINVAL}] il filesystem su cui si opera non supporta
      \func{splice}, oppure nessuno dei file descriptor è una \textit{pipe},
      oppure si 
      è dato un valore a \param{off\_in} o \param{off\_out} ma il
      corrispondente file è un dispositivo che non supporta la funzione
      \func{lseek}.
    \item[\errcode{ENOMEM}] non c'è memoria sufficiente per l'operazione
      richiesta.
    \item[\errcode{ESPIPE}] o \param{off\_in} o \param{off\_out} non sono
      \val{NULL} ma il corrispondente file descriptor è una \textit{pipe}.
  \end{errlist}
}
\end{funcproto}


La funzione esegue un trasferimento di \param{len} byte dal file descriptor
\param{fd\_in} al file descriptor \param{fd\_out}, uno dei quali deve essere
una \textit{pipe}; l'altro file descriptor può essere qualunque, questo
significa che può essere, oltre che un file di dati, anche un altra
\textit{pipe}, o un socket.  Come accennato una \textit{pipe} non è altro che
un buffer in \textit{kernel space}, per cui a seconda che essa sia usata
per \param{fd\_in} o \param{fd\_out} si avrà rispettivamente la copia dei dati
dal buffer al file o viceversa.

In caso di successo la funzione ritorna il numero di byte trasferiti, che può
essere, come per le normali funzioni di lettura e scrittura su file, inferiore
a quelli richiesti; un valore negativo indicherà un errore mentre un valore
nullo indicherà che non ci sono dati da trasferire (ad esempio si è giunti
alla fine del file in lettura). Si tenga presente che, a seconda del verso del
trasferimento dei dati, la funzione si comporta nei confronti del file
descriptor che fa riferimento al file ordinario, come \func{read} o
\func{write}, e pertanto potrà anche bloccarsi (a meno che non si sia aperto
il suddetto file in modalità non bloccante).

I due argomenti \param{off\_in} e \param{off\_out} consentono di specificare,
come per l'analogo \param{offset} di \func{sendfile}, la posizione all'interno
del file da cui partire per il trasferimento dei dati. Come per
\func{sendfile} un valore nullo indica di usare la posizione corrente sul
file, ed essa sarà aggiornata automaticamente secondo il numero di byte
trasferiti. Un valore non nullo invece deve essere un puntatore ad una
variabile intera che indica la posizione da usare; questa verrà aggiornata, al
ritorno della funzione, al byte successivo all'ultimo byte trasferito.
Ovviamente soltanto uno di questi due argomenti, e più precisamente quello che
fa riferimento al file descriptor non associato alla \textit{pipe}, può essere
specificato come valore non nullo.

Infine l'argomento \param{flags} consente di controllare alcune
caratteristiche del funzionamento della funzione; il contenuto è una maschera
binaria e deve essere specificato come OR aritmetico dei valori riportati in
tab.~\ref{tab:splice_flag}. Alcuni di questi valori vengono utilizzati anche
dalle funzioni \func{vmsplice} e \func{tee} per cui la tabella riporta le
descrizioni complete di tutti i valori possibili anche quando, come per
\const{SPLICE\_F\_GIFT}, questi non hanno effetto su \func{splice}.

\begin{table}[htb]
  \centering
  \footnotesize
  \begin{tabular}[c]{|l|p{10cm}|}
    \hline
    \textbf{Valore} & \textbf{Significato} \\
    \hline
    \hline
    \constd{SPLICE\_F\_MOVE} & Suggerisce al kernel di spostare le pagine
                               di memoria contenenti i dati invece di
                               copiarle: per una maggiore efficienza
                               \func{splice} usa quando possibile i
                               meccanismi della memoria virtuale per
                               eseguire i trasferimenti di dati. In maniera
                               analoga a \func{mmap}), qualora le pagine non
                               possano essere spostate dalla \textit{pipe} o
                               il buffer non corrisponda a pagine intere
                               esse saranno comunque copiate. Viene usato
                               soltanto da \func{splice}.\\ 
    \constd{SPLICE\_F\_NONBLOCK}& Richiede di operare in modalità non
                                  bloccante; questo flag influisce solo sulle
                                  operazioni che riguardano l'I/O da e verso la
                                  \textit{pipe}. Nel caso di \func{splice}
                                  questo significa che la funzione potrà
                                  comunque bloccarsi nell'accesso agli altri
                                  file descriptor (a meno che anch'essi non
                                  siano stati aperti in modalità non
                                  bloccante).\\
    \constd{SPLICE\_F\_MORE} & Indica al kernel che ci sarà l'invio di
                               ulteriori dati in una \func{splice}
                               successiva, questo è un suggerimento utile
                               che viene usato quando \param{fd\_out} è un
                               socket. Questa opzione consente di utilizzare
                               delle opzioni di gestione dei socket che
                               permettono di ottimizzare le trasmissioni via
                               rete (si veda la descrizione di
                               \const{TCP\_CORK} in
                               sez.~\ref{sec:sock_tcp_udp_options} e quella
                               di \const{MSG\_MORE} in
                               sez.~\ref{sec:net_sendmsg}).  Attualmente
                               viene usato solo da \func{splice}, potrà essere
                               implementato in futuro anche per
                               \func{vmsplice} e \func{tee}.\\
    \constd{SPLICE\_F\_GIFT} & Le pagine di memoria utente sono
                               ``\textsl{donate}'' al kernel; questo
                               significa che la cache delle pagine e i dati
                               su disco potranno differire, e che
                               l'applicazione non potrà modificare
                               quest'area di memoria. 
                               Se impostato una seguente \func{splice} che
                               usa \const{SPLICE\_F\_MOVE} potrà spostare le 
                               pagine con successo, altrimenti esse dovranno
                               essere copiate; per usare questa opzione i
                               dati dovranno essere opportunamente allineati
                               in posizione ed in dimensione alle pagine di
                               memoria. Viene usato soltanto da
                               \func{vmsplice}.\\
    \hline
  \end{tabular}
  \caption{Le costanti che identificano i bit della maschera binaria
    dell'argomento \param{flags} di \func{splice}, \func{vmsplice} e
    \func{tee}.} 
  \label{tab:splice_flag}
\end{table}


Per capire meglio il funzionamento di \func{splice} vediamo un esempio con un
semplice programma che usa questa funzione per effettuare la copia di un file
su un altro senza utilizzare buffer in \textit{user space}. Lo scopo del
programma è quello di eseguire la copia dei dati con \func{splice}, questo
significa che si dovrà usare la funzione due volte, prima per leggere i dati
dal file di ingresso e poi per scriverli su quello di uscita, appoggiandosi ad
una \textit{pipe}: lo schema del flusso dei dati è illustrato in
fig.~\ref{fig:splicecp_data_flux}.

\begin{figure}[htb]
  \centering
  \includegraphics[height=3.5cm]{img/splice_copy}
  \caption{Struttura del flusso di dati usato dal programma \texttt{splicecp}.}
  \label{fig:splicecp_data_flux}
\end{figure}

Il programma si chiama \texttt{splicecp.c} ed il codice completo è disponibile
coi sorgenti allegati alla guida, il corpo principale del programma, che non
contiene la sezione di gestione delle opzioni, le funzioni di ausilio, le
aperture dei file di ingresso e di uscita passati come argomenti e quella
della \textit{pipe} intermedia, è riportato in fig.~\ref{fig:splice_example}.

\begin{figure}[!htb]
  \footnotesize \centering
  \begin{minipage}[c]{\codesamplewidth}
    \includecodesample{listati/splicecp.c}
  \end{minipage}
  \normalsize
  \caption{Esempio di codice che usa \func{splice} per effettuare la copia di
    un file.}
  \label{fig:splice_example}
\end{figure}

Il ciclo principale (\texttt{\small 13-38}) inizia con la lettura dal file
sorgente tramite la prima \func{splice} (\texttt{\small 14-15}), in questo
caso si è usato come primo argomento il file descriptor del file sorgente e
come terzo quello del capo in scrittura della \textit{pipe}. Il funzionamento
delle \textit{pipe} e l'uso della coppia di file descriptor ad esse associati
è trattato in dettaglio in sez.~\ref{sec:ipc_unix}; non ne parleremo qui dato
che nell'ottica dell'uso di \func{splice} questa operazione corrisponde
semplicemente al trasferimento dei dati dal file al buffer in \textit{kernel
  space}.

La lettura viene eseguita in blocchi pari alla dimensione specificata
dall'opzione \texttt{-s} (il default è 4096); essendo in questo caso
\func{splice} equivalente ad una \func{read} sul file, se ne controlla il
valore di uscita in \var{nread} che indica quanti byte sono stati letti, se
detto valore è nullo (\texttt{\small 16}) questo significa che si è giunti
alla fine del file sorgente e pertanto l'operazione di copia è conclusa e si
può uscire dal ciclo arrivando alla conclusione del programma (\texttt{\small
  59}). In caso di valore negativo (\texttt{\small 17-24}) c'è stato un
errore ed allora si ripete la lettura (\texttt{\small 16}) se questo è dovuto
ad una interruzione, o altrimenti si esce con un messaggio di errore
(\texttt{\small 21-23}).

Una volta completata con successo la lettura si avvia il ciclo di scrittura
(\texttt{\small 25-37}); questo inizia (\texttt{\small 26-27}) con la
seconda \func{splice} che cerca di scrivere gli \var{nread} byte letti, si
noti come in questo caso il primo argomento faccia di nuovo riferimento alla
\textit{pipe} (in questo caso si usa il capo in lettura, per i dettagli si
veda al solito sez.~\ref{sec:ipc_unix}) mentre il terzo sia il file descriptor
del file di destinazione.

Di nuovo si controlla il numero di byte effettivamente scritti restituito in
\var{nwrite} e in caso di errore al solito si ripete la scrittura se questo è
dovuto a una interruzione o si esce con un messaggio negli altri casi
(\texttt{\small 28-35}). Infine si chiude il ciclo di scrittura sottraendo
(\texttt{\small 37}) il numero di byte scritti a quelli di cui è richiesta la
scrittura,\footnote{in questa parte del ciclo \var{nread}, il cui valore
  iniziale è dato dai byte letti dalla precedente chiamata a \func{splice},
  viene ad assumere il significato di byte da scrivere.} così che il ciclo di
scrittura venga ripetuto fintanto che il valore risultante sia maggiore di
zero, indice che la chiamata a \func{splice} non ha esaurito tutti i dati
presenti sul buffer.

Si noti come il programma sia concettualmente identico a quello che si sarebbe
scritto usando \func{read} al posto della prima \func{splice} e \func{write}
al posto della seconda, utilizzando un buffer in \textit{user space} per
eseguire la copia dei dati, solo che in questo caso non è stato necessario
allocare nessun buffer e non si è trasferito nessun dato in \textit{user
  space}.  Si noti anche come si sia usata la combinazione
\texttt{SPLICE\_F\_MOVE | SPLICE\_F\_MORE } per l'argomento \param{flags} di
\func{splice}, infatti anche se un valore nullo avrebbe dato gli stessi
risultati, l'uso di questi flag, che si ricordi servono solo a dare
suggerimenti al kernel, permette in genere di migliorare le prestazioni.

Come accennato con l'introduzione di \func{splice} sono state realizzate anche
altre due \textit{system call}, \func{vmsplice} e \func{tee}, che utilizzano
la stessa infrastruttura e si basano sullo stesso concetto di manipolazione e
trasferimento di dati attraverso un buffer in \textit{kernel space}; benché
queste non attengono strettamente ad operazioni di trasferimento dati fra file
descriptor, le tratteremo qui, essendo strettamente correlate fra loro.

La prima funzione, \funcd{vmsplice}, è la più simile a \func{splice} e come
indica il suo nome consente di trasferire i dati dalla memoria virtuale di un
processo (ad esempio per un file mappato in memoria) verso una \textit{pipe};
il suo prototipo è:

\begin{funcproto}{
\fhead{fcntl.h} 
\fhead{sys/uio.h}
\fdecl{long vmsplice(int fd, const struct iovec *iov, unsigned long nr\_segs,\\
\phantom{long vmsplice(}unsigned int flags)}
\fdesc{Trasferisce dati dalla memoria di un processo verso una \textit{pipe}.} 
}

{La funzione ritorna il numero di byte trasferiti in caso di successo e $-1$
  per un errore, nel qual caso \var{errno} assumerà uno dei valori:
  \begin{errlist}
    \item[\errcode{EBADF}] o \param{fd} non è un file descriptor valido o non
      fa riferimento ad una \textit{pipe}.
    \item[\errcode{EINVAL}] si è usato un valore nullo per \param{nr\_segs}
      oppure si è usato \const{SPLICE\_F\_GIFT} ma la memoria non è allineata.
    \item[\errcode{ENOMEM}] non c'è memoria sufficiente per l'operazione
      richiesta.
  \end{errlist}
}
\end{funcproto}

La \textit{pipe} indicata da \param{fd} dovrà essere specificata tramite il
file descriptor corrispondente al suo capo aperto in scrittura (di nuovo si
faccia riferimento a sez.~\ref{sec:ipc_unix}), mentre per indicare quali
segmenti della memoria del processo devono essere trasferiti verso di essa si
dovrà utilizzare un vettore di strutture \struct{iovec} (vedi
fig.~\ref{fig:file_iovec}), esattamente con gli stessi criteri con cui le si
usano per l'I/O vettorizzato, indicando gli indirizzi e le dimensioni di
ciascun segmento di memoria su cui si vuole operare; le dimensioni del
suddetto vettore devono essere passate nell'argomento \param{nr\_segs} che
indica il numero di segmenti di memoria da trasferire.  Sia per il vettore che
per il valore massimo di \param{nr\_segs} valgono le stesse limitazioni
illustrate in sez.~\ref{sec:file_multiple_io}.

In caso di successo la funzione ritorna il numero di byte trasferiti sulla
\textit{pipe}. In generale, se i dati una volta creati non devono essere
riutilizzati (se cioè l'applicazione che chiama \func{vmsplice} non
modificherà più la memoria trasferita), è opportuno utilizzare
per \param{flag} il valore \const{SPLICE\_F\_GIFT}; questo fa sì che il kernel
possa rimuovere le relative pagine dalla cache della memoria virtuale, così
che queste possono essere utilizzate immediatamente senza necessità di
eseguire una copia dei dati che contengono.

La seconda funzione aggiunta insieme a \func{splice} è \func{tee}, che deve il
suo nome all'omonimo comando in \textit{user space}, perché in analogia con
questo permette di duplicare i dati in ingresso su una \textit{pipe} su
un'altra \textit{pipe}. In sostanza, sempre nell'ottica della manipolazione
dei dati su dei buffer in \textit{kernel space}, la funzione consente di
eseguire una copia del contenuto del buffer stesso. Il prototipo di
\funcd{tee} è il seguente:

\begin{funcproto}{
\fhead{fcntl.h}
\fdecl{long tee(int fd\_in, int fd\_out, size\_t len, unsigned int
    flags)}
\fdesc{Duplica i dati da una \textit{pipe} ad un'altra.} 
}

{La funzione ritorna restituisce il numero di byte copiati in caso di successo
  e $-1$ per un errore, nel qual caso \var{errno} assumerà uno dei valori:
  \begin{errlist}
    \item[\errcode{EINVAL}] o uno fra \param{fd\_in} e \param{fd\_out} non fa
      riferimento ad una \textit{pipe} o entrambi fanno riferimento alla
      stessa \textit{pipe}.
    \item[\errcode{ENOMEM}] non c'è memoria sufficiente per l'operazione
      richiesta.
  \end{errlist}
}
\end{funcproto}

La funzione copia \param{len} byte del contenuto di una \textit{pipe} su di
un'altra; \param{fd\_in} deve essere il capo in lettura della \textit{pipe}
sorgente e \param{fd\_out} il capo in scrittura della \textit{pipe}
destinazione; a differenza di quanto avviene con \func{read} i dati letti con
\func{tee} da \param{fd\_in} non vengono \textsl{consumati} e restano
disponibili sulla \textit{pipe} per una successiva lettura (di nuovo per il
comportamento delle \textit{pipe} si veda sez.~\ref{sec:ipc_unix}). Al
momento\footnote{quello della stesura di questo paragrafo, avvenuta il Gennaio
  2010, in futuro potrebbe essere implementato anche \const{SPLICE\_F\_MORE}.}
il solo valore utilizzabile per \param{flag}, fra quelli elencati in
tab.~\ref{tab:splice_flag}, è \const{SPLICE\_F\_NONBLOCK} che rende la
funzione non bloccante.

La funzione restituisce il numero di byte copiati da una \textit{pipe}
all'altra (o $-1$ in caso di errore), un valore nullo indica che non ci sono
byte disponibili da copiare e che il capo in scrittura della \textit{pipe} è
stato chiuso; si tenga presente però che questo non avviene se si è impostato
il flag \const{SPLICE\_F\_NONBLOCK}, in tal caso infatti si avrebbe un errore
di \errcode{EAGAIN}. Un esempio di realizzazione del comando \texttt{tee}
usando questa funzione, ripreso da quello fornito nella pagina di manuale e
dall'esempio allegato al patch originale, è riportato in
fig.~\ref{fig:tee_example}. Il programma consente di copiare il contenuto
dello \textit{standard input} sullo \textit{standard output} e su un file
specificato come argomento, il codice completo si trova nel file
\texttt{tee.c} dei sorgenti allegati alla guida.

\begin{figure}[!htb]
  \footnotesize \centering
  \begin{minipage}[c]{\codesamplewidth}
    \includecodesample{listati/tee.c}
  \end{minipage}
  \normalsize
  \caption{Esempio di codice che usa \func{tee} per copiare i dati dello
    standard input sullo standard output e su un file.}
  \label{fig:tee_example}
\end{figure}

La prima parte del programma, che si è omessa per brevità, si cura
semplicemente di controllare che sia stato fornito almeno un argomento (il
nome del file su cui scrivere), di aprirlo e che sia lo standard input che lo
standard output corrispondano ad una \textit{pipe}.

Il ciclo principale (\texttt{\small 11-32}) inizia con la chiamata a
\func{tee} che duplica il contenuto dello standard input sullo standard output
(\texttt{\small 13}), questa parte è del tutto analoga ad una lettura ed
infatti come nell'esempio di fig.~\ref{fig:splice_example} si controlla il
valore di ritorno della funzione in \var{len}; se questo è nullo significa che
non ci sono più dati da leggere e si chiude il ciclo (\texttt{\small 14}), se
è negativo c'è stato un errore, ed allora si ripete la chiamata se questo è
dovuto ad una interruzione (\texttt{\small 15-48}) o si stampa un messaggio
di errore e si esce negli altri casi (\texttt{\small 18-21}).

Una volta completata la copia dei dati sullo \textit{standard output} si
possono estrarre dallo \textit{standard input} e scrivere sul file, di nuovo
su usa un ciclo di scrittura (\texttt{\small 24-31}) in cui si ripete una
chiamata a \func{splice} (\texttt{\small 25}) fintanto che non si sono scritti
tutti i \var{len} byte copiati in precedenza con \func{tee} (il funzionamento
è identico all'analogo ciclo di scrittura del precedente esempio di
fig.~\ref{fig:splice_example}).

Infine una nota finale riguardo \func{splice}, \func{vmsplice} e \func{tee}:
occorre sottolineare che benché finora si sia parlato di trasferimenti o copie
di dati in realtà nella implementazione di queste \textit{system call} non è
affatto detto che i dati vengono effettivamente spostati o copiati, il kernel
infatti realizza le \textit{pipe} come un insieme di puntatori\footnote{per
  essere precisi si tratta di un semplice buffer circolare, un buon articolo
  sul tema si trova su \url{http://lwn.net/Articles/118750/}.}  alle pagine di
memoria interna che contengono i dati, per questo una volta che i dati sono
presenti nella memoria del kernel tutto quello che viene fatto è creare i
suddetti puntatori ed aumentare il numero di referenze; questo significa che
anche con \func{tee} non viene mai copiato nessun byte, vengono semplicemente
copiati i puntatori.

% TODO?? dal 2.6.25 splice ha ottenuto il supporto per la ricezione su rete


% TODO trattare qui copy_file_range (vedi http://lwn.net/Articles/659523/),
% introdotta nel kernel 4.5

\subsection{Gestione avanzata dell'accesso ai dati dei file}
\label{sec:file_fadvise}

Nell'uso generico dell'interfaccia per l'accesso al contenuto dei file le
operazioni di lettura e scrittura non necessitano di nessun intervento di
supervisione da parte dei programmi, si eseguirà una \func{read} o una
\func{write}, i dati verranno passati al kernel che provvederà ad effettuare
tutte le operazioni (e a gestire il \textit{caching} dei dati) per portarle a
termine in quello che ritiene essere il modo più efficiente.

Il problema è che il concetto di migliore efficienza impiegato dal kernel è
relativo all'uso generico, mentre esistono molti casi in cui ci sono esigenze
specifiche dei singoli programmi, che avendo una conoscenza diretta di come
verranno usati i file, possono necessitare di effettuare delle ottimizzazioni
specifiche, relative alle proprie modalità di I/O sugli stessi. Tratteremo in
questa sezione una serie funzioni che consentono ai programmi di ottimizzare
il loro accesso ai dati dei file e controllare la gestione del relativo
\textit{caching}.

\itindbeg{read-ahead}

Una prima funzione che può essere utilizzata per modificare la gestione
ordinaria dell'I/O su un file è \funcd{readahead} (questa è una funzione
specifica di Linux, introdotta con il kernel 2.4.13, e non deve essere usata
se si vogliono scrivere programmi portabili), che consente di richiedere una
lettura anticipata del contenuto dello stesso in cache, così che le seguenti
operazioni di lettura non debbano subire il ritardo dovuto all'accesso al
disco; il suo prototipo è:

\begin{funcproto}{
\fhead{fcntl.h}
\fdecl{ssize\_t readahead(int fd, off64\_t *offset, size\_t count)}
\fdesc{Esegue una lettura preventiva del contenuto di un file in cache.} 
}

{La funzione ritorna $0$ in caso di successo e $-1$ per un errore, nel qual
  caso \var{errno} assumerà uno dei valori: 
  \begin{errlist}
    \item[\errcode{EBADF}] l'argomento \param{fd} non è un file descriptor
      valido o non è aperto in lettura.
    \item[\errcode{EINVAL}] l'argomento \param{fd} si riferisce ad un tipo di
      file che non supporta l'operazione (come una \textit{pipe} o un socket).
  \end{errlist}
}
\end{funcproto}

La funzione richiede che venga letto in anticipo il contenuto del file
\param{fd} a partire dalla posizione \param{offset} e per un ammontare di
\param{count} byte, in modo da portarlo in cache.  La funzione usa la memoria
virtuale ed il meccanismo della paginazione per cui la lettura viene eseguita
in blocchi corrispondenti alle dimensioni delle pagine di memoria, ed i valori
di \param{offset} e \param{count} vengono arrotondati di conseguenza.

La funzione estende quello che è un comportamento normale del kernel che,
quando si legge un file, aspettandosi che l'accesso prosegua, esegue sempre
una lettura preventiva di una certa quantità di dati; questo meccanismo di
lettura anticipata viene chiamato \textit{read-ahead}, da cui deriva il nome
della funzione. La funzione \func{readahead}, per ottimizzare gli accessi a
disco, effettua la lettura in cache della sezione richiesta e si blocca
fintanto che questa non viene completata.  La posizione corrente sul file non
viene modificata ed indipendentemente da quanto indicato con \param{count} la
lettura dei dati si interrompe una volta raggiunta la fine del file.

Si può utilizzare questa funzione per velocizzare le operazioni di lettura
all'interno di un programma tutte le volte che si conosce in anticipo quanti
dati saranno necessari nelle elaborazioni successive. Si potrà così
concentrare in un unico momento (ad esempio in fase di inizializzazione) la
lettura dei dati da disco, così da ottenere una migliore velocità di risposta
nelle operazioni successive.

\itindend{read-ahead}

Il concetto di \func{readahead} viene generalizzato nello standard
POSIX.1-2001 dalla funzione \func{posix\_fadvise} (anche se
l'argomento \param{len} è stato modificato da \type{size\_t} a \type{off\_t}
nella revisione POSIX.1-2003 TC1) che consente di ``\textsl{avvisare}'' il
kernel sulle modalità con cui si intende accedere nel futuro ad una certa
porzione di un file, così che esso possa provvedere le opportune
ottimizzazioni; il prototipo di \funcd{posix\_fadvise}\footnote{la funzione è
  stata introdotta su Linux solo a partire dal kernel 2.5.60, ed è disponibile
  soltanto se è stata definita la macro \macro{\_XOPEN\_SOURCE} ad valore di
  almeno \texttt{600} o la macro \macro{\_POSIX\_C\_SOURCE} ad valore di
  almeno \texttt{200112L}.} è:


\begin{funcproto}{
\fhead{fcntl.h}
\fdecl{int posix\_fadvise(int fd, off\_t offset, off\_t len, int advice)}
\fdesc{Dichiara al kernel le future modalità di accesso ad un file.}
}

{La funzione ritorna $0$ in caso di successo e $-1$ per un errore, nel qual
  caso \var{errno} assumerà uno dei valori: 
  \begin{errlist}
    \item[\errcode{EBADF}] l'argomento \param{fd} non è un file descriptor
      valido.
    \item[\errcode{EINVAL}] il valore di \param{advice} non è valido o
      \param{fd} si riferisce ad un tipo di file che non supporta l'operazione
      (come una \textit{pipe} o un socket).
    \item[\errcode{ESPIPE}] previsto dallo standard se \param{fd} è una
      \textit{pipe} o un socket (ma su Linux viene restituito
      \errcode{EINVAL}).
  \end{errlist}
}
\end{funcproto}

La funzione dichiara al kernel le modalità con cui intende accedere alla
regione del file indicato da \param{fd} che inizia alla posizione
\param{offset} e si estende per \param{len} byte. Se per \param{len} si usa un
valore nullo la regione coperta sarà da \param{offset} alla fine del file, ma
questo è vero solo per le versioni più recenti, fino al kernel 2.6.6 il valore
nullo veniva interpretato letteralmente. Le modalità sono indicate
dall'argomento \param{advice} che è una maschera binaria dei valori illustrati
in tab.~\ref{tab:posix_fadvise_flag}, che riprendono il significato degli
analoghi già visti in sez.~\ref{sec:file_memory_map} per
\func{madvise}.\footnote{dato che si tratta dello stesso tipo di funzionalità,
  in questo caso applicata direttamente al sistema ai contenuti di un file
  invece che alla sua mappatura in memoria.} Si tenga presente comunque che la
funzione dà soltanto un avvertimento, non esiste nessun vincolo per il kernel,
che utilizza semplicemente l'informazione.

\begin{table}[htb]
  \centering
  \footnotesize
  \begin{tabular}[c]{|l|p{10cm}|}
    \hline
    \textbf{Valore} & \textbf{Significato} \\
    \hline
    \hline
    \constd{POSIX\_FADV\_NORMAL}  & Non ci sono avvisi specifici da fare
                                   riguardo le modalità di accesso, il
                                   comportamento sarà identico a quello che si
                                   avrebbe senza nessun avviso.\\ 
    \constd{POSIX\_FADV\_SEQUENTIAL}& L'applicazione si aspetta di accedere di
                                   accedere ai dati specificati in maniera
                                   sequenziale, a partire dalle posizioni più
                                   basse.\\ 
    \constd{POSIX\_FADV\_RANDOM}  & I dati saranno letti in maniera
                                   completamente causale.\\
    \constd{POSIX\_FADV\_NOREUSE} & I dati saranno acceduti una sola volta.\\ 
    \constd{POSIX\_FADV\_WILLNEED}& I dati saranno acceduti a breve.\\ 
    \constd{POSIX\_FADV\_DONTNEED}& I dati non saranno acceduti a breve.\\ 
    \hline
  \end{tabular}
  \caption{Valori delle costanti usabili per l'argomento \param{advice} di
    \func{posix\_fadvise}, che indicano la modalità con cui si intende accedere
    ad un file.}
  \label{tab:posix_fadvise_flag}
\end{table}

Come \func{madvise} anche \func{posix\_fadvise} si appoggia al sistema della
memoria virtuale ed al meccanismo standard del \textit{read-ahead} utilizzato
dal kernel; in particolare utilizzando il valore
\const{POSIX\_FADV\_SEQUENTIAL} si raddoppia la dimensione dell'ammontare di
dati letti preventivamente rispetto al default, aspettandosi appunto una
lettura sequenziale che li utilizzerà, mentre con \const{POSIX\_FADV\_RANDOM}
si disabilita del tutto il suddetto meccanismo, dato che con un accesso del
tutto casuale è inutile mettersi a leggere i dati immediatamente successivi
gli attuali; infine l'uso di \const{POSIX\_FADV\_NORMAL} consente di
riportarsi al comportamento di default.

Le due modalità \const{POSIX\_FADV\_NOREUSE} e \const{POSIX\_FADV\_WILLNEED}
fino al kernel 2.6.18 erano equivalenti, a partire da questo kernel la prima
viene non ha più alcun effetto, mentre la seconda dà inizio ad una lettura in
cache della regione del file indicata.  La quantità di dati che verranno letti
è ovviamente limitata in base al carico che si viene a creare sul sistema
della memoria virtuale, ma in genere una lettura di qualche megabyte viene
sempre soddisfatta (ed un valore superiore è solo raramente di qualche
utilità). In particolare l'uso di \const{POSIX\_FADV\_WILLNEED} si può
considerare l'equivalente POSIX di \func{readahead}.

Infine con \const{POSIX\_FADV\_DONTNEED} si dice al kernel di liberare le
pagine di cache occupate dai dati presenti nella regione di file indicata.
Questa è una indicazione utile che permette di alleggerire il carico sulla
cache, ed un programma può utilizzare periodicamente questa funzione per
liberare pagine di memoria da dati che non sono più utilizzati per far posto a
nuovi dati utili; la pagina di manuale riporta l'esempio dello streaming di
file di grosse dimensioni, dove le pagine occupate dai dati già inviati
possono essere tranquillamente scartate.

Sia \func{posix\_fadvise} che \func{readahead} attengono alla ottimizzazione
dell'accesso in lettura; lo standard POSIX.1-2001 prevede anche una funzione
specifica per le operazioni di scrittura, \funcd{posix\_fallocate} (la
funzione è stata introdotta a partire dalle glibc 2.1.94), che consente di
preallocare dello spazio disco per assicurarsi che una seguente scrittura non
fallisca, il suo prototipo, anch'esso disponibile solo se si definisce la
macro \macro{\_XOPEN\_SOURCE} ad almeno 600, è:

\begin{funcproto}{
\fhead{fcntl.h}
\fdecl{int posix\_fallocate(int fd, off\_t offset, off\_t len)}
\fdesc{Richiede la allocazione di spazio disco per un file.} 
}

{La funzione ritorna $0$ in caso di successo e direttamente un codice di
  errore altrimenti, in tal caso \var{errno} non viene impostato, e si otterrà
  direttamente uno dei valori:
  \begin{errlist}
    \item[\errcode{EBADF}] l'argomento \param{fd} non è un file descriptor
      valido o non è aperto in scrittura.
    \item[\errcode{EINVAL}] o \param{offset} o \param{len} sono minori di
      zero.
    \item[\errcode{EFBIG}] il valore di (\param{offset} + \param{len}) eccede
      la dimensione massima consentita per un file.
    \item[\errcode{ENODEV}] l'argomento \param{fd} non fa riferimento ad un
      file regolare.
    \item[\errcode{ENOSPC}] non c'è sufficiente spazio disco per eseguire
      l'operazione. 
    \item[\errcode{ESPIPE}] l'argomento \param{fd} è una \textit{pipe}.
  \end{errlist}
}
\end{funcproto}

La funzione assicura che venga allocato sufficiente spazio disco perché sia
possibile scrivere sul file indicato dall'argomento \param{fd} nella regione
che inizia dalla posizione \param{offset} e si estende per \param{len} byte;
se questa regione si estende oltre la fine del file le dimensioni di
quest'ultimo saranno incrementate di conseguenza. Dopo aver eseguito con
successo la funzione è garantito che una successiva scrittura nella regione
indicata non fallirà per mancanza di spazio disco. La funzione non ha nessun
effetto né sul contenuto, né sulla posizione corrente del file.

Ci si può chiedere a cosa possa servire una funzione come
\func{posix\_fallocate} dato che è sempre possibile ottenere l'effetto voluto
eseguendo esplicitamente sul file la scrittura di una serie di zeri (usando
\funcd{pwrite} per evitare spostamenti della posizione corrente sul file) per
l'estensione di spazio necessaria qualora il file debba essere esteso o abbia
dei buchi.\footnote{si ricordi che occorre scrivere per avere l'allocazione e
  che l'uso di \func{truncate} per estendere un file creerebbe soltanto uno
  \textit{sparse file} (vedi sez.~\ref{sec:file_lseek}) senza una effettiva
  allocazione dello spazio disco.}  In realtà questa è la modalità con cui la
funzione veniva realizzata nella prima versione fornita dalla \acr{glibc}, per
cui la funzione costituiva in sostanza soltanto una standardizzazione delle
modalità di esecuzione di questo tipo di allocazioni.

Questo metodo, anche se funzionante, comporta però l'effettiva esecuzione una
scrittura su tutto lo spazio disco necessario, da fare al momento della
richiesta di allocazione, pagandone il conseguente prezzo in termini di
prestazioni; il tutto quando in realtà servirebbe solo poter riservare lo
spazio per poi andarci a scrivere, una sola volta, quando il contenuto finale
diventa effettivamente disponibile.  Per poter fare tutto questo è però
necessario il supporto da parte del kernel, e questo è divenuto disponibile
solo a partire dal kernel 2.6.23 in cui è stata introdotta la nuova
\textit{system call} \func{fallocate},\footnote{non è detto che la funzione
  sia disponibile per tutti i filesystem, ad esempio per XFS il supporto è
  stato introdotto solo a partire dal kernel 2.6.25.}  che consente di
realizzare direttamente all'interno del kernel l'allocazione dello spazio
disco così da poter realizzare una versione di \func{posix\_fallocate} con
prestazioni molto più elevate; nella \acr{glibc} la nuova \textit{system call}
viene sfruttata per la realizzazione di \func{posix\_fallocate} a partire
dalla versione 2.10.

Trattandosi di una funzione di servizio, ed ovviamente disponibile
esclusivamente su Linux, inizialmente \funcd{fallocate} non era stata definita
come funzione di libreria,\footnote{pertanto poteva essere invocata soltanto
  in maniera indiretta con l'ausilio di \func{syscall}, vedi
  sez.~\ref{sec:proc_syscall}, come \code{long fallocate(int fd, int mode,
      loff\_t offset, loff\_t len)}.} ma a partire dalla \acr{glibc} 2.10 è
  stato fornito un supporto esplicito; il suo prototipo è:

\begin{funcproto}{
\fhead{fcntl.h} 
\fdecl{int fallocate(int fd, int mode, off\_t offset, off\_t len)}
\fdesc{Prealloca dello spazio disco per un file.} 
}

{La funzione ritorna $0$ in caso di successo e $-1$ per un errore, nel qual
  caso \var{errno} assumerà uno dei valori: 
  \begin{errlist}
    \item[\errcode{EBADF}] \param{fd} non fa riferimento ad un file descriptor
      valido aperto in scrittura.
    \item[\errcode{EFBIG}] la somma di \param{offset} e \param{len} eccede le
      dimensioni massime di un file. 
    \item[\errcode{EINVAL}] \param{offset} è minore di zero o \param{len} è
      minore o uguale a zero. 
    \item[\errcode{ENODEV}] \param{fd} non fa riferimento ad un file ordinario
      o a una directory. 
    \item[\errcode{EPERM}] il file è immutabile o \textit{append-only} (vedi
      sez.~\ref{sec:file_perm_overview}).
    \item[\errcode{ENOSYS}] il filesystem contenente il file associato
      a \param{fd} non supporta \func{fallocate}.
    \item[\errcode{EOPNOTSUPP}] il filesystem contenente il file associato
      a \param{fd} non supporta l'operazione \param{mode}.
  \end{errlist}
  ed inoltre \errval{EINTR}, \errval{EIO} e \errval{ENOSPC} nel loro significato
  generico.}
\end{funcproto}

La funzione prende gli stessi argomenti di \func{posix\_fallocate} con lo
stesso significato, a cui si aggiunge l'argomento \param{mode} che indica le
modalità di allocazione; se questo è nullo il comportamento è identico a
quello di \func{posix\_fallocate} e si può considerare \func{fallocate} come
l'implementazione ottimale della stessa a livello di kernel.

Inizialmente l'unico altro valore possibile per \param{mode} era
\const{FALLOC\_FL\_KEEP\_SIZE} che richiede che la dimensione del file
(quella ottenuta nel campo \var{st\_size} di una struttura \struct{stat} dopo
una chiamata a \texttt{fstat}) non venga modificata anche quando la somma
di \param{offset} e \param{len} eccede la dimensione corrente, che serve
quando si deve comunque preallocare dello spazio per scritture in append. In
seguito sono stati introdotti altri valori, riassunti in
tab.\ref{tab:fallocate_mode}, per compiere altre operazioni relative alla
allocazione dello spazio disco dei file.

\begin{table}[htb]
  \centering
  \footnotesize
  \begin{tabular}[c]{|l|p{10cm}|}
    \hline
    \textbf{Valore} & \textbf{Significato} \\
    \hline
    \hline
    \constd{FALLOC\_FL\_INSERT}     & .\\
    \constd{FALLOC\_FL\_COLLAPSE\_RANGE}& .\\ 
    \constd{FALLOC\_FL\_KEEP\_SIZE} & Mantiene invariata la dimensione del
                                     file, pur allocando lo spazio disco anche
                                     oltre la dimensione corrente del file.\\
    \constd{FALLOC\_FL\_PUNCH\_HOLE}& Crea un \textsl{buco} nel file (vedi
                                     sez.~\ref{sec:file_lseek}) rendendolo una
                                     \textit{sparse file} (dal kernel
                                     2.6.38).\\  
    \constd{FALLOC\_FL\_ZERO\_RANGE}& .\\ 
    \hline
  \end{tabular}
  \caption{Valori delle costanti usabili per l'argomento \param{mode} di
    \func{fallocate}.}
  \label{tab:fallocate_mode}
\end{table}

In particolare con \const{FALLOC\_FL\_PUNCH\_HOLE} è possibile scartare il
contenuto della sezione di file indicata da \param{offser} e \param{len},
creando un \textsl{buco} (si ricordi quanto detto in
sez.~\ref{sec:file_lseek}); i blocchi del file interamente contenuti
nell'intervallo verranno disallocati, la parte di intervallo contenuta
parzialmente in altri blocchi verrà riempita con zeri e la lettura dal file
restituirà degli zeri per tutto l'intervallo indicato. In sostanza si rende il
file uno \textit{sparse file} a posteriori.

% vedi http://lwn.net/Articles/226710/ e http://lwn.net/Articles/240571/
% http://kernelnewbies.org/Linux_2_6_23

% TODO aggiungere FALLOC_FL_ZERO_RANGE e FALLOC_FL_COLLAPSE_RANGE, inseriti
% nel kernel 3.15 (sul secondo vedi http://lwn.net/Articles/589260/), vedi
% anche http://lwn.net/Articles/629965/

% TODO aggiungere FALLOC_FL_INSERT vedi  http://lwn.net/Articles/629965/


% TODO non so dove trattarli, ma dal 2.6.39 ci sono i file handle, vedi
% http://lwn.net/Articles/432757/ (probabilmente da associare alle
% at-functions) 


% LocalWords:  dell'I locking multiplexing cap sez system call socket BSD GID
% LocalWords:  descriptor client deadlock NONBLOCK EAGAIN polling select kernel
% LocalWords:  pselect like sys unistd int fd readfds writefds exceptfds struct
% LocalWords:  timeval errno EBADF EINTR EINVAL ENOMEM sleep tab signal void of
% LocalWords:  CLR ISSET SETSIZE POSIX read NULL nell'header l'header glibc fig
% LocalWords:  libc header psignal sigmask SOURCE XOPEN timespec sigset race DN
% LocalWords:  condition sigprocmask tut self trick oldmask poll XPG pollfd l'I
% LocalWords:  ufds unsigned nfds RLIMIT NOFILE EFAULT ndfs events revents hung
% LocalWords:  POLLIN POLLRDNORM POLLRDBAND POLLPRI POLLOUT POLLWRNORM POLLERR
% LocalWords:  POLLWRBAND POLLHUP POLLNVAL POLLMSG SysV stream ASYNC SETOWN FAQ
% LocalWords:  GETOWN fcntl SETFL SIGIO SETSIG Stevens driven siginfo sigaction
% LocalWords:  all'I nell'I Frequently Unanswered Question SIGHUP lease holder
% LocalWords:  breaker truncate write SETLEASE arg RDLCK WRLCK UNLCK GETLEASE
% LocalWords:  uid capabilities capability EWOULDBLOCK notify dall'OR ACCESS st
% LocalWords:  pread readv MODIFY pwrite writev ftruncate creat mknod mkdir buf
% LocalWords:  symlink rename DELETE unlink rmdir ATTRIB chown chmod utime lio
% LocalWords:  MULTISHOT thread linkando librt layer aiocb asyncronous control
% LocalWords:  block ASYNCHRONOUS lseek fildes nbytes reqprio PRIORITIZED sigev
% LocalWords:  PRIORITY SCHEDULING opcode listio sigevent signo value function
% LocalWords:  aiocbp ENOSYS append error const EINPROGRESS fsync return ssize
% LocalWords:  DSYNC fdatasync SYNC cancel ECANCELED ALLDONE CANCELED suspend
% LocalWords:  NOTCANCELED list nent timout sig NOP WAIT NOWAIT size count iov
% LocalWords:  iovec vector EOPNOTSUPP EISDIR len memory mapping mapped swap NB
% LocalWords:  mmap length prot flags off MAP FAILED ANONYMOUS EACCES SHARED SH
% LocalWords:  only ETXTBSY DENYWRITE ENODEV filesystem EPERM EXEC noexec table
% LocalWords:  ENFILE lenght segment violation SIGSEGV FIXED msync munmap copy
% LocalWords:  DoS Denial Service EXECUTABLE NORESERVE LOCKED swapping stack fs
% LocalWords:  GROWSDOWN ANON POPULATE prefaulting SIGBUS fifo VME fork old SFD
% LocalWords:  exec atime ctime mtime mprotect addr mremap address new Failed
% LocalWords:  long MAYMOVE realloc VMA virtual Ingo Molnar remap pages pgoff
% LocalWords:  dall' fault cache linker prelink advisory discrectionary lock fl
% LocalWords:  flock shared exclusive operation dup inode linked NFS cmd ENOLCK
% LocalWords:  EDEADLK whence SEEK CUR type pid GETLK SETLK SETLKW HP EACCESS
% LocalWords:  switch bsd lockf mandatory SVr sgid group root mount mand TRUNC
% LocalWords:  SVID UX Documentation sendfile dnotify inotify NdA ppoll fds add
% LocalWords:  init EMFILE FIONREAD ioctl watch char pathname uint mask ENOSPC
% LocalWords:  CLOSE NOWRITE MOVE MOVED FROM TO rm wd event page ctl acquired
% LocalWords:  attribute Universe epoll Solaris kqueue level triggered Jonathan
% LocalWords:  Lemon BSDCON edge Libenzi kevent backporting epfd EEXIST ENOENT
% LocalWords:  MOD wait EPOLLIN EPOLLOUT EPOLLRDHUP SOCK EPOLLPRI EPOLLERR one
% LocalWords:  EPOLLHUP EPOLLET EPOLLONESHOT shot maxevents ctlv ALL DONT HPUX
% LocalWords:  FOLLOW ONESHOT ONLYDIR FreeBSD EIO caching sysctl instances name
% LocalWords:  watches IGNORED ISDIR OVERFLOW overflow UNMOUNT queued cookie ls
% LocalWords:  NUL sizeof casting printevent nread limits sysconf SC wrapper Di
% LocalWords:  splice result argument DMA controller zerocopy Linus Larry Voy
% LocalWords:  Jens Anxboe vmsplice seek ESPIPE GIFT TCP CORK MSG splicecp nr
% LocalWords:  nwrite segs patch readahead posix fadvise TC advice FADV NORMAL
% LocalWords:  SEQUENTIAL NOREUSE WILLNEED DONTNEED streaming fallocate EFBIG
% LocalWords:  POLLRDHUP half close pwait Gb madvise MADV ahead REMOVE tmpfs it
% LocalWords:  DONTFORK DOFORK shmfs preadv pwritev syscall linux loff head XFS
% LocalWords:  MERGEABLE EOVERFLOW prealloca hole FALLOC KEEP stat fstat union
% LocalWords:  conditions sigwait CLOEXEC signalfd sizemask SIGKILL SIGSTOP ssi
% LocalWords:  sigwaitinfo FifoReporter Windows ptr sigqueue named timerfd TFD
% LocalWords:  clockid CLOCK MONOTONIC REALTIME itimerspec interval Resource
% LocalWords:  ABSTIME gettime temporarily unavailable SIGINT SIGQUIT SIGTERM
% LocalWords:  sigfd fifofd break siginf names starting echo Message from Got
% LocalWords:  message kill received means exit TLOCK ULOCK EPOLLWAKEUP


%%% Local Variables: 
%%% mode: latex
%%% TeX-master: "gapil"
%%% End: 

